% Generated mechanically from the frozen sdga.tex target.
% Do not edit this generated file; edit the bound locale source instead.

% OK, start here.
%

\setcounter{chapter}{23}
\chapter{微分分次层}



\phantomsection
\label{sdga-section-phantom}


\section{引言}
\label{sdga-section-introduction}

\noindent
本章继续《微分分次代数》第 \ref{dga-section-introduction} 节开始的
讨论。一篇综述论文是 \cite{Keller-survey}。




\section{约定}
\label{sdga-section-conventions}

\noindent
本章继续采用如下约定：{\it 环}是指有单位元 $1$ 的交换环。若 $R$
是环，则{\it $R$-代数 $A$}是一个 $R$-模 $A$，连同一个
$R$-双线性映射 $A \times A \to A$（乘法），使得乘法结合且具有
单位元。换言之，这是有单位元的结合 $R$-代数，并且结构映射
$R \to A$ 的像位于 $A$ 的中心中。








\section{分次代数层}
\label{sdga-section-ga}

\noindent
请跳过本节。

\begin{definition}
\label{sdga-definition-ga}
设 $(\mathcal{C}, \mathcal{O})$ 是环化位点。$(\mathcal{C},\mathcal O)$
上的{\it 分次 $\mathcal{O}$-代数层}或{\it 分次代数层}，由以
$n \in \mathbf{Z}$ 为指标的一族 $\mathcal{O}$-模
$\mathcal{A}^n$ 以及 $\mathcal{O}$-双线性映射
$$
\mathcal{A}^n \times \mathcal{A}^m \to \mathcal{A}^{n + m},\quad
(a, b) \longmapsto ab
$$
给出；这些映射称为乘法映射，并满足下列性质：
\begin{enumerate}
\item 乘法是结合的；以及
\item $\mathcal{A}^0$ 有一个全局截面 $1$，它是乘法的双边单位元。
\end{enumerate}
常把这样的结构记作 $\mathcal{A}$。一个{\it 分次
$\mathcal{O}$-代数同态} $f : \mathcal{A} \to \mathcal{B}$，
是与乘法映射相容的一族 $\mathcal{O}$-模映射
$f^n : \mathcal{A}^n \to \mathcal{B}^n$
。
\end{definition}

\noindent
给定一个分次 $\mathcal{O}$-代数 $\mathcal{A}$ 以及一个对象
$U \in \Ob(\mathcal{C})$，使用记号
$$
\mathcal{A}(U) =
\Gamma(U, \mathcal{A}) =
\bigoplus\nolimits_{n \in \mathbf{Z}} \mathcal{A}^n(U)
$$
这是一个分次 $\mathcal{O}(U)$-代数。

\begin{remark}
\label{sdga-remark-functoriality-ga}
设 $(f, f^\sharp) : (\Sh(\mathcal{C}), \mathcal{O}_\mathcal{C})
\to (\Sh(\mathcal{D}), \mathcal{O}_\mathcal{D})$
是环化拓扑斯的态射。则有：
\begin{enumerate}
\item 设 $\mathcal{A}$ 是分次 $\mathcal{O}_\mathcal{C}$-代数。
$\mathcal{A}$ 的乘法映射诱导乘法映射
$f_*\mathcal{A}^n \times f_*\mathcal{A}^m \to f_*\mathcal{A}^{n + m}$
，并且经由 $f^\sharp$，可把它们看作
$\mathcal{O}_\mathcal{D}$-双线性映射。把由此得到的分次
$\mathcal{O}_\mathcal{D}$-代数记作 $f_*\mathcal{A}$。
\item 设 $\mathcal{B}$ 是分次 $\mathcal{O}_\mathcal{D}$-代数。
$\mathcal{B}$ 的乘法映射诱导乘法映射
$f^*\mathcal{B}^n \times f^*\mathcal{B}^m \to f^*\mathcal{B}^{n + m}$
，并且利用 $f^\sharp$，可把它们看作
$\mathcal{O}_\mathcal{C}$-双线性映射。把由此得到的分次
$\mathcal{O}_\mathcal{C}$-代数记作 $f^*\mathcal{B}$。
\item 分次 $\mathcal{O}_\mathcal{C}$-代数同态
$f^*\mathcal{B} \to \mathcal{A}$ 的集合，与分次
$\mathcal{O}_\mathcal{C}$-代数同态
$\mathcal{B} \to f_*\mathcal{A}$ 的集合一一对应。
\end{enumerate}
第 (3) 部分立即由模层上的 $f^*$ 与 $f_*$ 之间通常的伴随得到。
\end{remark}




\section{分次模层}
\label{sdga-section-graded-modules}

\noindent
请跳过本节。

\begin{definition}
\label{sdga-definition-gm}
设 $(\mathcal{C}, \mathcal{O})$ 是环化位点，$\mathcal{A}$ 是
$(\mathcal{C}, \mathcal{O})$ 上的分次代数层。一个（右）
{\it 分次 $\mathcal{A}$-模}，或 $\mathcal{A}$ 上的（右）
{\it 分次模}，由以 $n \in \mathbf{Z}$ 为指标的一族
$\mathcal{O}$-模 $\mathcal{M}^n$ 以及 $\mathcal{O}$-双线性映射
$$
\mathcal{M}^n \times \mathcal{A}^m \to \mathcal{M}^{n + m},\quad
(x, a) \longmapsto xa
$$
给出；这些映射称为乘法映射，并满足下列性质：
\begin{enumerate}
\item 乘法满足 $(xa)a' = x(aa')$；
\item 对每个 $n$，$\mathcal{A}^0$ 的单位截面 $1$
在 $\mathcal{M}^n$ 上的作用都是恒等作用。
\end{enumerate}
我们常说“设 $\mathcal{M}$ 是一个分次 $\mathcal{A}$-模”来表示上述情形。
一个{\it 分次 $\mathcal{A}$-模同态}
$f : \mathcal{M} \to \mathcal{N}$，是与乘法映射相容的一族
$\mathcal{O}$-模映射 $f^n : \mathcal{M}^n \to \mathcal{N}^n$。
（右）分次 $\mathcal{A}$-模构成的范畴记作
$\textit{Mod}(\mathcal{A})$。
\end{definition}

\noindent
可以完全类似地定义{\it 左分次模}，但本章默认所说的模都是右模。

\medskip\noindent
给定一个分次 $\mathcal{A}$-模 $\mathcal{M}$ 和一个对象
$U \in \Ob(\mathcal{C})$，我们使用记号
$$
\mathcal{M}(U) =
\Gamma(U, \mathcal{M}) =
\bigoplus\nolimits_{n \in \mathbf{Z}} \mathcal{M}^n(U)
$$
这是一个（右）分次 $\mathcal{A}(U)$-模。

\begin{lemma}
\label{sdga-lemma-gm-abelian}
设 $(\mathcal{C}, \mathcal{O})$ 是环化位点，$\mathcal{A}$ 是分次
$\mathcal{O}$-代数。范畴 $\textit{Mod}(\mathcal{A})$ 是阿贝尔范畴，
并具有下列性质：
\begin{enumerate}
\item $\textit{Mod}(\mathcal{A})$ 有任意直和；
\item $\textit{Mod}(\mathcal{A})$ 有任意余极限；
\item $\textit{Mod}(\mathcal{A})$ 中的滤过余极限是正合的；
\item $\textit{Mod}(\mathcal{A})$ 有任意积；
\item $\textit{Mod}(\mathcal{A})$ 有任意极限。
\end{enumerate}
函子
$$
\textit{Mod}(\mathcal{A}) \longrightarrow \textit{Mod}(\mathcal{O}),\quad
\mathcal{M} \longmapsto \mathcal{M}^n
$$
把一个分次 $\mathcal{A}$-模映到其第 $n$ 项；它与所有极限和余极限交换。
\end{lemma}

\noindent
该引理说明，我们可以逐项取极限和余极限。它还说明（或者，如果你愿意，
蕴含）遗忘函子
$$
\textit{Mod}(\mathcal{A})
\longrightarrow
\text{分次 }\mathcal{O}\text{-模}
$$
与所有极限和余极限交换。

\begin{proof}
把引理陈述中的函子记作
$\text{gr}^n : \textit{Mod}(\mathcal{A}) \to \textit{Mod}(\mathcal{O})$
。考虑分次 $\mathcal{A}$-模同态
$f : \mathcal{M} \to \mathcal{N}$。把 $f$ 看成分次
$\mathcal{O}$-模映射时，它的核与余核还带有定义
\ref{sdga-definition-gm} 中的乘法映射。因此，它们也是
$\textit{Mod}(\mathcal{A})$ 中的核与余核。由此
$\textit{Mod}(\mathcal{A})$ 是阿贝尔范畴，并且取核与余核和
$\text{gr}^n$ 交换。

\medskip\noindent
为了证明极限和余极限存在，只需证明积和直和存在；参见《范畴》引理
\ref{categories-lemma-limits-products-equalizers} 和
\ref{categories-lemma-colimits-coproducts-coequalizers}。
同样的引理还表明，只要 $\text{gr}^n$ 与直和及积交换，
就可推出它与极限及余极限交换。

\medskip\noindent
设 $\mathcal{M}_t$（$t \in T$）是一族分次 $\mathcal{A}$-模。
可以考虑这样一个分次 $\mathcal{A}$-模：其次数 $n$ 项是
$\bigoplus_{t \in T} \mathcal{M}_t^n$（乘法映射是显然的）。
读者容易验证，这是 $\textit{Mod}(\mathcal{A})$ 中的直和。
积的情形类似。

\medskip\noindent
注意，对每个 $n$，$\text{gr}^n$ 都是正合函子；而
$\textit{Mod}(\mathcal{A})$ 中的复形
$\mathcal{M}_1 \to \mathcal{M}_2 \to \mathcal{M}_3$ 正合，当且仅当
$\text{gr}^n\mathcal{M}_1 \to \text{gr}^n\mathcal{M}_2 \to
\text{gr}^n\mathcal{M}_3$ 对每个 $n$ 都在
$\textit{Mod}(\mathcal{O})$ 中正合。因此，由于
$\textit{Mod}(\mathcal{O})$ 中的滤过余极限是正合的，我们得到 (3)；
它是格罗滕迪克阿贝尔范畴，参见《位点上的上同调》第
\ref{sites-cohomology-section-unbounded} 节。
\end{proof}






\section{分次模层的分次范畴}
\label{sdga-section-gm-gr-cat}

\noindent
请跳过本节。本节是《微分分次代数》例
\ref{dga-example-gm-gr-cat} 的类似物。关于分次范畴的约定，
请参见《微分分次代数》第 \ref{dga-section-graded} 节。

\medskip\noindent
设 $(\mathcal{C}, \mathcal{O})$ 是环化位点，$\mathcal{A}$ 是
$(\mathcal{C}, \mathcal{O})$ 上的分次代数层。我们将构造
$R = \Gamma(\mathcal{C}, \mathcal{O})$
上的分次范畴 $\textit{Mod}^{gr}(\mathcal{A})$，使其对应范畴
$(\textit{Mod}^{gr}(\mathcal{A}))^0$ 是分次 $\mathcal{A}$-模范畴。
取右分次 $\mathcal{A}$-模作为 $\textit{Mod}^{gr}(\mathcal{A})$
的对象（见第 \ref{sdga-section-graded-modules} 节）。给定分次
$\mathcal{A}$-模 $\mathcal{L}$ 和 $\mathcal{M}$，置
$$
\Hom_{\textit{Mod}^{gr}(\mathcal{A})}(\mathcal{L}, \mathcal{M}) =
\bigoplus\nolimits_{n \in \mathbf{Z}} \Hom^n(\mathcal{L}, \mathcal{M})
$$
，其中 $\Hom^n(\mathcal{L}, \mathcal{M})$ 是次数 $n$ 齐次的右
$\mathcal{A}$-模映射 $f : \mathcal{L} \to \mathcal{M}$ 所成的集合。
更确切地说，$f$ 由对 $i \in \mathbf{Z}$ 给出的一族映射
$f : \mathcal{L}^i \to \mathcal{M}^{i + n}$ 组成，并与乘法映射相容。
用分量表示，
$$
\Hom^n(\mathcal{L}, \mathcal{M})
\subset
\prod\nolimits_{p + q = n}
\Hom_\mathcal{O}(\mathcal{L}^{-q}, \mathcal{M}^p)
$$
（注意指标的反转）是由满足
$$
f_{p, q}(m a) = f_{p - i, q + i}(m)a
$$
的 $f = (f_{p, q})$ 组成的子集；这里 $a$ 是 $\mathcal{A}^i$
的局部截面，$m$ 是 $\mathcal{L}^{-q-i}$ 的局部截面。
对分次 $\mathcal{A}$-模 $\mathcal{K}$、$\mathcal{L}$、
$\mathcal{M}$，通过映射
$$
\Hom^m(\mathcal{L}, \mathcal{M}) \times
\Hom^n(\mathcal{K}, \mathcal{L}) \longrightarrow
\Hom^{n + m}(\mathcal{K}, \mathcal{M})
$$
定义 $\textit{Mod}^{gr}(\mathcal{A})$ 中的复合；它就是右
$\mathcal{A}$-模映射的通常复合：$(g,f)\mapsto g\circ f$。





\section{分次模层的张量积}
\label{sdga-section-tensor-product}

\noindent
请跳过本节。本节是《微分分次代数》第
\ref{dga-section-tensor-product} 节部分内容的类似物。

\medskip\noindent
设 $(\mathcal{C}, \mathcal{O})$ 是环化位点，$\mathcal{A}$ 是
$(\mathcal{C},\mathcal{O})$ 上的分次代数层。设 $\mathcal{M}$ 是右分次
$\mathcal{A}$-模，$\mathcal{N}$ 是左分次 $\mathcal{A}$-模。
定义{\it 张量积}
$\mathcal{M} \otimes_\mathcal{A} \mathcal{N}$
为这样的分次 $\mathcal{O}$-模：它的次数 $n$ 项是
$$
(\mathcal{M} \otimes_\mathcal{A} \mathcal{N})^n =
\Coker\left(
\bigoplus\nolimits_{r + s + t = n} \mathcal{M}^r \otimes_\mathcal{O}
\mathcal{A}^s \otimes_\mathcal{O} \mathcal{N}^t
\longrightarrow
\bigoplus\nolimits_{p + q = n} \mathcal{M}^p \otimes_\mathcal{O} \mathcal{N}^q
\right)
$$
，其中该映射把
$\mathcal{M}^r\otimes_\mathcal{O}\mathcal{A}^s
\otimes_\mathcal{O}\mathcal{N}^t$ 的局部截面
$x\otimes a\otimes y$ 映到 $xa\otimes y-x\otimes ay$。
按此定义，
$(\mathcal{M} \otimes_\mathcal{A} \mathcal{N})^n$
是预层
$U \mapsto (\mathcal{M}(U) \otimes_{\mathcal{A}(U)} \mathcal{N}(U))^n$
的层化；这里分次模的张量积按《微分分次代数》第
\ref{dga-section-tensor-product} 节中的定义。

\medskip\noindent
固定左分次 $\mathcal{A}$-模 $\mathcal{N}$，便得到函子
$$
- \otimes_\mathcal{A} \mathcal{N} :
\textit{Mod}(\mathcal{A})
\longrightarrow
\text{Gr}(\textit{Mod}(\mathcal{O})) =
\text{分次 }\mathcal{O}\text{-模}
$$
关于记号 $\text{Gr}(-)$，请参见《同调代数》定义
\ref{homology-definition-graded}。分次 $\mathcal{O}$-模的分次范畴记作
$\text{Gr}^{gr}(\textit{Mod}(\mathcal{O}))$；见《微分分次代数》例
\ref{dga-example-graded-category-graded-objects}。
上述函子可以提升为分次范畴之间的函子
$$
- \otimes_\mathcal{A} \mathcal{N} :
\textit{Mod}^{gr}(\mathcal{A})
\longrightarrow
\text{Gr}^{gr}(\textit{Mod}(\mathcal{O}))
$$
：把从 $\mathcal{M}$ 到 $\mathcal{M}'$ 的次数 $n$ 同态映到
从 $\mathcal{M}\otimes_\mathcal{A}\mathcal{N}$ 到
$\mathcal{M}'\otimes_\mathcal{A}\mathcal{N}$ 的诱导次数 $n$ 映射。






\section{分次模层的内 Hom}
\label{sdga-section-internal-hom-graded}

\noindent
我们强烈建议读者跳过本节。

\medskip\noindent
我们需要第 \ref{sdga-section-gm-gr-cat} 节中构造的层化版本。
设 $(\mathcal{C}, \mathcal{O})$、$\mathcal{A}$、
$\mathcal{M}$、$\mathcal{L}$ 如第 \ref{sdga-section-gm-gr-cat} 节所述。定义
$$
\SheafHom^{gr}_\mathcal{A}(\mathcal{M}, \mathcal{L})
$$
为这样的分次 $\mathcal{O}$-模：其次数 $n$ 项
$$
\SheafHom^n_\mathcal{A}(\mathcal{M}, \mathcal{L})
\subset
\prod\nolimits_{p + q = n}
\SheafHom_\mathcal{O}(\mathcal{L}^{-q}, \mathcal{M}^p)
$$
是由满足
$$
f_{p, q}(m a) = f_{p - i, q + i}(m)a
$$
的局部截面 $f=(f_{p,q})$ 组成的子层；这里 $a$ 是
$\mathcal{A}^i$ 的局部截面，$m$ 是 $\mathcal{L}^{-q-i}$
的局部截面。如同第 \ref{sdga-section-gm-gr-cat} 节，有复合映射
$$
\SheafHom^{gr}_\mathcal{A}(\mathcal{L}, \mathcal{M}) \otimes_\mathcal{O}
\SheafHom^{gr}_\mathcal{A}(\mathcal{K}, \mathcal{L})
\longrightarrow
\SheafHom^{gr}_\mathcal{A}(\mathcal{K}, \mathcal{M})
$$
，其中左端是第 \ref{sdga-section-tensor-product} 节所定义的分次
$\mathcal{O}$-模张量积。该映射由复合映射
$$
\SheafHom^m_\mathcal{A}(\mathcal{L}, \mathcal{M}) \otimes_\mathcal{O}
\SheafHom^n_\mathcal{A}(\mathcal{K}, \mathcal{L}) \longrightarrow
\SheafHom^{n + m}_\mathcal{A}(\mathcal{K}, \mathcal{M})
$$
给出；后者（局部地）按通常复合定义。

\medskip\noindent
按这些定义，有
$$
\Hom_{\textit{Mod}^{gr}(\mathcal{A})}(\mathcal{L}, \mathcal{M}) =
\Gamma(\mathcal{C}, \SheafHom^{gr}_\mathcal{A}(\mathcal{L}, \mathcal{M}))
$$
作为分次 $R$-模的等式，并且它与复合相容。







\section{分次双模层与张量--Hom 伴随}
\label{sdga-section-graded-bimodules}

\noindent
请跳过本节。

\begin{definition}
\label{sdga-definition-bimodule}
设 $(\mathcal{C},\mathcal{O})$ 是环化位点，$\mathcal{A}$ 和
$\mathcal{B}$ 是 $(\mathcal{C},\mathcal{O})$ 上的分次代数层。
一个{\it 分次 $(\mathcal{A},\mathcal{B})$-双模}，由以
$n\in\mathbf{Z}$ 为指标的一族 $\mathcal{O}$-模
$\mathcal{M}^n$ 以及 $\mathcal{O}$-双线性映射
$$
\mathcal{M}^n \times \mathcal{B}^m \to \mathcal{M}^{n + m},\quad
(x, b) \longmapsto xb
$$
和
$$
\mathcal{A}^n \times \mathcal{M}^m \to \mathcal{M}^{n + m},\quad
(a, x) \longmapsto ax
$$
给出；这些映射称为乘法映射，并满足下列性质：
\begin{enumerate}
\item 乘法满足 $a(a'x)=(aa')x$ 和 $(xb)b'=x(bb')$；
\item $(ax)b = a(xb)$；
\item $\mathcal{A}^0$ 的单位截面 $1$ 通过乘法作恒等作用；
\item $\mathcal{B}^0$ 的单位截面 $1$ 通过乘法作恒等作用。
\end{enumerate}
我们常把这样的结构记作 $\mathcal{M}$。一个{\it 分次
$(\mathcal{A},\mathcal{B})$-双模同态}
$f:\mathcal{M}\to\mathcal{N}$，是与乘法映射相容的一族
$\mathcal{O}$-模映射 $f^n:\mathcal{M}^n\to\mathcal{N}^n$。
\end{definition}

\noindent
给定分次 $(\mathcal{A},\mathcal{B})$-双模 $\mathcal{M}$ 和对象
$U\in\Ob(\mathcal{C})$，我们使用记号
$$
\mathcal{M}(U) =
\Gamma(U, \mathcal{M}) =
\bigoplus\nolimits_{n \in \mathbf{Z}} \mathcal{M}^n(U)
$$
这是分次 $(\mathcal{A}(U),\mathcal{B}(U))$-双模。

\medskip\noindent
设 $(\mathcal{C},\mathcal{O})$ 是环化位点，$\mathcal{A}$ 和
$\mathcal{B}$ 是 $(\mathcal{C},\mathcal{O})$ 上的分次代数层。
设 $\mathcal{M}$ 是右分次 $\mathcal{A}$-模，$\mathcal{N}$ 是分次
$(\mathcal{A},\mathcal{B})$-双模。在这种情形下，第
\ref{sdga-section-tensor-product} 节定义的分次张量积
$$
\mathcal{M} \otimes_\mathcal{A} \mathcal{N}
$$
是右分次 $\mathcal{B}$-模，其乘法映射是显然的。该构造定义一个函子
以及一个分次范畴之间的函子
$$
\otimes_\mathcal{A} \mathcal{N} :
\textit{Mod}(\mathcal{A})
\longrightarrow
\textit{Mod}(\mathcal{B})
\quad\text{且}\quad
\otimes_\mathcal{A} \mathcal{N} :
\textit{Mod}^{gr}(\mathcal{A})
\longrightarrow
\textit{Mod}^{gr}(\mathcal{B})
$$
：把从 $\mathcal{M}$ 到 $\mathcal{M}'$ 的次数 $n$ 同态映到
从 $\mathcal{M}\otimes_\mathcal{A}\mathcal{N}$ 到
$\mathcal{M}'\otimes_\mathcal{A}\mathcal{N}$ 的诱导次数 $n$ 映射。

\medskip\noindent
设 $(\mathcal{C},\mathcal{O})$ 是环化位点，$\mathcal{A}$ 和
$\mathcal{B}$ 是 $(\mathcal{C},\mathcal{O})$ 上的分次代数层。
设 $\mathcal{N}$ 是分次 $(\mathcal{A},\mathcal{B})$-双模，
$\mathcal{L}$ 是右分次 $\mathcal{B}$-模。在这种情形下，第
\ref{sdga-section-internal-hom-graded} 节定义的分次内 Hom
$$
\SheafHom_\mathcal{B}^{gr}(\mathcal{N}, \mathcal{L})
$$
是右分次 $\mathcal{A}$-模，其乘法映射为\footnote{我们在这里约定，
此式不涉及任何符号。}
$$
\SheafHom^n_\mathcal{B}(\mathcal{N}, \mathcal{L})
\times \mathcal{A}^m
\longrightarrow
\SheafHom^{n + m}_\mathcal{B}(\mathcal{N}, \mathcal{L})
$$
：它把 $U$ 上
$\SheafHom^n_\mathcal{B}(\mathcal{N},\mathcal{L})$ 的截面
$f=(f_{p,q})$ 和 $U$ 上 $\mathcal{A}^m$ 的截面 $a$，映到
$U$ 上 $\SheafHom^{n+m}_\mathcal{B}(\mathcal{N},\mathcal{L})$
的截面 $fa$；后者定义为映射族
$$
\mathcal{N}^{-q - m}|_U \xrightarrow{a \cdot -}
\mathcal{N}^{-q}|_U \xrightarrow{f_{p, q}}
\mathcal{M}^p|_U
$$
。略去其良定义性的验证。该构造定义一个函子以及一个分次范畴之间的函子
$$
\SheafHom_\mathcal{B}^{gr}(\mathcal{N}, -) :
\textit{Mod}(\mathcal{B})
\longrightarrow
\textit{Mod}(\mathcal{A})
\quad\text{且}\quad
\SheafHom_\mathcal{B}^{gr}(\mathcal{N}, -) :
\textit{Mod}^{gr}(\mathcal{B})
\longrightarrow
\textit{Mod}^{gr}(\mathcal{A})
$$
：把从 $\mathcal{L}$ 到 $\mathcal{L}'$ 的次数 $n$ 同态映到
从 $\SheafHom_\mathcal{B}^{gr}(\mathcal{N},\mathcal{L})$ 到
$\SheafHom_\mathcal{B}^{gr}(\mathcal{N},\mathcal{L}')$
的诱导次数 $n$ 映射。

\begin{lemma}
\label{sdga-lemma-tensor-hom-adjunction-gr}
设 $(\mathcal{C},\mathcal{O})$ 是环化位点，$\mathcal{A}$ 和
$\mathcal{B}$ 是 $(\mathcal{C},\mathcal{O})$ 上的分次代数层。
设 $\mathcal{M}$ 是右分次 $\mathcal{A}$-模，$\mathcal{N}$ 是分次
$(\mathcal{A},\mathcal{B})$-双模，$\mathcal{L}$ 是右分次
$\mathcal{B}$-模。采用上述约定，有
$$
\Hom_{\textit{Mod}^{gr}(\mathcal{B})}(
\mathcal{M} \otimes_\mathcal{A} \mathcal{N}, \mathcal{L}) =
\Hom_{\textit{Mod}^{gr}(\mathcal{A})}(
\mathcal{M}, \SheafHom_\mathcal{B}^{gr}(\mathcal{N}, \mathcal{L}))
$$
以及
$$
\SheafHom_\mathcal{B}^{gr}(
\mathcal{M} \otimes_\mathcal{A} \mathcal{N}, \mathcal{L}) =
\SheafHom_\mathcal{A}^{gr}(
\mathcal{M}, \SheafHom_\mathcal{B}^{gr}(\mathcal{N}, \mathcal{L}))
$$
；这些等式关于 $\mathcal{M}$、$\mathcal{N}$、$\mathcal{L}$
都具有函子性。
\end{lemma}

\begin{proof}
略。提示：把两端都解释为 $\mathcal{A}$-双线性的分次映射
$\psi : \mathcal{M} \times \mathcal{N} \to \mathcal{L}$
即可；这些映射在右侧关于 $\mathcal{B}$ 是线性的。
\end{proof}

\noindent
设 $(\mathcal{C},\mathcal{O})$ 是环化位点，$\mathcal{A}$ 和
$\mathcal{B}$ 是 $(\mathcal{C},\mathcal{O})$ 上的分次代数层。
作为上述情形的特例，设给定分次 $\mathcal{O}$-代数同态
$\varphi:\mathcal{A}\to\mathcal{B}$。于是得到一个函子和一个
分次范畴之间的函子
$$
\otimes_{\mathcal{A}, \varphi} \mathcal{B} :
\textit{Mod}(\mathcal{A})
\longrightarrow
\textit{Mod}(\mathcal{B})
\quad\text{且}\quad
\otimes_{\mathcal{A}, \varphi} \mathcal{B} :
\textit{Mod}^{gr}(\mathcal{A})
\longrightarrow
\textit{Mod}^{gr}(\mathcal{B})
$$
。另一方面，有限制标量函子
$$
res_\varphi :
\textit{Mod}(\mathcal{B})
\longrightarrow
\textit{Mod}(\mathcal{A})
\quad\text{且}\quad
res_\varphi :
\textit{Mod}^{gr}(\mathcal{B})
\longrightarrow
\textit{Mod}^{gr}(\mathcal{A})
$$
。由上一引理可以证明这些函子彼此伴随（如限制标量与基变换的通常情形）。
具体地，把 $\mathcal{B}$ 视为分次 $(\mathcal{A},\mathcal{B})$-双模时
记作 ${}_\mathcal{A}\mathcal{B}_\mathcal{B}$。则对任意右分次
$\mathcal{B}$-模 $\mathcal{L}$，有
$$
\SheafHom_\mathcal{B}^{gr}({}_\mathcal{A}\mathcal{B}_\mathcal{B}, \mathcal{L})
= res_\varphi(\mathcal{L})
$$
作为右分次 $\mathcal{A}$-模的等式。因此，引理
\ref{sdga-lemma-tensor-hom-adjunction-gr} 给出函子性同构
$$
\Hom_{\textit{Mod}^{gr}(\mathcal{B})}(
\mathcal{M} \otimes_{\mathcal{A}, \varphi} \mathcal{B}, \mathcal{L}) =
\Hom_{\textit{Mod}^{gr}(\mathcal{A})}(
\mathcal{M}, res_\varphi(\mathcal{L}))
$$
。若从上下文可知 $\varphi$，通常在此公式中省略对它的依赖。
以同样方式得到分次 $\mathcal{O}$-模的等式
$$
\SheafHom^{gr}_\mathcal{B}(
\mathcal{M} \otimes_\mathcal{A} \mathcal{B}, \mathcal{L}) =
\SheafHom_\mathcal{A}^{gr}(\mathcal{M}, \mathcal{L})
$$
。




\section{分次模层的拉回与直像}
\label{sdga-section-functoriality-graded}

\noindent
我们建议读者跳过本节。

\medskip\noindent
设 $(f, f^\sharp) : (\Sh(\mathcal{C}), \mathcal{O}_\mathcal{C})
\to (\Sh(\mathcal{D}), \mathcal{O}_\mathcal{D})$
是环化拓扑斯的态射。设 $\mathcal{A}$ 是分次
$\mathcal{O}_\mathcal{C}$-代数，$\mathcal{B}$ 是分次
$\mathcal{O}_\mathcal{D}$-代数。设给定分次
$f^{-1}\mathcal{O}_\mathcal{D}$-代数同态
$$
\varphi : f^{-1}\mathcal{B} \to \mathcal{A}
$$
。由限制标量与扩张标量的伴随，这等价于一个分次
$\mathcal{O}_\mathcal{C}$-代数同态
$\varphi:f^*\mathcal{B}\to\mathcal{A}$；也等价地，可以把
$\varphi$ 看作分次 $\mathcal{O}_\mathcal{D}$-代数同态
$$
\varphi : \mathcal{B} \to f_*\mathcal{A}
$$
。参见评注 \ref{sdga-remark-functoriality-ga}。

\medskip\noindent
先定义函子
$$
f_* :
\textit{Mod}(\mathcal{A})
\longrightarrow
\textit{Mod}(\mathcal{B})
$$
。给定分次 $\mathcal{A}$-模 $\mathcal{M}$，定义
$f_*\mathcal{M}$ 为这样的分次 $\mathcal{B}$-模：其次数 $n$ 项是
$f_*\mathcal{M}^n$，乘法取为
$$
f_*\mathcal{M}^n \times \mathcal{B}^m
\xrightarrow{(\text{id}, \varphi^m)}
f_*\mathcal{M}^n \times f_*\mathcal{A}^m
\xrightarrow{f_*\mu_{n, m}}
f_*\mathcal{M}^{n + m}
$$
，其中 $\mu_{n,m}:\mathcal{M}^n\times\mathcal{A}^m
\to\mathcal{M}^{n+m}$ 是 $\mathcal{M}$ 作为 $\mathcal{A}$-模的
乘法映射。这里使用了 $f_*$ 与积交换这一事实。该构造关于
$\mathcal{M}$ 显然具有函子性，因而得到所需函子。

\medskip\noindent
再定义函子
$$
f^* :
\textit{Mod}(\mathcal{B})
\longrightarrow
\textit{Mod}(\mathcal{A})
$$
为下列函子的复合：
$$
\textit{Mod}(\mathcal{B})
\xrightarrow{f^{-1}}
\textit{Mod}(f^{-1}\mathcal{B})
\xrightarrow{ - \otimes_{f^{-1}\mathcal{B}} \mathcal{A}}
\textit{Mod}(\mathcal{A})
$$
首先，给定分次 $\mathcal{B}$-模 $\mathcal{N}$，定义
$f^{-1}\mathcal{N}$ 为这样的分次 $f^{-1}\mathcal{B}$-模：
其次数 $n$ 项是 $f^{-1}\mathcal{N}^n$，乘法取为
$$
f^{-1}\nu_{n, m} :
f^{-1}\mathcal{N}^n \times f^{-1}\mathcal{B}^m
\longrightarrow
f^{-1}\mathcal{N}^{n + m}
$$
，其中 $\nu_{n,m}:\mathcal{N}^n\times\mathcal{B}^m
\to\mathcal{N}^{n+m}$ 是 $\mathcal{N}$ 作为 $\mathcal{B}$-模的
乘法映射。这里使用了 $f^{-1}$ 与积交换这一事实。该构造关于
$\mathcal{N}$ 显然具有函子性，从而得到函子 $f^{-1}$。
然后利用第 \ref{sdga-section-graded-bimodules} 节关于张量积的讨论，
定义函子
$$
- \otimes_{f^{-1}\mathcal{B}} \mathcal{A} :
\textit{Mod}(f^{-1}\mathcal{B})
\longrightarrow
\textit{Mod}(\mathcal{A})
$$
。最后，如上所预示，置
$$
f^*\mathcal{N} =
f^{-1}\mathcal{N} \otimes_{f^{-1}\mathcal{B}, \varphi} \mathcal{A}
$$
。

\medskip\noindent
函子 $f_*$ 和 $f^*$ 很容易提升为分次范畴之间的函子
$$
f_* :
\textit{Mod}^{gr}(\mathcal{A})
\longrightarrow
\textit{Mod}^{gr}(\mathcal{B})
\quad\text{且}\quad
f^* :
\textit{Mod}^{gr}(\mathcal{B})
\longrightarrow
\textit{Mod}^{gr}(\mathcal{A})
$$
；它们在底层对象上的作用不变，而在次数 $n$ 的齐次态射上由上述构造的
函子性定义。

\begin{lemma}
\label{sdga-lemma-adjunction-push-pull-gr}
在上述情形下，有
$$
\Hom_{\textit{Mod}^{gr}(\mathcal{B})}(
\mathcal{N}, f_*\mathcal{M}) =
\Hom_{\textit{Mod}^{gr}(\mathcal{A})}(
f^*\mathcal{N}, \mathcal{M})
$$
\end{lemma}

\begin{proof}
略。提示：先利用阿贝尔群层上 $f^{-1}$ 与 $f_*$ 的伴随，证明
$f^{-1}$ 与 $f_*$ 作为 $\textit{Mod}(\mathcal{B})$ 和
$\textit{Mod}(f^{-1}\mathcal{B})$ 之间的函子彼此伴随。
再使用第 \ref{sdga-section-graded-bimodules} 节给出的基变换与限制标量的伴随。
\end{proof}





\section{局部化与分次模层}
\label{sdga-section-localize-graded}

\noindent
我们建议读者跳过本节。

\medskip\noindent
设 $(\mathcal{C},\mathcal{O})$ 是环化位点，$U\in\Ob(\mathcal{C})$。
记
$$
j :
(\Sh(\mathcal{C}/U), \mathcal{O}_U)
\longrightarrow
(\Sh(\mathcal{C}), \mathcal{O})
$$
为相应的局部化态射（《位点上的模》第
\ref{sites-modules-section-localize} 节）。下面将使用如下事实：
对 $\mathcal{O}_U$-模 $\mathcal{M}_i$（$i=1,2$）和
$\mathcal{O}$-模 $\mathcal{A}$，有典范映射
$$
j_! :
\Hom_{\mathcal{O}_U}(
\mathcal{M}_1 \otimes_{\mathcal{O}_U} \mathcal{A}|_U, \mathcal{M}_2)
\longrightarrow
\Hom_\mathcal{O}(
j_!\mathcal{M}_1 \otimes_\mathcal{O} \mathcal{A}, j_!\mathcal{M}_2)
$$
。事实上，由《位点上的模》引理
\ref{sites-modules-lemma-j-shriek-and-tensor}，有
$j_!(\mathcal{M}_1 \otimes_{\mathcal{O}_U} \mathcal{A}|_U) =
j_!\mathcal{M}_1 \otimes_\mathcal{O} \mathcal{A}$。

\medskip\noindent
设 $\mathcal{A}$ 是分次 $\mathcal{O}$-代数。用
$\mathcal{A}_U$ 表示 $\mathcal{A}$ 在 $\mathcal{C}/U$ 上的限制；
换言之，$\mathcal{A}_U=j^*\mathcal{A}=j^{-1}\mathcal{A}$。
第 \ref{sdga-section-functoriality-graded} 节已经构造伴随函子
$$
j_* :
\textit{Mod}^{gr}(\mathcal{A}_U)
\longrightarrow
\textit{Mod}^{gr}(\mathcal{A})
\quad\text{且}\quad
j^* :
\textit{Mod}^{gr}(\mathcal{A})
\longrightarrow
\textit{Mod}^{gr}(\mathcal{A}_U)
$$
，其中 $j^*$ 是 $j_*$ 的左伴随。我们断言，此外还有正合函子
$$
j_! :
\textit{Mod}^{gr}(\mathcal{A}_U)
\longrightarrow
\textit{Mod}^{gr}(\mathcal{A})
$$
，它是 $j^*$ 的左伴随。具体地，给定分次 $\mathcal{A}_U$-模
$\mathcal{M}$，定义 $j_!\mathcal{M}$ 为这样的分次 $\mathcal{A}$-模：
其次数 $n$ 项为 $j_!\mathcal{M}^n$，乘法映射取为
$$
j_!\mu_{n, m} :
j_!\mathcal{M}^n \times \mathcal{A}^m \to
j_!\mathcal{M}^{n + m}
$$
，其中 $\mu_{m,n}:\mathcal{M}^n\times\mathcal{A}^m
\to\mathcal{M}^{n+m}$ 是给定的乘法映射。给定分次
$\mathcal{A}_U$-模之间次数 $n$ 的齐次映射
$f:\mathcal{M}\to\mathcal{M}'$，得到次数 $n$ 的齐次映射
$j_!f:j_!\mathcal{M}\to j_!\mathcal{M}'$。这样便得到所需函子。

\begin{lemma}
\label{sdga-lemma-extension-by-zero-graded}
在上述情形下，有
$$
\Hom_{\textit{Mod}^{gr}(\mathcal{A})}(
j_!\mathcal{M}, \mathcal{N}) =
\Hom_{\textit{Mod}^{gr}(\mathcal{A}_U)}(
\mathcal{M}, j^*\mathcal{N})
$$
\end{lemma}

\begin{proof}
由《位点上的模》第 \ref{sites-modules-section-localize} 节的讨论，
$\mathcal{O}$-模上的函子 $j_!$ 与 $j^*$ 彼此伴随。因此，若只看
$\mathcal{O}$-模结构，我们知道
$$
\Hom_{\text{Gr}^{gr}(\textit{Mod}(\mathcal{O}))}(
j_!\mathcal{M}, \mathcal{N}) =
\Hom_{\text{Gr}^{gr}(\textit{Mod}(\mathcal{O}_U))}(
\mathcal{M}, j^*\mathcal{N})
$$
（回忆，$\text{Gr}^{gr}(\textit{Mod}(\mathcal{O}))$ 表示分次
$\mathcal{O}$-模的分次范畴）。还须检验这些等同把左端的
$\mathcal{A}$-模映射对应到右端的 $\mathcal{A}_U$-模映射。
为此，设 $\mathcal{O}_U$-线性映射
$f^n : \mathcal{M}^n \to j^*\mathcal{N}^{n + d}$
对应于 $\mathcal{O}$-线性映射
$g^n : j_!\mathcal{M}^n \to \mathcal{N}^{n + d}$
。只需证明下图交换
$$
\xymatrix{
\mathcal{M}^n \otimes_{\mathcal{O}_U} \mathcal{A}_U^m
\ar[r]_{f^n \otimes 1} \ar[d] &
j^*\mathcal{N}^{n + d} \otimes_{\mathcal{O}_U} \mathcal{A}_U^m \ar[d] \\
\mathcal{M}^{n + m} \ar[r]^{f^{n + m}} &
j^*\mathcal{N}^{n + m + d}
}
$$
当且仅当下图交换：
$$
\xymatrix{
j_!\mathcal{M}^n \otimes_\mathcal{O} \mathcal{A}^m
\ar[r]_{g^n \otimes 1} \ar[d] &
\mathcal{N}^{n + d} \otimes_\mathcal{O} \mathcal{A}_U^m \ar[d] \\
j_!\mathcal{M}^{n + m} \ar[r]^{g^{n + m}} &
\mathcal{N}^{n + m + d}
}
$$
不过，我们知道
\begin{align*}
\Hom_{\mathcal{O}_U}(\mathcal{M}^n \otimes_{\mathcal{O}_U} \mathcal{A}_U^m,
j^*\mathcal{N}^{n + d + m})
& =
\Hom_\mathcal{O}(j_!(\mathcal{M}^n \otimes_{\mathcal{O}_U} \mathcal{A}_U^m),
\mathcal{N}^{n + d + m}) \\
& =
\Hom_\mathcal{O}(j_!\mathcal{M}^n \otimes_\mathcal{O} \mathcal{A}^m,
\mathcal{N}^{n + d + m})
\end{align*}
，这由上面已经使用的《位点上的模》引理
\ref{sites-modules-lemma-j-shriek-and-tensor} 得到。略去如下验证：
第一个群中第一幅图交换的障碍，映到最后一个群中第二幅图交换的障碍。
\end{proof}

\begin{lemma}
\label{sdga-lemma-tensor-with-extension-by-zero}
在上述情形下，设 $\mathcal{M}$ 是右分次 $\mathcal{A}_U$-模，
$\mathcal{N}$ 是左分次 $\mathcal{A}$-模。则
$$
j_!\mathcal{M} \otimes_\mathcal{A} \mathcal{N} =
j_!(\mathcal{M} \otimes_{\mathcal{A}_U} \mathcal{N}|_U)
$$
作为分次 $\mathcal{O}$-模的等式，并且关于 $\mathcal{M}$ 和
$\mathcal{N}$ 具有函子性。
\end{lemma}

\begin{proof}
回忆，$j_!\mathcal{M}\otimes_\mathcal{A}\mathcal{N}$ 的次数 $n$
分量是下列典范映射的余核：
$$
\bigoplus\nolimits_{r + s + t = n}
j_!\mathcal{M}^r \otimes_\mathcal{O}
\mathcal{A}^s \otimes_\mathcal{O}
\mathcal{N}^t
\longrightarrow
\bigoplus\nolimits_{p + q = n}
j_!\mathcal{M}^p \otimes_\mathcal{O} \mathcal{N}^q
$$
。参见第 \ref{sdga-section-tensor-product} 节。由《位点上的模》引理
\ref{sites-modules-lemma-j-shriek-and-tensor}，这等于下列映射的余核：
$$
\bigoplus\nolimits_{r + s + t = n}
j_!(\mathcal{M}^r \otimes_{\mathcal{O}_U}
\mathcal{A}^s|_U \otimes_{\mathcal{O}_U}
\mathcal{N}^t|_U)
\longrightarrow
\bigoplus\nolimits_{p + q = n}
j_!(\mathcal{M}^p \otimes_{\mathcal{O}_U} \mathcal{N}^q|_U)
$$
，结论即得。另一种证明是重做上述引理证明中的米田论证。
\end{proof}








\section{分次模层上的平移函子}
\label{sdga-section-shift}

\noindent
我们强烈建议读者跳过本节。分次代数上的分次模层正是《微分分次代数》
评注 \ref{dga-remark-graded-shift-functors} 所讨论现象的一个例子。

\medskip\noindent
设 $(\mathcal{C},\mathcal{O})$ 是环化位点，$\mathcal{A}$ 是
$(\mathcal{C},\mathcal{O})$ 上的分次代数层，$\mathcal{M}$ 是分次
$\mathcal{A}$-模，$k\in\mathbf{Z}$。定义 $\mathcal{M}$ 的
{\it 第 $k$ 次平移}，记作 $\mathcal{M}[k]$，为这样的分次
$\mathcal{A}$-模：其第 $n$ 个分次部分为
$$
(\mathcal{M}[k])^n = \mathcal{M}^{n + k}
$$
，即 $\mathcal{M}$ 的第 $n+k$ 个分次部分。它的乘法映射
$$
(\mathcal{M}[k])^n \times \mathcal{A}^m
\longrightarrow
(\mathcal{M}[k])^{n + m}
$$
直接取为 $\mathcal{M}$ 的乘法映射
$$
\mathcal{M}^{n + k} \times \mathcal{A}^m
\longrightarrow
\mathcal{M}^{n + m + k}
$$
。显然，这定义了函子 $[k]$，并且 $[k+l]=[k]\circ[l]$，还有
$$
\Hom_{\textit{Mod}^{gr}(\mathcal{A})}(\mathcal{L}, \mathcal{M}[k]) =
\Hom_{\textit{Mod}^{gr}(\mathcal{A})}(\mathcal{L}, \mathcal{M})[k]
$$
（不引入符号），并且关于 $\mathcal{M}$ 和 $\mathcal{L}$ 具有函子性。
因此，正如《微分分次代数》评注
\ref{dga-remark-graded-shift-functors} 所讨论的那样，分次
$\mathcal{A}$-模的分次范畴确实可以从分次 $\mathcal{A}$-模的通常范畴
和平移函子恢复出来。

\begin{lemma}
\label{sdga-lemma-gm-grothendieck-abelian}
设 $(\mathcal{C},\mathcal{O})$ 是环化位点，$\mathcal{A}$ 是分次
$\mathcal{O}$-代数。范畴 $\textit{Mod}(\mathcal{A})$
是格罗滕迪克阿贝尔范畴。
\end{lemma}

\begin{proof}
由引理 \ref{sdga-lemma-gm-abelian} 以及格罗滕迪克阿贝尔范畴的定义
（《单射对象》定义 \ref{injectives-definition-grothendieck-conditions}），
只须证明 $\textit{Mod}(\mathcal{A})$ 有生成元。我们断言
$$
\mathcal{G} = \bigoplus\nolimits_{k, U} j_{U!}\mathcal{A}_U[k]
$$
是一个生成元，其中直和遍历 $\mathcal{C}$ 的所有对象 $U$ 和
$k\in\mathbf{Z}$。事实上，给定分次 $\mathcal{A}$-模 $\mathcal{M}$，
若从 $\mathcal{G}$ 到 $\mathcal{M}$ 没有非零映射，则对所有
$k$ 和 $U$ 都有
$$
\Hom_{\textit{Mod}(\mathcal{A})}(j_{U!}\mathcal{A}_U[k], \mathcal{M}) =
\Hom_{\textit{Mod}(\mathcal{A}_U)}(\mathcal{A}_U[k], \mathcal{M}|_U) =
\Gamma(U, \mathcal{M}^{-k})
$$
等于零。因此 $\mathcal{M}$ 为零。
\end{proof}









\section{微分分次代数层}
\label{sdga-section-dga}

\noindent
本节是《微分分次代数》第 \ref{dga-section-dga} 节的类似物。

\begin{definition}
\label{sdga-definition-dga}
设 $(\mathcal{C},\mathcal{O})$ 是环化位点。$(\mathcal{C},\mathcal{O})$
上的{\it 微分分次 $\mathcal{O}$-代数层}或{\it 微分分次代数层}，
是一个 $\mathcal{O}$-模上链复形 $\mathcal{A}^\bullet$，并带有
$\mathcal{O}$-双线性映射
$$
\mathcal{A}^n \times \mathcal{A}^m \to \mathcal{A}^{n + m},\quad
(a, b) \longmapsto ab
$$
；这些映射称为乘法映射，并满足下列性质：
\begin{enumerate}
\item 乘法是结合的；
\item $\mathcal{A}^0$ 有一个全局截面 $1$，它是乘法的双边单位元；
\item 对 $U\in\Ob(\mathcal{C})$、$a\in\mathcal{A}^n(U)$ 和
$b\in\mathcal{A}^m(U)$，有
$$
\text{d}^{n + m}(ab) = \text{d}^n(a)b + (-1)^n a\text{d}^m(b)
$$
\end{enumerate}
我们常把这样的结构记作 $(\mathcal{A},\text{d})$。从
$(\mathcal{A},\text{d})$ 到 $(\mathcal{B},\text{d})$ 的
{\it 微分分次 $\mathcal{O}$-代数同态}，是与乘法映射相容的
$\mathcal{O}$-模复形映射
$f:\mathcal{A}^\bullet\to\mathcal{B}^\bullet$。
\end{definition}

\noindent
给定微分分次 $\mathcal{O}$-代数 $(\mathcal{A},\text{d})$ 和对象
$U\in\Ob(\mathcal{C})$，我们使用记号
$$
\mathcal{A}(U) =
\Gamma(U, \mathcal{A}) =
\bigoplus\nolimits_{n \in \mathbf{Z}} \mathcal{A}^n(U)
$$
这是微分分次 $\mathcal{O}(U)$-代数。

\medskip\noindent
我们将尽可能把微分分次 $\mathcal{O}$-代数
$(\mathcal{A},\text{d})$ 视为一个分次 $\mathcal{O}$-代数
$\mathcal{A}$，连同满足定义中莱布尼茨法则的次数 $1$ 算子
$\text{d}:\mathcal{A}\to\mathcal{A}$（这里把 $\mathcal{A}$
视为分次 $\mathcal{O}$-模）。

\begin{remark}
\label{sdga-remark-functoriality-dga}
设 $(f,f^\sharp):(\Sh(\mathcal{C}),\mathcal{O}_\mathcal{C})
\to (\Sh(\mathcal{D}), \mathcal{O}_\mathcal{D})$
是环化拓扑斯的态射。
\begin{enumerate}
\item 设 $(\mathcal{A},\text{d})$ 是微分分次
$\mathcal{O}_\mathcal{C}$-代数。其直像是微分分次
$\mathcal{O}_\mathcal{D}$-代数 $(f_*\mathcal{A},\text{d})$，
其中 $f_*\mathcal{A}$ 如评注 \ref{sdga-remark-functoriality-ga} 所述，
而映射 $f_*\mathcal{A}^n\to f_*\mathcal{A}^{n+1}$ 上的
$\text{d}=f_*\text{d}$。略去莱布尼茨法则成立的验证。
\item 设 $\mathcal{B}$ 是微分分次 $\mathcal{O}_\mathcal{D}$-代数。
其拉回是微分分次 $\mathcal{O}_\mathcal{C}$-代数
$(f^*\mathcal{B},\text{d})$，其中 $f^*\mathcal{B}$ 如评注
\ref{sdga-remark-functoriality-ga} 所述，而映射
$f^*\mathcal{B}^n\to f^*\mathcal{B}^{n+1}$ 上的
$\text{d}=f^*\text{d}$。略去莱布尼茨法则成立的验证。
\item 微分分次 $\mathcal{O}_\mathcal{C}$-代数同态
$f^*\mathcal{B}\to\mathcal{A}$ 的集合，与微分分次
$\mathcal{O}_\mathcal{D}$-代数同态
$\mathcal{B}\to f_*\mathcal{A}$ 的集合一一对应。
\end{enumerate}
第 (3) 条立即由模层上的 $f^*$ 与 $f_*$ 之间通常的伴随得到。
\end{remark}
















\section{微分分次模层}
\label{sdga-section-modules}

\noindent
本节是《微分分次代数》第 \ref{dga-section-modules} 节的类似物。

\begin{definition}
\label{sdga-definition-dgm}
设 $(\mathcal{C},\mathcal{O})$ 是环化位点，
$(\mathcal{A},\text{d})$ 是 $(\mathcal{C},\mathcal{O})$ 上的
微分分次代数层。一个（右）{\it 微分分次 $\mathcal{A}$-模}，
或 $\mathcal{A}$ 上的（右）{\it 微分分次模}，是一个上链复形
$\mathcal{M}^\bullet$，并带有 $\mathcal{O}$-双线性映射
$$
\mathcal{M}^n \times \mathcal{A}^m \to \mathcal{M}^{n + m},\quad
(x, a) \longmapsto xa
$$
；这些映射称为乘法映射，并满足下列性质：
\begin{enumerate}
\item 乘法满足 $(xa)a'=x(aa')$；
\item 对每个 $n$，$\mathcal{A}^0$ 的单位截面 $1$ 在
$\mathcal{M}^n$ 上作恒等作用；
\item 对 $U\in\Ob(\mathcal{C})$、$x\in\mathcal{M}^n(U)$ 和
$a\in\mathcal{A}^m(U)$，有
$$
\text{d}^{n + m}(xa) = \text{d}^n(x)a + (-1)^n x\text{d}^m(a)
$$
\end{enumerate}
我们常说“设 $\mathcal{M}$ 是微分分次 $\mathcal{A}$-模”来表示上述情形。
从 $\mathcal{M}$ 到 $\mathcal{N}$ 的
{\it 微分分次 $\mathcal{A}$-模同态}，是与乘法映射相容的
$\mathcal{O}$-模复形映射
$f:\mathcal{M}^\bullet\to\mathcal{N}^\bullet$。
（右）微分分次 $\mathcal{A}$-模构成的范畴记作
$\textit{Mod}(\mathcal{A},\text{d})$。
\end{definition}

\noindent
可以完全类似地定义{\it 左微分分次模}，但本章默认所说的模都是右模。

\medskip\noindent
给定微分分次 $\mathcal{A}$-模 $\mathcal{M}$ 和对象
$U\in\Ob(\mathcal{C})$，我们使用记号
$$
\mathcal{M}(U) =
\Gamma(U, \mathcal{M}) =
\bigoplus\nolimits_{n \in \mathbf{Z}} \mathcal{M}^n(U)
$$
这是一个（右）微分分次 $\mathcal{A}(U)$-模。

\begin{lemma}
\label{sdga-lemma-dgm-abelian}
设 $(\mathcal{C},\mathcal{O})$ 是环化位点，
$(\mathcal{A},\text{d})$ 是微分分次 $\mathcal{O}$-代数。
范畴 $\textit{Mod}(\mathcal{A},\text{d})$ 是阿贝尔范畴，并具有下列性质：
\begin{enumerate}
\item $\textit{Mod}(\mathcal{A},\text{d})$ 有任意直和；
\item $\textit{Mod}(\mathcal{A},\text{d})$ 有任意余极限；
\item $\textit{Mod}(\mathcal{A},\text{d})$ 中的滤过余极限是正合的；
\item $\textit{Mod}(\mathcal{A},\text{d})$ 有任意积；
\item $\textit{Mod}(\mathcal{A},\text{d})$ 有任意极限。
\end{enumerate}
遗忘函子
$$
\textit{Mod}(\mathcal{A}, \text{d})
\longrightarrow
\textit{Mod}(\mathcal{A})
$$
把微分分次 $\mathcal{A}$-模映到其底层分次模；它与所有极限和余极限交换。
\end{lemma}

\begin{proof}
把引理陈述中的函子记作
$F : \textit{Mod}(\mathcal{A}, \text{d}) \to \textit{Mod}(\mathcal{A})$
。注意，范畴 $\textit{Mod}(\mathcal{A})$ 具有性质 (1)--(5)；
见引理 \ref{sdga-lemma-gm-abelian}。

\medskip\noindent
考虑分次 $\mathcal{A}$-模同态 $f:\mathcal{M}\to\mathcal{N}$。
把 $f$ 看成分次 $\mathcal{A}$-模映射时，其核与余核还带有定义
\ref{sdga-definition-dgm} 中的微分。因此，它们也是
$\textit{Mod}(\mathcal{A},\text{d})$ 中的核与余核。由此
$\textit{Mod}(\mathcal{A},\text{d})$ 是阿贝尔范畴，并且取核与余核和
$F$ 交换。

\medskip\noindent
为了证明极限和余极限存在，只须证明积和直和存在；参见《范畴》引理
\ref{categories-lemma-limits-products-equalizers} 和
\ref{categories-lemma-colimits-coproducts-coequalizers}。
同样的引理表明，只要 $F$ 与直和及积交换，就可推出它与极限及余极限交换。

\medskip\noindent
设 $\mathcal{M}_t$（$t\in T$）是一族微分分次 $\mathcal{A}$-模。
把直和 $\bigoplus\mathcal{M}_t$ 看成分次 $\mathcal{A}$-模。
由于分次模的直和是逐项作出的，显然
$\bigoplus\mathcal{M}_t$ 带有微分。读者容易验证，这是
$\textit{Mod}(\mathcal{A},\text{d})$ 中的直和。积的情形类似。

\medskip\noindent
注意 $F$ 是正合函子，并且 $\textit{Mod}(\mathcal{A},\text{d})$
中的复形 $\mathcal{M}_1\to\mathcal{M}_2\to\mathcal{M}_3$ 正合，当且仅当
$F(\mathcal{M}_1) \to F(\mathcal{M}_2) \to F(\mathcal{M}_3)$
在 $\textit{Mod}(\mathcal{A})$ 中正合。由于
$\textit{Mod}(\mathcal{A})$ 中的滤过余极限正合，由此得到 (3)。
\end{proof}

\noindent
结合引理 \ref{sdga-lemma-dgm-abelian} 和 \ref{sdga-lemma-gm-abelian}，
得到阿贝尔范畴之间的正合忠实函子
$$
\textit{Mod}(\mathcal{A}, \text{d}) \longrightarrow \text{Comp}(\mathcal{O})
$$
。对微分分次 $\mathcal{A}$-模 $\mathcal{M}$，定义其上同调
$\mathcal{O}$-模 $H^i(\mathcal{M})$ 为与 $\mathcal{M}$ 对应的
$\mathcal{O}$-模复形的上同调。因此，微分分次 $\mathcal{A}$-模的短正合列
$0 \to \mathcal{K} \to \mathcal{L} \to \mathcal{M} \to 0$
给出上同调模的长正合列
\begin{equation}
\label{sdga-equation-les}
H^n(\mathcal{K}) \to H^n(\mathcal{L}) \to H^n(\mathcal{M}) \to
H^{n + 1}(\mathcal{K})
\end{equation}
；见《同调代数》引理 \ref{homology-lemma-long-exact-sequence-cochain}。

\medskip\noindent
此外，从现在起沿用模复形的全部术语。例如，若对所有
$k\in\mathbf{Z}$ 都有 $H^k(\mathcal{M})=0$，则称微分分次
$\mathcal{A}$-模 $\mathcal{M}$ 是{\it 无环的}。若微分分次
$\mathcal{A}$-模同态 $\mathcal{M}\to\mathcal{N}$ 对所有
$k\in\mathbf{Z}$ 都诱导同构
$H^k(\mathcal{M})\to H^k(\mathcal{N})$，则称它为{\it 拟同构}。
其余依此类推。


















\section{模的微分分次范畴}
\label{sdga-section-dgm-dg-cat}

\noindent
本节是《微分分次代数》例 \ref{dga-example-dgm-dg-cat} 的类似物。
关于微分分次范畴的约定，请参见《微分分次代数》第
\ref{dga-section-dga-categories} 节。

\medskip\noindent
设 $(\mathcal{C},\mathcal{O})$ 是环化位点，
$(\mathcal{A},\text{d})$ 是 $(\mathcal{C},\mathcal{O})$ 上的
微分分次代数层。我们将构造微分分次范畴
$$
\textit{Mod}^{dg}(\mathcal{A}, \text{d})
$$
，它在 $R=\Gamma(\mathcal{C},\mathcal{O})$ 上，且其对应的复形范畴
是微分分次 $\mathcal{A}$-模范畴：
$$
\textit{Mod}(\mathcal{A}, \text{d}) =
\text{Comp}(\textit{Mod}^{dg}(\mathcal{A}, \text{d}))
$$
。取右微分分次 $\mathcal{A}$-模作为
$\textit{Mod}^{dg}(\mathcal{A},\text{d})$ 的对象；见第
\ref{sdga-section-modules} 节。给定微分分次 $\mathcal{A}$-模
$\mathcal{L}$ 和 $\mathcal{M}$，置
$$
\Hom_{\textit{Mod}^{dg}(\mathcal{A}, \text{d})}(\mathcal{L}, \mathcal{M}) =
\Hom_{\textit{Mod}^{gr}(\mathcal{A})}(\mathcal{L}, \mathcal{M}) =
\bigoplus\nolimits_{n \in \mathbf{Z}} \Hom^n(\mathcal{L}, \mathcal{M})
$$
，把它作为分次 $R$-模；见第 \ref{sdga-section-gm-gr-cat} 节。
换言之，第 $n$ 个分次部分 $\Hom^n(\mathcal{L},\mathcal{M})$
是次数 $n$ 齐次的右 $\mathcal{A}$-模映射所成的 $R$-模。
对元素 $f\in\Hom^n(\mathcal{L},\mathcal{M})$，置
$$
\text{d}(f) =
\text{d}_\mathcal{M} \circ f - (-1)^n f \circ \text{d}_\mathcal{L}
$$
。为使此式有意义，把 $\text{d}_\mathcal{M}$ 和
$\text{d}_\mathcal{L}$ 看成分次 $\mathcal{O}$-模映射，并使用分次
$\mathcal{O}$-模映射的复合。显然，$\text{d}(f)$ 作为分次
$\mathcal{O}$-模映射是次数 $n+1$ 齐次的，并且关于
$\mathcal{A}$ 是线性的，因为对 $\mathcal{M}$ 和 $\mathcal{A}$
的齐次局部截面 $x$ 和 $a$，有
\begin{align*}
\text{d}(f)(xa)
& =
\text{d}_\mathcal{M}(f(x) a) - (-1)^n f (\text{d}_\mathcal{L}(xa)) \\
& =
\text{d}_\mathcal{M}(f(x)) a + (-1)^{\deg(x) + n} f(x) \text{d}(a) 
- (-1)^n f(\text{d}_\mathcal{L}(x)) a - (-1)^{n + \deg(x)} f(x) \text{d}(a) \\
& = \text{d}(f)(x) a
\end{align*}
，如所需（注意，若 $\Hom$ 上的微分定义中没有该符号，这一计算便不成立）。

\medskip\noindent
对微分分次 $\mathcal{A}$-模 $\mathcal{K}$、$\mathcal{L}$、
$\mathcal{M}$，第 \ref{sdga-section-gm-gr-cat} 节已经通过层映射的通常复合
定义了复合
$$
\Hom^m(\mathcal{L}, \mathcal{M}) \times
\Hom^n(\mathcal{K}, \mathcal{L}) \longrightarrow
\Hom^{n + m}(\mathcal{K}, \mathcal{M})
$$
。这定义微分分次模之间的映射
$$
\Hom_{\textit{Mod}^{dg}(\mathcal{A}, \text{d})}(\mathcal{L}, \mathcal{M})
\otimes_R
\Hom_{\textit{Mod}^{dg}(\mathcal{A}, \text{d})}(\mathcal{K}, \mathcal{L})
\longrightarrow
\Hom_{\textit{Mod}^{dg}(\mathcal{A}, \text{d})}(\mathcal{K}, \mathcal{M})
$$
，如《微分分次代数》定义 \ref{dga-definition-dga-category} 所要求，
因为
\begin{align*}
\text{d}(g \circ f) & =
\text{d}_\mathcal{M} \circ g \circ f
- (-1)^{n + m} g \circ f \circ \text{d}_\mathcal{K} \\
& =
\left(\text{d}_\mathcal{M} \circ g - (-1)^m g \circ \text{d}_L\right) \circ f
+ (-1)^m g \circ \left(\text{d}_\mathcal{L} \circ f
- (-1)^n f \circ \text{d}_\mathcal{K}\right) \\
& =
\text{d}(g) \circ f + (-1)^m g \circ \text{d}(f)
\end{align*}
，这里 $f$ 的次数为 $n$，$g$ 的次数为 $m$，正如所需。






\section{微分分次模层的张量积}
\label{sdga-section-tensor-product-dg}

\noindent
本节是《微分分次代数》第
\ref{dga-section-tensor-product} 节部分内容的类似物。

\medskip\noindent
设 $(\mathcal{C},\mathcal{O})$ 是环化位点，
$(\mathcal{A},\text{d})$ 是 $(\mathcal{C},\mathcal{O})$ 上的
微分分次代数层。设 $\mathcal{M}$ 是右微分分次 $\mathcal{A}$-模，
$\mathcal{N}$ 是左微分分次 $\mathcal{A}$-模。在这种情形下，
如下定义{\it 张量积} $\mathcal{M}\otimes_\mathcal{A}\mathcal{N}$。
作为分次 $\mathcal{O}$-模，它由第 \ref{sdga-section-tensor-product} 节的
构造给出，并带有微分
$$
\text{d}_{\mathcal{M} \otimes_\mathcal{A} \mathcal{N}} :
(\mathcal{M} \otimes_\mathcal{A} \mathcal{N})^n
\longrightarrow
(\mathcal{M} \otimes_\mathcal{A} \mathcal{N})^{n + 1}
$$
，其定义规则为
$$
\text{d}_{\mathcal{M} \otimes_\mathcal{A} \mathcal{N}}(x \otimes y) =
\text{d}_\mathcal{M}(x) \otimes y +
(-1)^{\deg(x)}x \otimes \text{d}_\mathcal{N}(y)
$$
，其中 $x$ 和 $y$ 分别是 $\mathcal{M}$ 和 $\mathcal{N}$ 的齐次局部截面。
为了看出其良定义性，必须证明
$\text{d}_{\mathcal{M}\otimes_\mathcal{A}\mathcal{N}}$ 把形如
$xa\otimes y-x\otimes ay$ 的元素映到零；这里 $x$、$a$、$y$
分别是 $\mathcal{M}$、$\mathcal{A}$、$\mathcal{N}$ 的齐次局部截面。
计算得
\begin{align*}
&
\text{d}_{\mathcal{M} \otimes_\mathcal{A} \mathcal{N}}(
xa \otimes y - x \otimes ay) \\
& =
\text{d}_\mathcal{M}(xa) \otimes y + (-1)^{\deg(x) + \deg(a)}
xa \otimes \text{d}_\mathcal{N}(y)
-\text{d}_\mathcal{M}(x) \otimes ay - (-1)^{\deg(x)}
x \otimes \text{d}_\mathcal{N}(ay) \\
& =
\text{d}_\mathcal{M}(x)a \otimes y + (-1)^{\deg(x)}x\text{d}(a) \otimes y
+ (-1)^{\deg(x) + \deg(a)} xa \otimes \text{d}_\mathcal{N}(y) \\
&
-\text{d}_\mathcal{M}(x) \otimes ay - (-1)^{\deg(x)}
x \otimes \text{d}(a)y - (-1)^{\deg(x) +
\deg(a)} x\otimes a\text{d}_\mathcal{N}(y)
\end{align*}
。注意下列元素
$$
\text{d}_\mathcal{M}(x)a \otimes y - \text{d}_\mathcal{M}(x) \otimes ay,\quad
x\text{d}(a) \otimes y - x \otimes \text{d}(a)y,\quad
\text{且}\quad
xa \otimes \text{d}_\mathcal{N}(y) - x\otimes a\text{d}_\mathcal{N}(y)
$$
在 $\mathcal{M}\otimes_\mathcal{A}\mathcal{N}$ 中都映到零，故得结论。
略去下式的验证：
$\text{d}_{\mathcal{M} \otimes_\mathcal{A} \mathcal{N}} \circ
\text{d}_{\mathcal{M} \otimes_\mathcal{A} \mathcal{N}} = 0$.

\medskip\noindent
固定左微分分次 $\mathcal{A}$-模 $\mathcal{N}$，便得到函子
$$
- \otimes_\mathcal{A} \mathcal{N} :
\textit{Mod}(\mathcal{A}, \text{d})
\longrightarrow
\text{Comp}(\mathcal{O})
$$
，其右端是 $\mathcal{O}$-模复形范畴。它可以提升为微分分次范畴之间的函子
$$
- \otimes_\mathcal{A} \mathcal{N} :
\textit{Mod}^{dg}(\mathcal{A}, \text{d})
\longrightarrow
\text{Comp}^{dg}(\mathcal{O})
$$
。在底层分次对象上，把次数 $n$ 的同态
$f:\mathcal{M}\to\mathcal{M}'$ 映到次数 $n$ 的映射
$f \otimes \text{id}_\mathcal{N} :
\mathcal{M} \otimes_\mathcal{A} \mathcal{N} \to
\mathcal{M}' \otimes_\mathcal{A} \mathcal{N}$，因为这正是第
\ref{sdga-section-tensor-product} 节中的做法。为证明该构造成立，
必须验证映射
$$
\Hom_{\textit{Mod}^{dg}(\mathcal{A}, \text{d})}(\mathcal{M}, \mathcal{M}')
\longrightarrow
\Hom_{\text{Comp}^{dg}(\mathcal{O})}(
\mathcal{M} \otimes_\mathcal{A} \mathcal{N},
\mathcal{M}' \otimes_\mathcal{A} \mathcal{N})
$$
与微分相容。对上述 $f$，须证明
$$
(\text{d}_{\mathcal{M}'} \circ f - (-1)^n f \circ \text{d}_\mathcal{M})
\otimes \text{id}_\mathcal{N}
$$
等于
$$
\text{d}_{\mathcal{M}' \otimes_\mathcal{A} \mathcal{N}}
\circ (f \otimes \text{id}_\mathcal{N})
- (-1)^n (f \otimes \text{id}_\mathcal{N}) \circ
\text{d}_{\mathcal{M} \otimes_\mathcal{A} \mathcal{N}}
$$
。计算这些算子在形如 $x\otimes y$ 的局部截面上的作用，其中
$x$ 和 $y$ 分别是 $\mathcal{M}$ 和 $\mathcal{N}$ 的齐次局部截面。
第一个算子给出
$$
(\text{d}_{\mathcal{M}'}(f(x)) - (-1)^n f(\text{d}_\mathcal{M}(x))) \otimes y
$$
，第二个算子给出
\begin{align*}
&\text{d}_{\mathcal{M}' \otimes_\mathcal{A} \mathcal{N}}(f(x) \otimes y)
- (-1)^n (f \otimes \text{id}_\mathcal{N})(
\text{d}_{\mathcal{M} \otimes_\mathcal{A} \mathcal{N}}(x \otimes y) \\
& =
\text{d}_{\mathcal{M}'}(f(x)) \otimes y +
(-1)^{\deg(x) + n}f(x) \otimes \text{d}_\mathcal{N}(y) \\
&
-(-1)^n f(\text{d}_\mathcal{M}(x)) \otimes y
-(-1)^n (-1)^{\deg(x)}f(x) \otimes \text{d}_\mathcal{N}(y)
\end{align*}
，这确实是同一个局部截面。










\section{微分分次模层的内 Hom}
\label{sdga-section-internal-hom-dg}

\noindent
我们需要第 \ref{sdga-section-dgm-dg-cat} 节中构造的层化版本。
设 $(\mathcal{C},\mathcal{O})$、$\mathcal{A}$、
$\mathcal{M}$、$\mathcal{L}$ 如第 \ref{sdga-section-dgm-dg-cat} 节所述。定义
$$
\SheafHom^{dg}_\mathcal{A}(\mathcal{M}, \mathcal{L}) =
\SheafHom^{gr}_\mathcal{A}(\mathcal{M}, \mathcal{L}) =
\bigoplus\nolimits_{n \in \mathbf{Z}}
\SheafHom^n_\mathcal{A}(\mathcal{M}, \mathcal{L})
$$
为分次 $\mathcal{O}$-模；见第
\ref{sdga-section-internal-hom-graded} 节。换言之，$U$ 上第 $n$ 个
分次部分 $\SheafHom^n_\mathcal{A}(\mathcal{L},\mathcal{M})$
的截面 $f$，是次数 $n$ 齐次的右 $\mathcal{A}_U$-模映射
$\mathcal{L}|_U\to\mathcal{M}|_U$。对这样的 $f$，置
$$
\text{d}(f) =
\text{d}_\mathcal{M}|_U \circ f - (-1)^n f \circ \text{d}_\mathcal{L}|_U
$$
。为使此式有意义，把 $\text{d}_\mathcal{M}|_U$ 和
$\text{d}_\mathcal{L}|_U$ 看成分次 $\mathcal{O}_U$-模映射，
并使用分次 $\mathcal{O}_U$-模映射的复合。显然，
$\text{d}(f)$ 作为分次 $\mathcal{O}_U$-模映射是次数 $n+1$
齐次的。使用与第 \ref{sdga-section-dgm-dg-cat} 节完全相同的计算，
可见 $\text{d}(f)$ 关于 $\mathcal{A}_U$ 是线性的。

\medskip\noindent
如同第 \ref{sdga-section-dgm-dg-cat} 节，有复合映射
$$
\SheafHom^{dg}_\mathcal{A}(\mathcal{L}, \mathcal{M}) \otimes_\mathcal{O}
\SheafHom^{dg}_\mathcal{A}(\mathcal{K}, \mathcal{L})
\longrightarrow
\SheafHom^{dg}_\mathcal{A}(\mathcal{K}, \mathcal{M})
$$
，其中左端是第 \ref{sdga-section-tensor-product-dg} 节定义的微分分次
$\mathcal{O}$-模张量积。该映射由复合映射
$$
\SheafHom^m(\mathcal{L}, \mathcal{M}) \otimes_\mathcal{O}
\SheafHom^n(\mathcal{K}, \mathcal{L}) \longrightarrow
\SheafHom^{n + m}(\mathcal{K}, \mathcal{M})
$$
给出；后者（局部地）按通常复合定义。在局部截面上使用与第
\ref{sdga-section-dgm-dg-cat} 节完全相同的计算，可见该复合映射是
微分分次 $\mathcal{O}$-模同态。

\medskip\noindent
按这些定义，有
$$
\Hom_{\textit{Mod}^{dg}(\mathcal{A})}(\mathcal{L}, \mathcal{M}) =
\Gamma(\mathcal{C}, \SheafHom^{dg}_\mathcal{A}(\mathcal{L}, \mathcal{M}))
$$
作为分次 $R$-模的等式，并且它与复合相容。









\section{微分分次双模层与张量--Hom 伴随}
\label{sdga-section-dg-bimodules}

\noindent
本节是《微分分次代数》第
\ref{dga-section-tensor-product} 节部分内容的类似物。

\begin{definition}
\label{sdga-definition-dg-bimodule}
设 $(\mathcal{C},\mathcal{O})$ 是环化位点，$\mathcal{A}$ 和
$\mathcal{B}$ 是 $(\mathcal{C},\mathcal{O})$ 上的微分分次代数层。
一个{\it 微分分次 $(\mathcal{A},\mathcal{B})$-双模}，由
$\mathcal{O}$-模复形 $\mathcal{M}^\bullet$ 以及
$\mathcal{O}$-双线性映射
$$
\mathcal{M}^n \times \mathcal{B}^m \to \mathcal{M}^{n + m},\quad
(x, b) \longmapsto xb
$$
和
$$
\mathcal{A}^n \times \mathcal{M}^m \to \mathcal{M}^{n + m},\quad
(a, x) \longmapsto ax
$$
给出；这些映射称为乘法映射，并满足下列性质：
\begin{enumerate}
\item 乘法满足 $a(a'x)=(aa')x$ 和 $(xb)b'=x(bb')$；
\item $(ax)b=a(xb)$；
\item $\text{d}(ax) = \text{d}(a) x + (-1)^{\deg(a)}a \text{d}(x)$ 和
$\text{d}(xb) = \text{d}(x) b + (-1)^{\deg(x)}x \text{d}(b)$；
\item $\mathcal{A}^0$ 的单位截面 $1$ 通过乘法作恒等作用；
\item $\mathcal{B}^0$ 的单位截面 $1$ 通过乘法作恒等作用。
\end{enumerate}
我们常把这样的结构记作 $\mathcal{M}$，有时写成
${}_\mathcal{A}\mathcal{M}_\mathcal{B}$。一个
{\it 微分分次 $(\mathcal{A},\mathcal{B})$-双模同态}
$f:\mathcal{M}\to\mathcal{N}$，是与乘法映射相容的
$\mathcal{O}$-模复形映射
$f:\mathcal{M}^\bullet\to\mathcal{N}^\bullet$。
\end{definition}

\noindent
给定微分分次 $(\mathcal{A},\mathcal{B})$-双模 $\mathcal{M}$
和对象 $U\in\Ob(\mathcal{C})$，我们使用记号
$$
\mathcal{M}(U) =
\Gamma(U, \mathcal{M}) =
\bigoplus\nolimits_{n \in \mathbf{Z}} \mathcal{M}^n(U)
$$
这是微分分次 $(\mathcal{A}(U),\mathcal{B}(U))$-双模。

\medskip\noindent
注意，微分分次 $(\mathcal{A},\mathcal{B})$-双模 $\mathcal{M}$
等价于一个右微分分次 $\mathcal{B}$-模，同时也是左微分分次
$\mathcal{A}$-模，且两者的分次和微分一致，并且
$\mathcal{A}$-模结构与 $\mathcal{B}$-模结构交换。精确陈述如下。

\begin{lemma}
\label{sdga-lemma-what-makes-a-bimodule-dg}
设 $(\mathcal{C},\mathcal{O})$ 是环化位点，$\mathcal{A}$ 和
$\mathcal{B}$ 是 $(\mathcal{C},\mathcal{O})$ 上的微分分次代数层，
$\mathcal{N}$ 是右微分分次 $\mathcal{B}$-模。$\mathcal{N}$ 上与给定
微分分次 $\mathcal{B}$-模结构相容的
$(\mathcal{A},\mathcal{B})$-双模结构，与微分分次
$\mathcal{O}$-代数同态
$$
\mathcal{A}
\longrightarrow
\SheafHom^{dg}_\mathcal{B}(\mathcal{N}, \mathcal{N})
$$
一一对应。
\end{lemma}

\begin{proof}
略。
\end{proof}

\noindent
设 $(\mathcal{C},\mathcal{O})$ 是环化位点，$\mathcal{A}$ 和
$\mathcal{B}$ 是 $(\mathcal{C},\mathcal{O})$ 上的微分分次代数层。
设 $\mathcal{M}$ 是右微分分次 $\mathcal{A}$-模，$\mathcal{N}$ 是
微分分次 $(\mathcal{A},\mathcal{B})$-双模。在这种情形下，第
\ref{sdga-section-tensor-product-dg} 节定义的微分分次张量积
$$
\mathcal{M} \otimes_\mathcal{A} \mathcal{N}
$$
是右微分分次 $\mathcal{B}$-模，其乘法映射如第
\ref{sdga-section-graded-bimodules} 节所述。该构造定义一个函子和一个
分次范畴之间的函子
$$
\otimes_\mathcal{A} \mathcal{N} :
\textit{Mod}(\mathcal{A}, \text{d})
\longrightarrow
\textit{Mod}(\mathcal{B}, \text{d})
\quad\text{且}\quad
\otimes_\mathcal{A} \mathcal{N} :
\textit{Mod}^{dg}(\mathcal{A}, \text{d})
\longrightarrow
\textit{Mod}^{dg}(\mathcal{B}, \text{d})
$$
：把从 $\mathcal{M}$ 到 $\mathcal{M}'$ 的次数 $n$ 同态映到
从 $\mathcal{M}\otimes_\mathcal{A}\mathcal{N}$ 到
$\mathcal{M}'\otimes_\mathcal{A}\mathcal{N}$ 的诱导次数 $n$ 映射。

\medskip\noindent
设 $(\mathcal{C},\mathcal{O})$ 是环化位点，$\mathcal{A}$ 和
$\mathcal{B}$ 是 $(\mathcal{C},\mathcal{O})$ 上的微分分次代数层。
设 $\mathcal{N}$ 是微分分次 $(\mathcal{A},\mathcal{B})$-双模，
$\mathcal{L}$ 是右微分分次 $\mathcal{B}$-模。在这种情形下，第
\ref{sdga-section-internal-hom-dg} 节定义的微分分次内 Hom
$$
\SheafHom_\mathcal{B}^{dg}(\mathcal{N}, \mathcal{L})
$$
是右微分分次 $\mathcal{A}$-模，其中右分次 $\mathcal{A}$-模结构
按第 \ref{sdga-section-graded-bimodules} 节定义。也可以用下列复合来定义乘法：
$$
\SheafHom_\mathcal{B}^{dg}(\mathcal{N}, \mathcal{L})
\otimes_\mathcal{O} \mathcal{A}
\to
\SheafHom_\mathcal{B}^{dg}(\mathcal{N}, \mathcal{L})
\otimes_\mathcal{O} \SheafHom_\mathcal{B}^{dg}(\mathcal{N}, \mathcal{N})
\to
\SheafHom_\mathcal{B}^{dg}(\mathcal{N}, \mathcal{L})
$$
其中第一个箭头来自引理 \ref{sdga-lemma-what-makes-a-bimodule-dg}，
第二个箭头是第 \ref{sdga-section-internal-hom-dg} 节的复合。
由于这两个箭头都与微分相容，所得对象确为微分分次
$\mathcal{A}$-模。该构造定义一个函子和一个微分分次范畴之间的函子
$$
\SheafHom_\mathcal{B}^{dg}(\mathcal{N}, -) :
\textit{Mod}(\mathcal{B}, \text{d})
\longrightarrow
\textit{Mod}(\mathcal{A})
\quad\text{且}\quad
\SheafHom_\mathcal{B}^{dg}(\mathcal{N}, -) :
\textit{Mod}^{dg}(\mathcal{B}, \text{d})
\longrightarrow
\textit{Mod}^{dg}(\mathcal{A}, \text{d})
$$
：把从 $\mathcal{L}$ 到 $\mathcal{L}'$ 的次数 $n$ 同态映到
从 $\SheafHom_\mathcal{B}^{dg}(\mathcal{N},\mathcal{L})$ 到
$\SheafHom_\mathcal{B}^{dg}(\mathcal{N},\mathcal{L}')$
的诱导次数 $n$ 映射。

\begin{lemma}
\label{sdga-lemma-tensor-hom-adjunction-dg}
设 $(\mathcal{C},\mathcal{O})$ 是环化位点，$\mathcal{A}$ 和
$\mathcal{B}$ 是 $(\mathcal{C},\mathcal{O})$ 上的微分分次代数层。
设 $\mathcal{M}$ 是右微分分次 $\mathcal{A}$-模，$\mathcal{N}$ 是
微分分次 $(\mathcal{A},\mathcal{B})$-双模，$\mathcal{L}$ 是右微分分次
$\mathcal{B}$-模。采用上述约定，有
$$
\Hom_{\textit{Mod}^{dg}(\mathcal{B}, \text{d})}(
\mathcal{M} \otimes_\mathcal{A} \mathcal{N}, \mathcal{L}) =
\Hom_{\textit{Mod}^{dg}(\mathcal{A}, \text{d})}(
\mathcal{M}, \SheafHom_\mathcal{B}^{dg}(\mathcal{N}, \mathcal{L}))
$$
以及
$$
\SheafHom_\mathcal{B}^{dg}(
\mathcal{M} \otimes_\mathcal{A} \mathcal{N}, \mathcal{L}) =
\SheafHom_\mathcal{A}^{dg}(
\mathcal{M}, \SheafHom_\mathcal{B}^{dg}(\mathcal{N}, \mathcal{L}))
$$
；这些等式关于 $\mathcal{M}$、$\mathcal{N}$、$\mathcal{L}$
都具有函子性。
\end{lemma}

\begin{proof}
略。提示：在分次层面，引理 \ref{sdga-lemma-tensor-hom-adjunction-gr}
已经证明了这一点。因此只须检验这些同构与微分相容；
这可通过局部截面层面的计算完成。
\end{proof}

\noindent
设 $(\mathcal{C},\mathcal{O})$ 是环化位点，$\mathcal{A}$ 和
$\mathcal{B}$ 是 $(\mathcal{C},\mathcal{O})$ 上的微分分次代数层。
作为上述情形的特例，设给定微分分次 $\mathcal{O}$-代数同态
$\varphi:\mathcal{A}\to\mathcal{B}$。于是得到一个函子和一个
微分分次范畴之间的函子
$$
\otimes_{\mathcal{A}, \varphi} \mathcal{B} :
\textit{Mod}(\mathcal{A}, \text{d})
\longrightarrow
\textit{Mod}(\mathcal{B}, \text{d})
\quad\text{且}\quad
\otimes_{\mathcal{A}, \varphi} \mathcal{B} :
\textit{Mod}^{dg}(\mathcal{A}, \text{d})
\longrightarrow
\textit{Mod}^{dg}(\mathcal{B}, \text{d})
$$
。另一方面，有限制标量函子
$$
res_\varphi :
\textit{Mod}(\mathcal{B}, \text{d})
\longrightarrow
\textit{Mod}(\mathcal{A}, \text{d})
\quad\text{且}\quad
res_\varphi :
\textit{Mod}^{dg}(\mathcal{B}, \text{d})
\longrightarrow
\textit{Mod}^{dg}(\mathcal{A}, \text{d})
$$
。由上一引理可以证明这些函子彼此伴随（如限制标量与基变换的通常情形）。
具体地，把 $\mathcal{B}$ 视为微分分次
$(\mathcal{A},\mathcal{B})$-双模时，记作
${}_\mathcal{A}\mathcal{B}_\mathcal{B}$。于是对任意右微分分次
$\mathcal{B}$-模 $\mathcal{L}$，有
$$
\SheafHom_\mathcal{B}^{dg}({}_\mathcal{A}\mathcal{B}_\mathcal{B}, \mathcal{L})
= res_\varphi(\mathcal{L})
$$
作为右微分分次 $\mathcal{A}$-模的等式。因此，引理
\ref{sdga-lemma-tensor-hom-adjunction-gr} 给出函子性同构
$$
\Hom_{\textit{Mod}^{dg}(\mathcal{B}, \text{d})}(
\mathcal{M} \otimes_{\mathcal{A}, \varphi} \mathcal{B}, \mathcal{L}) =
\Hom_{\textit{Mod}^{dg}(\mathcal{A}, \text{d})}(
\mathcal{M}, res_\varphi(\mathcal{L}))
$$
。若从上下文可知 $\varphi$，通常在此公式中省略对它的依赖。
以同样方式得到分次 $\mathcal{O}$-模的等式
$$
\SheafHom^{dg}_\mathcal{B}(
\mathcal{M} \otimes_\mathcal{A} \mathcal{B}, \mathcal{L}) =
\SheafHom_\mathcal{A}^{dg}(\mathcal{M}, \mathcal{L})
$$
。












\section{微分分次模层的拉回与直像}
\label{sdga-section-functoriality-dg}

\noindent
设 $(f,f^\sharp):(\Sh(\mathcal{C}),\mathcal{O}_\mathcal{C})
\to (\Sh(\mathcal{D}), \mathcal{O}_\mathcal{D})$
是环化拓扑斯的态射。设 $\mathcal{A}$ 是微分分次
$\mathcal{O}_\mathcal{C}$-代数，$\mathcal{B}$ 是微分分次
$\mathcal{O}_\mathcal{D}$-代数。设给定微分分次
$f^{-1}\mathcal{O}_\mathcal{D}$-代数同态
$$
\varphi : f^{-1}\mathcal{B} \to \mathcal{A}
$$
。由限制标量与扩张标量的伴随，这等价于微分分次
$\mathcal{O}_\mathcal{C}$-代数同态
$\varphi:f^*\mathcal{B}\to\mathcal{A}$；也等价地，可把
$\varphi$ 看作微分分次 $\mathcal{O}_\mathcal{D}$-代数同态
$$
\varphi : \mathcal{B} \to f_*\mathcal{A}
$$
。参见评注 \ref{sdga-remark-functoriality-dga}。

\medskip\noindent
先定义函子
$$
f_* :
\textit{Mod}(\mathcal{A}, \text{d})
\longrightarrow
\textit{Mod}(\mathcal{B}, \text{d})
$$
。给定微分分次 $\mathcal{A}$-模 $\mathcal{M}$，定义
$f_*\mathcal{M}$ 为第 \ref{sdga-section-functoriality-graded} 节构造的
分次 $\mathcal{B}$-模，其微分由映射
$f_*d:f_*\mathcal{M}^n\to f_*\mathcal{M}^{n+1}$ 给出。
该构造关于 $\mathcal{M}$ 显然具有函子性，从而得到所需函子。

\medskip\noindent
再定义函子
$$
f^* :
\textit{Mod}(\mathcal{B}, \text{d})
\longrightarrow
\textit{Mod}(\mathcal{A}, \text{d})
$$
。给定微分分次 $\mathcal{B}$-模 $\mathcal{N}$，定义
$f^*\mathcal{N}$ 为第 \ref{sdga-section-functoriality-graded} 节构造的
分次 $\mathcal{A}$-模。回忆
$$
f^*\mathcal{N} = f^{-1}\mathcal{N} \otimes_{f^{-1}\mathcal{B}} \mathcal{A}
$$
。由于 $f^{-1}\mathcal{N}$ 带有微分
$f^{-1}\text{d}:f^{-1}\mathcal{N}^n\to f^{-1}\mathcal{N}^{n+1}$，
可把该张量积视为第 \ref{sdga-section-dg-bimodules} 节所讨论张量积的例子，
因而得到微分。该构造关于 $\mathcal{N}$ 显然具有函子性，
从而得到函子 $f^*$。

\medskip\noindent
函子 $f_*$ 和 $f^*$ 很容易提升为微分分次范畴之间的函子
$$
f_* :
\textit{Mod}^{dg}(\mathcal{A}, \text{d})
\longrightarrow
\textit{Mod}^{dg}(\mathcal{B}, \text{d})
\quad\text{且}\quad
f^* :
\textit{Mod}^{dg}(\mathcal{B}, \text{d})
\longrightarrow
\textit{Mod}^{dg}(\mathcal{A}, \text{d})
$$
；它们在底层对象上的作用不变，而在次数 $n$ 的齐次态射上由上述构造的
函子性定义。

\begin{lemma}
\label{sdga-lemma-adjunction-push-pull-dg}
在上述情形下，有
$$
\Hom_{\textit{Mod}^{dg}(\mathcal{B}, \text{d})}(
\mathcal{N}, f_*\mathcal{M}) =
\Hom_{\textit{Mod}^{dg}(\mathcal{A}, \text{d})}(
f^*\mathcal{N}, \mathcal{M})
$$
\end{lemma}

\begin{proof}
略。提示：由引理 \ref{sdga-lemma-adjunction-push-pull-gr}，
该结论对底层分次范畴成立。一次计算表明这些同构与微分相容。
\end{proof}















\section{局部化与微分分次模层}
\label{sdga-section-localize-dg}

\noindent
设 $(\mathcal{C},\mathcal{O})$ 是环化位点，$U\in\Ob(\mathcal{C})$。
记
$$
j :
(\Sh(\mathcal{C}/U), \mathcal{O}_U)
\longrightarrow
(\Sh(\mathcal{C}), \mathcal{O})
$$
为相应的局部化态射（《位点上的模》第
\ref{sites-modules-section-localize} 节）。下面将使用如下事实：
对 $\mathcal{O}_U$-模 $\mathcal{M}_i$（$i=1,2$）和
$\mathcal{O}$-模 $\mathcal{A}$，有典范映射
$$
j_! :
\Hom_{\mathcal{O}_U}(
\mathcal{M}_1 \otimes_{\mathcal{O}_U} \mathcal{A}|_U, \mathcal{M}_2)
\longrightarrow
\Hom_\mathcal{O}(
j_!\mathcal{M}_1 \otimes_\mathcal{O} \mathcal{A}, j_!\mathcal{M}_2)
$$
。事实上，由《位点上的模》引理
\ref{sites-modules-lemma-j-shriek-and-tensor}，有
$j_!(\mathcal{M}_1 \otimes_{\mathcal{O}_U} \mathcal{A}|_U) =
j_!\mathcal{M}_1 \otimes_\mathcal{O} \mathcal{A}$。

\medskip\noindent
设 $\mathcal{A}$ 是微分分次 $\mathcal{O}$-代数。用
$\mathcal{A}_U$ 表示 $\mathcal{A}$ 在 $\mathcal{C}/U$ 上的限制；
换言之，$\mathcal{A}_U=j^*\mathcal{A}=j^{-1}\mathcal{A}$。
第 \ref{sdga-section-functoriality-dg} 节已经构造伴随函子
$$
j_* :
\textit{Mod}^{dg}(\mathcal{A}_U, \text{d})
\longrightarrow
\textit{Mod}^{dg}(\mathcal{A}, \text{d})
\quad\text{且}\quad
j^* :
\textit{Mod}^{dg}(\mathcal{A}, \text{d})
\longrightarrow
\textit{Mod}^{dg}(\mathcal{A}_U, \text{d})
$$
，其中 $j^*$ 是 $j_*$ 的左伴随。我们断言，此外还有正合函子
$$
j_! :
\textit{Mod}^{dg}(\mathcal{A}_U, \text{d})
\longrightarrow
\textit{Mod}^{dg}(\mathcal{A}, \text{d})
$$
，它是 $j_*$ 的右伴随。具体地，给定微分分次
$\mathcal{A}_U$-模 $\mathcal{M}$，定义 $j_!\mathcal{M}$ 为第
\ref{sdga-section-localize-graded} 节构造的分次 $\mathcal{A}$-模，
其微分为
$j_!\text{d} : j_!\mathcal{M}^n \to j_!\mathcal{M}^{n + 1}$.
给定微分分次 $\mathcal{A}_U$-模之间次数 $n$ 的齐次映射
$f:\mathcal{M}\to\mathcal{M}'$，得到微分分次 $\mathcal{A}$-模之间
次数 $n$ 的齐次映射
$j_!f:j_!\mathcal{M}\to j_!\mathcal{M}'$。略去该构造与微分相容的
直接验证。这样便得到所需函子。

\begin{lemma}
\label{sdga-lemma-extension-by-zero-dg}
在上述情形下，有
$$
\Hom_{\textit{Mod}^{dg}(\mathcal{A}, \text{d})}(
j_!\mathcal{M}, \mathcal{N}) =
\Hom_{\textit{Mod}^{dg}(\mathcal{A}_U, \text{d})}(
\mathcal{M}, j^*\mathcal{N})
$$
\end{lemma}

\begin{proof}
略。提示：引理 \ref{sdga-lemma-extension-by-zero-graded} 已经说明该结论在
分次层面成立。因此只须检验所得同构与微分相容。
\end{proof}

\begin{lemma}
\label{sdga-lemma-tensor-with-extension-by-zero-dg}
在上述情形下，设 $\mathcal{M}$ 是右微分分次 $\mathcal{A}_U$-模，
$\mathcal{N}$ 是左微分分次 $\mathcal{A}$-模。则
$$
j_!\mathcal{M} \otimes_\mathcal{A} \mathcal{N} =
j_!(\mathcal{M} \otimes_{\mathcal{A}_U} \mathcal{N}|_U)
$$
作为 $\mathcal{O}$-模复形的等式，并且关于 $\mathcal{M}$ 和
$\mathcal{N}$ 具有函子性。
\end{lemma}

\begin{proof}
作为分次模，该结论由引理
\ref{sdga-lemma-tensor-with-extension-by-zero} 得到。
略去此同构与微分相容的验证。
\end{proof}








\section{微分分次模层上的平移函子}
\label{sdga-section-shift-dg}

\noindent
设 $(\mathcal{C},\mathcal{O})$ 是环化位点，$\mathcal{A}$ 是
$(\mathcal{C},\mathcal{O})$ 上的微分分次代数层，$\mathcal{M}$
是微分分次 $\mathcal{A}$-模，$k\in\mathbf{Z}$。如下定义
$\mathcal{M}$ 的{\it 第 $k$ 次平移}，记作 $\mathcal{M}[k]$：
\begin{enumerate}
\item 作为分次 $\mathcal{A}$-模，$\mathcal{M}[k]$ 如第
\ref{sdga-section-shift} 节所定义；
\item 微分
$d_{\mathcal{M}[k]} : (\mathcal{M}[k])^n \to (\mathcal{M}[k])^{n + 1}$
定义为
$(-1)^k\text{d}_\mathcal{M} : \mathcal{M}^{n + k} \to \mathcal{M}^{n + k + 1}$.
\end{enumerate}
对次数 $n$ 齐次的 $\mathcal{A}$-模同态
$f:\mathcal{L}\to\mathcal{M}$，令
$f[k]:\mathcal{L}[k]\to\mathcal{M}[k]$ 由与 $f$ 相同的分量映射给出。
则 $f[k]$ 是次数 $n$ 齐次的 $\mathcal{A}$-模映射。这给出映射
$$
\Hom_{\textit{Mod}^{dg}(\mathcal{A}, \text{d})}(\mathcal{L}, \mathcal{M})
\longrightarrow
\Hom_{\textit{Mod}^{dg}(\mathcal{A}, \text{d})}
(\mathcal{L}[k], \mathcal{M}[k])
$$
，并与微分相容（这是因为 $\mathcal{L}$ 与 $\mathcal{M}$ 的微分符号
改变了相同的量）。这些选择与《微分分次代数》定义
\ref{dga-definition-shift} 中的选择相容。显然，这定义微分分次范畴之间的函子
$$
[k] :
\textit{Mod}^{dg}(\mathcal{A}, \text{d})
\longrightarrow
\textit{Mod}^{dg}(\mathcal{A}, \text{d})
$$
，并且 $[k+l]=[k]\circ[l]$。

\medskip\noindent
我们断言，在第 \ref{sdga-section-shift} 节对底层分次模定义的同构
$$
\Hom_{\textit{Mod}^{dg}(\mathcal{A}, \text{d})}(\mathcal{L}, \mathcal{M}[k])
=
\Hom_{\textit{Mod}^{dg}(\mathcal{A}, \text{d})}(\mathcal{L}, \mathcal{M})[k]
$$
与微分相容。为说明这一点，设
$f:\mathcal{L}\to\mathcal{M}[k]$ 是次数 $n$ 齐次的右
$\mathcal{A}$-模映射；它是左端的次数 $n$ 元素。用
$f':\mathcal{L}\to\mathcal{M}$ 表示次数 $n+k$ 齐次且与 $f$
具有{\bf 相同}分量映射的 $\mathcal{A}$-模映射。按我们的约定，
它就是右端对应的次数 $n$ 元素。按左端微分的定义，得到
$$
\text{d}_{LHS}(f) =
\text{d}_{\mathcal{M}[k]} \circ f - (-1)^n f \circ \text{d}_\mathcal{L} =
(-1)^k\text{d}_\mathcal{M} \circ f - (-1)^n f \circ \text{d}_\mathcal{L}
$$
；而对右端微分，得到
$$
\text{d}_{RHS}(f') =
(-1)^k\left(
\text{d}_\mathcal{M} \circ f' - (-1)^{n + k} f' \circ \text{d}_\mathcal{L}
\right) =
(-1)^k\text{d}_\mathcal{M} \circ f' - (-1)^n f' \circ \text{d}_\mathcal{L}
$$
。这些映射具有相同的分量映射，证明完成。










\section{同伦范畴}
\label{sdga-section-homotopy}

\noindent
本节对应《微分分次代数》第 \ref{dga-section-homotopy} 节。

\begin{definition}
\label{sdga-definition-homotopy}
设 $(\mathcal{C}, \mathcal{O})$ 是环化位点，
$\mathcal{A}$ 是 $(\mathcal{C}, \mathcal{O})$ 上的微分分次代数层。设
$f, g : \mathcal{M} \to \mathcal{N}$
是微分分次 $\mathcal{A}$-模同态。{\it $f$ 与 $g$ 之间的同伦}
是次数为 $-1$ 的齐次分次 $\mathcal{A}$-模映射
$h : \mathcal{M} \to \mathcal{N}$，且满足
$$
f - g = \text{d}_\mathcal{N} \circ h + h \circ \text{d}_\mathcal{M}
$$
若存在这样的同伦，则称 $f$ 与 $g$ {\it 同伦}。
\end{definition}

\noindent
在该定义的情形下，若有映射
$a : \mathcal{K} \to \mathcal{M}$ 和
$c : \mathcal{N} \to \mathcal{L}$，则
$$
\begin{matrix}
h\text{ 是一个同伦} \\
\text{，介于 }f\text{ 与 }g\text{ 之间}
\end{matrix}
\quad
\Rightarrow
\quad
\begin{matrix}
c \circ h \circ a\text{ 是}\\
c \circ f \circ a\text{ 与 }c\circ g \circ a\text{ 之间的同伦}
\end{matrix}
$$
因此可以定义 $\textit{Mod}(\mathcal{A}, \text{d})$
中态射同伦类的复合。

\begin{definition}
\label{sdga-definition-complexes-notation}
设 $(\mathcal{C}, \mathcal{O})$ 是环化位点，
$\mathcal{A}$ 是 $(\mathcal{C}, \mathcal{O})$ 上的微分分次代数层。
{\it 同伦范畴}，记作 $K(\textit{Mod}(\mathcal{A}, \text{d}))$，
是如下范畴：其对象为 $\textit{Mod}(\mathcal{A}, \text{d})$
的对象，其态射为微分分次 $\mathcal{A}$-模同态的同伦类。
\end{definition}

\noindent
记号 $K(\textit{Mod}(\mathcal{A}, \text{d}))$ 并非标准记号，
但至少与 Stacks 项目其他地方对 $K(-)$ 的用法一致。

\medskip\noindent
在《微分分次代数》定义
\ref{dga-definition-homotopy-category-of-dga-category} 中，
我们定义了微分分次范畴 $\mathcal{S}$ 的复形范畴
$\text{Comp}(\mathcal{S})$ 和同伦范畴 $K(\mathcal{S})$。
将此定义应用于微分分次范畴
$\textit{Mod}^{dg}(\mathcal{A}, \text{d})$，得到
$$
\textit{Mod}(\mathcal{A}, \text{d}) =
\text{Comp}(\textit{Mod}^{dg}(\mathcal{A}, \text{d}))
$$
（参见第 \ref{sdga-section-dgm-dg-cat} 节的讨论），并得到
$$
K(\textit{Mod}(\mathcal{A}, \text{d})) =
K(\textit{Mod}^{dg}(\mathcal{A}, \text{d}))
$$
要验证后一等式，注意当 $h$ 如定义 \ref{sdga-definition-homotopy} 中所述时，
有等式
$$
\text{d}_{\Hom_{\textit{Mod}^{dg}(\mathcal{A}, \text{d})}
(\mathcal{M}, \mathcal{N})}(h) =
\text{d}_\mathcal{N} \circ h + h \circ \text{d}_\mathcal{M}
$$
细节从略。

\begin{lemma}
\label{sdga-lemma-homotopy-direct-sums}
设 $(\mathcal{C}, \mathcal{O})$ 是环化位点，
$\mathcal{A}$ 是 $(\mathcal{C}, \mathcal{O})$ 上的微分分次代数层。
同伦范畴 $K(\textit{Mod}(\mathcal{A}, \text{d}))$ 有直和与直积。
\end{lemma}

\begin{proof}
从略。提示：只需使用引理 \ref{sdga-lemma-dgm-abelian} 中的直和与直积。
这是可行的，因为我们已经看到这些函子与到分次 $\mathcal{A}$-模范畴
的遗忘函子交换，而且在阿贝尔群族的范畴上，
$\prod$ 与 $\bigoplus$ 都是正合函子。
\end{proof}










\section{锥与三角}
\label{sdga-section-conclude-triangulated}

\noindent
本节利用《微分分次代数》第 \ref{dga-section-review} 节的内容，
证明微分分次 $\mathcal{A}$-模范畴的同伦范畴是三角范畴。

\begin{lemma}
\label{sdga-lemma-axioms-AB}
设 $(\mathcal{C}, \mathcal{O})$ 是环化位点，
$\mathcal{A}$ 是 $(\mathcal{C}, \mathcal{O})$ 上的微分分次代数层。
微分分次范畴 $\textit{Mod}^{dg}(\mathcal{A}, \text{d})$
满足《微分分次代数》第 \ref{dga-section-review} 节的公理 (A) 与 (B)。
\end{lemma}

\begin{proof}
设给定微分分次 $\mathcal{A}$-模 $\mathcal{M}$ 与 $\mathcal{N}$。
以显然的方式定义微分分次 $\mathcal{A}$-模
$\mathcal{M} \oplus \mathcal{N}$。此时余投射
$i : \mathcal{M} \to \mathcal{M} \oplus \mathcal{N}$、
$j : \mathcal{N} \to \mathcal{M} \oplus \mathcal{N}$ 以及投射
$p : \mathcal{M} \oplus \mathcal{N} \to \mathcal{N}$、
$q : \mathcal{M} \oplus \mathcal{N} \to \mathcal{M}$
都是微分分次 $\mathcal{A}$-模态射。因此 $i,j,p,q$
均为次数 $0$ 的齐次闭态射，即 $\text{d}(i)=0$，等等。
故此直和是《微分分次代数》定义
\ref{dga-definition-dg-direct-sum} 意义下的微分分次和。
这证明了公理 (A)。

\medskip\noindent
公理 (B) 已在第 \ref{sdga-section-shift-dg} 节中证明。
\end{proof}

\noindent
设 $(\mathcal{C}, \mathcal{O})$ 是环化位点，
$\mathcal{A}$ 是 $(\mathcal{C}, \mathcal{O})$ 上的微分分次代数层。
回忆，$\textit{Mod}(\mathcal{A}, \text{d})$ 中的序列
$$
0 \to \mathcal{K} \to \mathcal{L} \to \mathcal{N} \to 0
$$
若在 $\textit{Mod}(\mathcal{A})$ 中分裂，换言之，若作为分次
$\mathcal{A}$-模序列分裂，则称为{\it 容许短正合序列}
（《微分分次代数》第 \ref{dga-section-review} 节）。
以 $s : \mathcal{N} \to \mathcal{L}$ 和
$\pi : \mathcal{L} \to \mathcal{K}$ 表示分次
$\mathcal{A}$-模意义下的分裂映射。结合引理 \ref{sdga-lemma-axioms-AB}
与《微分分次代数》引理 \ref{dga-lemma-get-triangle}，得到三角
$$
\mathcal{K} \to \mathcal{L} \to \mathcal{N} \to \mathcal{K}[1]
$$
其中《微分分次代数》引理 \ref{dga-lemma-get-triangle} 的证明中
的箭头 $\mathcal{N} \to \mathcal{K}[1]$ 构造为
$$
\delta = \pi \circ
\text{d}_{\Hom_{\textit{Mod}^{dg}(\mathcal{A}, \text{d})}
(\mathcal{L}, \mathcal{M})}(s) =
\pi \circ \text{d}_\mathcal{L} \circ s -
\pi \circ s \circ \text{d}_\mathcal{N} =
\pi \circ \text{d}_\mathcal{L} \circ s
$$
请原谅这个可怕的记号。无论如何，在我们的情形下，
《微分分次代数》引理 \ref{dga-lemma-get-triangle} 所构造的边界映射
$\delta$，在底层 $\mathcal{O}$-模复形上与 Stacks 项目各处用于
逐项分裂的复形短正合序列的通常边界映射一致；参见《导出范畴》
定义 \ref{derived-definition-split-ses}。

\begin{definition}
\label{sdga-definition-cone}
设 $(\mathcal{C}, \mathcal{O})$ 是环化位点，
$\mathcal{A}$ 是 $(\mathcal{C}, \mathcal{O})$ 上的微分分次代数层。
设 $f : \mathcal{K} \to \mathcal{L}$ 是微分分次
$\mathcal{A}$-模同态。$f$ 的{\it 锥}是按如下方式定义的微分分次
$\mathcal{A}$-模 $C(f)$：
\begin{enumerate}
\item 其底层 $\mathcal{O}$-模复形是相应的 $\mathcal{A}$-模复形映射
$f : \mathcal{K}^\bullet \to \mathcal{L}^\bullet$ 的锥，即
$C(f)^n = \mathcal{L}^n \oplus \mathcal{K}^{n + 1}$，微分为
$$
d_{C(f)} =
\left(
\begin{matrix}
\text{d}_\mathcal{L} & f \\
0 & -\text{d}_\mathcal{K}
\end{matrix}
\right)
$$
\item 乘法映射
$$
C(f)^n \times \mathcal{A}^m \to C(f)^{n + m}
$$
是乘法映射
$\mathcal{L}^n \times \mathcal{A}^m \to \mathcal{L}^{n + m}$ 与
$\mathcal{K}^{n + 1} \times \mathcal{A}^m \to \mathcal{K}^{n + 1 + m}$.
的直和。
\end{enumerate}
显然的映射诱导出典范的微分分次 $\mathcal{A}$-模同态
$i : \mathcal{L} \to C(f)$ 与 $p : C(f) \to \mathcal{K}[1]$。
\end{definition}

\noindent
注意，在该定义的情形下，序列
$$
0 \to \mathcal{L} \to C(f) \to \mathcal{K}[1] \to 0
$$
是容许短正合序列。

\begin{lemma}
\label{sdga-lemma-axiom-C}
设 $(\mathcal{C}, \mathcal{O})$ 是环化位点，
$\mathcal{A}$ 是 $(\mathcal{C}, \mathcal{O})$ 上的微分分次代数层。
微分分次范畴 $\textit{Mod}^{dg}(\mathcal{A}, \text{d})$
满足《微分分次代数》情形 \ref{dga-situation-ABC} 中表述的公理 (C)。
\end{lemma}

\begin{proof}
设 $f : \mathcal{K} \to \mathcal{L}$ 是微分分次
$\mathcal{A}$-模同态。由上文有容许短正合序列
$$
0 \to \mathcal{L} \to C(f) \to \mathcal{K}[1] \to 0
$$
为完成证明，只需说明与之相伴的边界映射
$$
\delta : \mathcal{K}[1] \to \mathcal{L}[1]
$$
（参见上文的讨论）等于 $f[1]$。对截面
$s : \mathcal{K}[1] \to C(f)$，在次数 $n$ 上取嵌入
$\mathcal{K}^{n + 1} \to C(f)^n$。映射 $\pi$ 在次数 $n$
上由投射 $C(f)^n \to \mathcal{L}^n$ 给出。因此最后需要计算
$$
\delta = \pi \circ \text{d}_{C(f)} \circ s
$$
（参见上文的讨论）。用矩阵记号，此式等于
$$
\left(
\begin{matrix}
1 & 0
\end{matrix}
\right)
\left(
\begin{matrix}
\text{d}_\mathcal{L} & f \\
0 & -\text{d}_\mathcal{K}
\end{matrix}
\right)
\left(
\begin{matrix}
0 \\
1
\end{matrix}
\right) = f,
$$
正如所需。
\end{proof}

\noindent
至此可知，《微分分次代数》第 \ref{dga-section-review} 节中
证明的所有引理对微分分次范畴
$\textit{Mod}^{dg}(\mathcal{A}, \text{d})$ 均成立。
特别地，有如下命题。

\begin{proposition}
\label{sdga-proposition-homotopy-category-triangulated}
设 $(\mathcal{C}, \mathcal{O})$ 是环化位点，
$\mathcal{A}$ 是 $(\mathcal{C}, \mathcal{O})$ 上的微分分次代数层。
同伦范畴 $K(\textit{Mod}(\mathcal{A}, \text{d}))$ 是三角范畴，其中
\begin{enumerate}
\item 平移函子是第 \ref{sdga-section-shift-dg} 节中构造的函子；
\item 杰出三角是 $K(\textit{Mod}(\mathcal{A}, \text{d}))$
中作为三角同构于形如
$$
\mathcal{K} \to \mathcal{L} \to \mathcal{N}
\xrightarrow{\delta} \mathcal{K}[1],\quad\quad
\delta = \pi \circ \text{d}_\mathcal{L} \circ s
$$
的三角；该三角由上文 $\textit{Mod}(\mathcal{A}, \text{d})$
中的容许短正合序列
$0 \to \mathcal{K} \to \mathcal{L} \to \mathcal{N} \to 0$ 构造。
\end{enumerate}
\end{proposition}

\begin{proof}
回忆 $K(\textit{Mod}(\mathcal{A}, \text{d})) =
K(\textit{Mod}^{dg}(\mathcal{A}, \text{d}))$，参见第
\ref{sdga-section-homotopy} 节。于是命题由引理 \ref{sdga-lemma-axioms-AB}、
\ref{sdga-lemma-axiom-C} 以及《微分分次代数》命题
\ref{dga-proposition-ABC-homotopy-category-triangulated} 得出。
\end{proof}

\begin{remark}
\label{sdga-remark-cone-identity}
设 $(\mathcal{C}, \mathcal{O})$ 是环化位点，
$\mathcal{A}$ 是 $(\mathcal{C}, \mathcal{O})$ 上的微分分次代数层。
把恒等映射 $\mathcal{A}\to\mathcal{A}$ 视为微分分次
$\mathcal{A}$-模映射，并令 $C=C(\text{id}_\mathcal{A})$ 为其锥。则
$$
\Hom_{\textit{Mod}(\mathcal{A}, \text{d})}(C, \mathcal{M}) =
\{(x, y) \in
\Gamma(\mathcal{C}, \mathcal{M}^0) \times
\Gamma(\mathcal{C}, \mathcal{M}^{-1}) \mid
x = \text{d}(y)\}
$$
其中从左到右的映射把 $f$ 送到偶对 $(x,y)$：
$x$ 是 $C^{-1}=\mathcal{A}^{-1}\oplus\mathcal{A}^0$
的整体截面 $(0,1)$ 的像，而 $y$ 是
$C^0=\mathcal{A}^0\oplus\mathcal{A}^1$
的整体截面 $(1,0)$ 的像。
\end{remark}

\begin{lemma}
\label{sdga-lemma-dgm-grothendieck-abelian}
设 $(\mathcal{C}, \mathcal{O})$ 是环化位点，
$(\mathcal{A}, \text{d})$ 是微分分次 $\mathcal{O}$-代数。
范畴 $\textit{Mod}(\mathcal{A}, \text{d})$ 是格罗滕迪克阿贝尔范畴。
\end{lemma}

\begin{proof}
由引理 \ref{sdga-lemma-dgm-abelian} 以及格罗滕迪克阿贝尔范畴的定义
（《内射对象》定义 \ref{injectives-definition-grothendieck-conditions}），
只需证明 $\textit{Mod}(\mathcal{A}, \text{d})$ 有生成元。
对 $\mathcal{C}$ 的每个对象 $U$，以 $C_U$ 表示如评注
\ref{sdga-remark-cone-identity} 中恒等映射
$\mathcal{A}_U\to\mathcal{A}_U$ 的锥。我们断言
$$
\mathcal{G} = \bigoplus\nolimits_{k, U} j_{U!}C_U[k]
$$
是生成元，其中直和遍历 $\mathcal{C}$ 的所有对象 $U$
与所有 $k\in\mathbf{Z}$。事实上，给定微分分次
$\mathcal{A}$-模 $\mathcal{M}$，若从 $\mathcal{G}$
到 $\mathcal{M}$ 没有非零映射，则对所有 $k$ 与 $U$ 有
\begin{align*}
\Hom_{\textit{Mod}(\mathcal{A})}(j_{U!}C_U[k], \mathcal{M}) \\
& =
\Hom_{\textit{Mod}(\mathcal{A}_U)}(C_U[k], \mathcal{M}|_U) \\
& =
\{(x, y) \in \mathcal{M}^{-k}(U) \times \mathcal{M}^{-k - 1}(U) \mid
x = \text{d}(y)\}
\end{align*}
等于零。因此 $\mathcal{M}$ 为零。
\end{proof}











\section{平坦消解}
\label{sdga-section-P-resolutions}

\noindent
本节对应《微分分次代数》第 \ref{dga-section-P-resolutions} 节。

\medskip\noindent
设 $(\mathcal{C}, \mathcal{O})$ 是环化位点，
$\mathcal{A}$ 是 $(\mathcal{C}, \mathcal{O})$ 上的微分分次代数层。
若右微分分次 $\mathcal{A}$-模 $\mathcal{P}$ 满足
\begin{enumerate}
\item 函子
$\mathcal{N} \mapsto \mathcal{P} \otimes_\mathcal{A} \mathcal{N}$
在分次左 $\mathcal{A}$-模范畴上正合；
\item 若 $\mathcal{N}$ 是无环微分分次左 $\mathcal{A}$-模，则
$\mathcal{P} \otimes_\mathcal{A} \mathcal{N}$ 无环；
\item 对任意环化拓扑斯态射
$(f, f^\sharp) : (\Sh(\mathcal{C}'), \mathcal{O}')
\to (\Sh(\mathcal{C}), \mathcal{O})$、任意微分分次
$\mathcal{O}'$-代数 $\mathcal{A}'$，以及任意微分分次
$f^{-1}\mathcal{O}_\mathcal{D}$-代数映射
$\varphi : f^{-1}\mathcal{A} \to \mathcal{A}'$，
则把拉回 $f^*\mathcal{P}$（第 \ref{sdga-section-functoriality-dg} 节）
视为微分分次 $\mathcal{A}'$-模时，它满足性质 (1) 与 (2)。
\end{enumerate}
则称 $\mathcal{P}$ {\it 良好}。第一个条件表示 $\mathcal{P}$
作为右分次 $\mathcal{A}$-模是平坦的；第二个条件表示
$\mathcal{P}$ 在 Spaltenstein 的意义下是 K-平坦的
（参见《位点上的上同调》第
\ref{sites-cohomology-section-flat} 节）；第三个条件表示
任意基变换后这些性质仍成立。

\medskip\noindent
也许令人意外的是，良好模有很多。

\begin{lemma}
\label{sdga-lemma-supply-good}
设 $(\mathcal{C}, \mathcal{O})$ 是环化位点，
$\mathcal{A}$ 是 $(\mathcal{C}, \mathcal{O})$ 上的微分分次代数层，
$U \in \Ob(\mathcal{C})$。则 $j_!\mathcal{A}_U$
是良好微分分次 $\mathcal{A}$-模。
\end{lemma}

\begin{proof}
设 $\mathcal{N}$ 是分次左 $\mathcal{A}$-模。
由引理 \ref{sdga-lemma-tensor-with-extension-by-zero}，作为分次模有
$$
j_!\mathcal{A}_U \otimes_\mathcal{A} \mathcal{N} =
j_!(\mathcal{A}_U \otimes_{\mathcal{A}_U} \mathcal{N}|_U) =
j_!(\mathcal{N}_U)
$$
由于限制到 $U$ 与 $j_!$ 都正合，这证明了条件 (1)。
利用引理 \ref{sdga-lemma-tensor-with-extension-by-zero-dg}，
同一论证也证明条件 (2)。

\medskip\noindent
考虑环化拓扑斯态射
$(f, f^\sharp) : (\Sh(\mathcal{C}'), \mathcal{O}')
\to (\Sh(\mathcal{C}), \mathcal{O})$、微分分次
$\mathcal{O}'$-代数 $\mathcal{A}'$，以及微分分次
$f^{-1}\mathcal{O}$-代数映射
$\varphi : f^{-1}\mathcal{A} \to \mathcal{A}'$。
需要证明
$$
f^*j_!\mathcal{A}_U =
f^{-1}j_!\mathcal{A}_U \otimes_{f^{-1}\mathcal{A}} \mathcal{A}'
$$
对带有微分分次 $\mathcal{O}'$-代数层 $\mathcal{A}'$ 的环化拓扑斯
$(\Sh(\mathcal{C}'), \mathcal{O}')$ 满足 (1) 与 (2)。
为证明这一点，可以把 $(\Sh(\mathcal{C}), \mathcal{O})$
和 $(\Sh(\mathcal{C}'), \mathcal{O}')$ 替换为等价的环化拓扑斯。
因此由《位点上的模》引理
\ref{sites-modules-lemma-morphism-ringed-topoi-comes-from-morphism-ringed-sites}
可设 $f$ 来自连续函子 $u:\mathcal{C}\to\mathcal{C}'$
给出的位点态射 $f:\mathcal{C}\to\mathcal{C}'$。
此时置 $U'=u(U)$，并以
$j' : \Sh(\mathcal{C}'/U') \to \Sh(\mathcal{C}')$
表示相应的局部化态射。得到环化拓扑斯态射的交换方块
$$
\xymatrix{
(\Sh(\mathcal{C}'/U'), \mathcal{O}'_{U'})
\ar[rr]_{(j', (j')^\sharp)} \ar[d]_{(f', (f')^\sharp)} & &
(\Sh(\mathcal{C}'), \mathcal{O}')
\ar[d]^{(f, f^\sharp)} \\
(\Sh(\mathcal{C}/U), \mathcal{O}_U)
\ar[rr]^{(j, j^\sharp)} & &
(\Sh(\mathcal{C}), \mathcal{O}).
}
$$
且有 $f'_*(j')^{-1}=j^{-1}f_*$。参见《位点上的模》引理
\ref{sites-modules-lemma-localize-morphism-ringed-sites}。
由伴随的唯一性得到 $f^{-1}j_!=j'_!(f')^{-1}$。因此
\begin{align*}
f^*j_!\mathcal{A}_U
& =
f^{-1}j_!\mathcal{A}_U \otimes_{f^{-1}\mathcal{A}} \mathcal{A}' \\
& =
j'_!(f')^{-1}\mathcal{A}_U \otimes_{f^{-1}\mathcal{A}} \mathcal{A}' \\
& =
j'_!\left(
(f')^{-1}\mathcal{A}_U \otimes_{f^{-1}\mathcal{A}|_{U'}}
\mathcal{A}'|_{U'}\right) \\
& =
j'_!\mathcal{A}'_{U'}
\end{align*}
第一个等式是 $j_!\mathcal{A}_U$ 拉回为
$\mathcal{A}'$ 上微分分次模的定义。第二个等式来自
$f^{-1}j_!=j'_!(f')^{-1}$。第三个等式来自把引理
\ref{sdga-lemma-tensor-with-extension-by-zero-dg} 应用于环化位点
$(\mathcal{C}',f^{-1}\mathcal{O})$：其微分分次代数层为
$f^{-1}\mathcal{A}$，所用微分分次模分别是
$\mathcal{C}'/U'$ 上的 $(f')^{-1}\mathcal{A}_U$
与 $\mathcal{C}'$ 上的 $\mathcal{A}'$。第四个等式成立是因为
当然有 $(f')^{-1}\mathcal{A}_U=f^{-1}\mathcal{A}|_{U'}$。
由此可见，拉回仍是同类的模，而我们已在上文为它证明条件 (1) 与 (2)。
\end{proof}

\begin{lemma}
\label{sdga-lemma-good-admissible-ses}
设 $(\mathcal{C}, \mathcal{O})$ 是环化位点，
$\mathcal{A}$ 是 $(\mathcal{C}, \mathcal{O})$ 上的微分分次代数层。设
$0 \to \mathcal{P} \to \mathcal{P}' \to \mathcal{P}'' \to 0$
是微分分次 $\mathcal{A}$-模的容许短正合序列。
若这三个模中任意两个良好，则第三个也良好。
\end{lemma}

\begin{proof}
对条件 (1)，结论是立即的，因为该序列在分次层面是直和。
对条件 (2)，注意对任意微分分次左 $\mathcal{A}$-模，序列
$$
0 \to 
\mathcal{P} \otimes_\mathcal{A} \mathcal{N} \to
\mathcal{P}' \otimes_\mathcal{A} \mathcal{N} \to
\mathcal{P}'' \otimes_\mathcal{A} \mathcal{N} \to 0
$$
是微分分次 $\mathcal{O}$-模的容许短正合序列
（遗忘微分后，张量积正是在分次模范畴中取得的）。
因此，若三者中任意两个作为 $\mathcal{O}$-模复形正合，
则第三个也正合。最后，同一论证表明，给定环化拓扑斯态射
$(f, f^\sharp) : (\Sh(\mathcal{C}'), \mathcal{O}')
\to (\Sh(\mathcal{C}), \mathcal{O})$、微分分次
$\mathcal{O}'$-代数 $\mathcal{A}'$，以及微分分次
$f^{-1}\mathcal{O}$-代数映射
$\varphi:f^{-1}\mathcal{A}\to\mathcal{A}'$，序列
$$
0 \to f^*\mathcal{P} \to f^*\mathcal{P}' \to f^*\mathcal{P}'' \to 0
$$
是微分分次 $\mathcal{A}'$-模的容许短正合序列，
而上述论证在这里同样适用。
\end{proof}

\begin{lemma}
\label{sdga-lemma-good-direct-sum}
设 $(\mathcal{C}, \mathcal{O})$ 是环化位点，
$\mathcal{A}$ 是 $(\mathcal{C}, \mathcal{O})$ 上的微分分次代数层。
良好微分分次 $\mathcal{A}$-模的任意直和是良好的。
良好微分分次 $\mathcal{A}$-模的滤过余极限也是良好的。
\end{lemma}

\begin{proof}
从略。提示：直和与滤过余极限均同张量积及拉回交换。
\end{proof}

\begin{lemma}
\label{sdga-lemma-good-quotient}
设 $(\mathcal{C}, \mathcal{O})$ 是环化位点，
$\mathcal{A}$ 是 $(\mathcal{C}, \mathcal{O})$ 上的微分分次代数层，
$\mathcal{M}$ 是微分分次 $\mathcal{A}$-模。存在微分分次
$\mathcal{A}$-模同态 $\mathcal{P}\to\mathcal{M}$，满足
\begin{enumerate}
\item $\mathcal{P} \to \mathcal{M}$ 是满射；
\item $\Ker(\text{d}_\mathcal{P}) \to \Ker(\text{d}_\mathcal{M})$
是满射；
\item $\mathcal{P}$ 良好。
\end{enumerate}
\end{lemma}

\begin{proof}
考虑三元组 $(U,k,x)$，其中 $U$ 是 $\mathcal{C}$ 的对象，
$k\in\mathbf{Z}$，而 $x$ 是 $U$ 上满足
$\text{d}_\mathcal{M}(x)=0$ 的 $\mathcal{M}^k$ 截面。
于是得到唯一的微分分次 $\mathcal{A}_U$-模态射
$\varphi_x : \mathcal{A}_U[-k] \to \mathcal{M}|_U$
把 $1$ 映到 $x$。它伴随于态射
$\psi_x : j_{U!}\mathcal{A}_U[-k] \to \mathcal{M}$.
注意，$1\in\mathcal{A}_U(U)$ 对应于次数为 $k$ 的截面
$1\in j_{U!}\mathcal{A}_U[-k](U)$；其微分为零，且
$\psi_x$ 把它映到 $x$。因此考虑映射
$$
\bigoplus\nolimits_{(U, k, x)} j_{U!}\mathcal{A}_U[-k]
\longrightarrow
\mathcal{M}
$$
便得到条件 (2) 与 (3)。事实上，对象
$j_{U!}\mathcal{A}_U[-k]$ 是良好的（引理 \ref{sdga-lemma-supply-good}），
而良好对象的任意直和仍良好（引理 \ref{sdga-lemma-good-direct-sum}）。

\medskip\noindent
接着考虑三元组 $(U,k,x)$，其中 $U$ 是 $\mathcal{C}$ 的对象，
$k\in\mathbf{Z}$，而 $x$ 是 $\mathcal{M}^k$ 的截面
（不必被微分消去）。此时可以如评注 \ref{sdga-remark-cone-identity}
中那样考虑恒等映射 $\mathcal{A}_U\to\mathcal{A}_U$ 的锥 $C_U$。
元素 $x$ 决定映射
$\varphi_x:C_U[-k-1]\to\mathcal{A}_U$，参见评注
\ref{sdga-remark-cone-identity}。由于有容许短正合序列
$$
0 \to \mathcal{A}_U \to C_U \to \mathcal{A}_U[1] \to 0
$$
由引理 \ref{sdga-lemma-good-admissible-ses} 和已经用过的引理
\ref{sdga-lemma-supply-good}，可知 $j_{U!}C_U$ 是良好模。
同上，对伴随于 $\varphi_x$ 的映射
$\psi_x:j_{U!}C_U\to\mathcal{M}$ 取直和所得的映射
$$
\bigoplus\nolimits_{(U, k, x)} j_{U!}C_U \longrightarrow \mathcal{M}
$$
是满射。再与第一段构造的映射取直和，即得结论。
\end{proof}

\begin{remark}
\label{sdga-remark-sheaf-graded-sets}
设 $(\mathcal{C},\mathcal{O})$ 是环化位点。
$\mathcal{C}$ 上的{\it 分次集合层}是配有集合层映射
$\deg:\mathcal{S}\to\underline{\mathbf{Z}}$ 的集合层
$\mathcal{S}$。以 $\mathcal{O}[\mathcal{S}]$ 表示分次
$\mathcal{O}$-模，即由分次集合层 $\mathcal{S}$
生成的自由 $\mathcal{O}$-模。更确切地，
$\mathcal{O}[\mathcal{S}]$ 的第 $n$ 个分次部分是规则
$$
U \longmapsto
\bigoplus\nolimits_{s \in \mathcal{S}(U),\ \deg(s) = n} s \cdot \mathcal{O}(U)
$$
的层化。赋予零微分后，也可以把它看作微分分次
$\mathcal{O}$-模。设 $\mathcal{A}$ 是分次 $\mathcal{O}$-代数层。
类似地定义 $\mathcal{A}[\mathcal{S}]$ 为分次
$\mathcal{A}$-模，其第 $n$ 个分次部分是规则
$$
U \longmapsto
\bigoplus\nolimits_{s \in \mathcal{S}(U)} s \cdot \mathcal{A}^{n - \deg(s)}(U)
$$
的层化。若 $\mathcal{A}$ 是微分分次 $\mathcal{O}$-代数，
则对所有 $s\in\mathcal{S}(U)$ 置 $\text{d}(s)=0$ 并层化，
从而把它变成微分分次 $\mathcal{O}$-模。
\end{remark}

\begin{lemma}
\label{sdga-lemma-free-graded-module-good}
设 $(\mathcal{C},\mathcal{O})$ 是环化位点，
$\mathcal{A}$ 是微分分次 $\mathcal{A}$-代数，
$\mathcal{S}$ 是 $\mathcal{C}$ 上的分次集合层。
则由 $\mathcal{S}$ 生成的自由分次模
$\mathcal{A}[\mathcal{S}]$，赋予评注
\ref{sdga-remark-sheaf-graded-sets} 中的微分后，
是良好微分分次 $\mathcal{A}$-模。
\end{lemma}

\begin{proof}
设 $\mathcal{N}$ 是分次左 $\mathcal{A}$-模。则有
$$
\mathcal{A}[\mathcal{S}] \otimes_\mathcal{A} \mathcal{N} =
\mathcal{O}[\mathcal{S}] \otimes_\mathcal{O} \mathcal{N} =
\mathcal{N}[\mathcal{S}]
$$
其中 $\mathcal{N}[\mathcal{S}$ 是分次 $\mathcal{O}$-模，
其次数 $n$ 部分是与预层
$$
U \longmapsto
\bigoplus\nolimits_{s \in \mathcal{S}(U)} s \cdot \mathcal{N}^{n - \deg(s)}(U)
$$
相伴的层。显然，$\mathcal{N}\to\mathcal{N}[\mathcal{S}]$
是正合函子，因此 $\mathcal{A}[\mathcal{S}$
作为分次 $\mathcal{A}$-模是平坦的。接着，设
$\mathcal{N}$ 是微分分次左 $\mathcal{A}$-模。则作为分次
$\mathcal{O}$-模有
$$
H^*(\mathcal{A}[\mathcal{S}] \otimes_\mathcal{A} \mathcal{N}) =
H^*(\mathcal{O}[\mathcal{S}] \otimes_\mathcal{O} \mathcal{N})
$$
由其在 $\mathcal{O}$ 上的平坦性，这等于
$$
H^*(\mathcal{N})[\mathcal{S}]
$$
。因此，若 $\mathcal{N}$ 无环，则
$\mathcal{A}[\mathcal{S}]\otimes_\mathcal{A}\mathcal{N}$ 无环。

\medskip\noindent
最后，考虑环化拓扑斯态射
$(f, f^\sharp) : (\Sh(\mathcal{C}'), \mathcal{O}')
\to (\Sh(\mathcal{C}), \mathcal{O})$
、微分分次 $\mathcal{O}'$-代数 $\mathcal{A}'$，以及微分分次
$f^{-1}\mathcal{O}$-代数映射
$\varphi:f^{-1}\mathcal{A}\to\mathcal{A}'$。不难看出
$$
f^*\mathcal{A}[\mathcal{S}] = \mathcal{A}'[f^{-1}\mathcal{S}]
$$
这就完成了该模良好的证明。
\end{proof}

\begin{lemma}
\label{sdga-lemma-resolve}
设 $(\mathcal{C},\mathcal{O})$ 是环化位点，
$\mathcal{A}$ 是 $(\mathcal{C},\mathcal{O})$ 上的微分分次代数层，
$\mathcal{M}$ 是微分分次 $\mathcal{A}$-模。存在微分分次
$\mathcal{A}$-模同态 $\mathcal{P}\to\mathcal{M}$，满足
\begin{enumerate}
\item $\mathcal{P}\to\mathcal{M}$ 是拟同构；
\item $\mathcal{P}$ 良好。
\end{enumerate}
\end{lemma}

\begin{proof}[第一种证明]
令 $\mathcal{S}_0$ 为分次集合层
（评注 \ref{sdga-remark-sheaf-graded-sets}），其次数 $n$ 部分为
$\Ker(\text{d}_\mathcal{M}^n)$。考虑微分分次模同态
$$
\mathcal{P}_0 = \mathcal{A}[\mathcal{S}_0] \longrightarrow \mathcal{M}
$$
其中左端如评注 \ref{sdga-remark-sheaf-graded-sets} 所述，而该映射把
$\mathcal{S}_0$ 的局部截面 $s$ 映到
$\mathcal{M}^{\deg(s)}$ 的相应局部截面
（它属于微分的核，故这确实是微分分次模映射）。
按构造，诱导的上同调层映射
$H^n(\mathcal{P}_0)\to H^n(\mathcal{M})$ 是满射。
我们将归纳构造映射
$$
\mathcal{P}_0 \to \mathcal{P}_1 \to \mathcal{P}_2 \to \ldots \to \mathcal{M}
$$
显然，对所有 $i$，$H^*(\mathcal{P}_i)\to H^*(\mathcal{M})$
都将是满射。给定 $\mathcal{P}_i\to\mathcal{M}$，
以 $\mathcal{S}_{i+1}$ 表示如下分次集合层：其次数 $n$ 部分为
$$
\Ker(\text{d}_{\mathcal{P}_i}^{n + 1})
\times_{\mathcal{M}^{n + 1}, \text{d}}
\mathcal{M}^n
$$
置
$$
\mathcal{P}_{i + 1} = \mathcal{P}_i \oplus \mathcal{A}[\mathcal{S}_{i + 1}]
$$
作为分次 $\mathcal{A}$-模，其微分及到 $\mathcal{M}$ 的映射定义如下：
\begin{enumerate}
\item 对 $\mathcal{P}_i$ 的局部截面，使用 $\mathcal{P}_i$
上的微分以及给定的到 $\mathcal{M}$ 的映射；
\item 对 $\mathcal{S}_{i+1}$ 的局部截面 $s=(p,m)$，
置 $\text{d}(s)$ 等于 $p$（视为 $\mathcal{P}_i$
的次数 $\deg(s)+1$ 截面），并把 $s$ 映到 $\mathcal{M}$ 中的 $m$；
\item 唯一延拓微分，使其满足莱布尼茨法则。
\end{enumerate}
这是有意义的，因为 $\text{d}(m)$ 是 $p$ 的像，且
$\text{d}(p)=0$。最后置
$\mathcal{P}=\colim\mathcal{P}_i$，并赋予诱导的到
$\mathcal{M}$ 的映射。

\medskip\noindent
映射 $\mathcal{P}\to\mathcal{M}$ 是拟同构：
$H^n(\mathcal{P})=\colim H^n(\mathcal{P}_i)$，而对每个 $i$，
映射 $H^n(\mathcal{P}_i)\to H^n(\mathcal{M})$ 是满射；
按构造，其核被映射
$H^n(\mathcal{P}_i)\to H^n(\mathcal{P}_{i+1})$ 消去。
每个 $\mathcal{P}_i$ 都良好，因为由引理
\ref{sdga-lemma-free-graded-module-good}，$\mathcal{P}_0$ 良好，
而每个 $\mathcal{P}_{i+1}$ 都是容许短正合序列
$0 \to \mathcal{P}_i \to \mathcal{P}_{i + 1} \to
\mathcal{A}[\mathcal{S}_{i + 1}] \to 0$
的中间项；其两端由归纳假设均良好。因此由引理
\ref{sdga-lemma-good-admissible-ses}，$\mathcal{P}_{i+1}$ 良好。
最后，由引理 \ref{sdga-lemma-good-direct-sum} 得到 $\mathcal{P}$ 良好。
\end{proof}

\begin{proof}[第二种证明]
我们强烈建议读者在阅读本证明前，先阅读《微分分次代数》引理
\ref{dga-lemma-resolve} 的证明。置 $\mathcal{M}=\mathcal{M}_0$。
归纳选取短正合序列
$$
0 \to \mathcal{M}_{i + 1} \to \mathcal{P}_i \to \mathcal{M}_i \to 0
$$
其中映射 $\mathcal{P}_i\to\mathcal{M}_i$ 按引理
\ref{sdga-lemma-good-quotient} 选取。这给出一个``消解''
$$
\ldots \to \mathcal{P}_2 \xrightarrow{f_2}
\mathcal{P}_1 \xrightarrow{f_1} \mathcal{P}_0 \to \mathcal{M} \to 0
$$
然后令 $\mathcal{P}$ 为如下定义的微分分次 $\mathcal{A}$-模：
\begin{enumerate}
\item 作为分次 $\mathcal{A}$-模，置
$\mathcal{P}=\bigoplus_{a\leq 0}\mathcal{P}_{-a}[-a]$，即其次数
$n$ 部分为
$\mathcal{P}^n=\bigoplus\nolimits_{a+b=n}\mathcal{P}_{-a}^b$；
\item $\mathcal{P}$ 上的微分按与双复形相伴的全复形的构造给出：
$$
\text{d}_\mathcal{P}(x) = f_{-a}(x) + (-1)^a \text{d}_{\mathcal{P}_{-a}}(x)
$$
其中 $x$ 是 $\mathcal{P}_{-a}^b$ 的局部截面。
\end{enumerate}
在这些约定下，$\mathcal{P}$ 确为微分分次 $\mathcal{A}$-模；
细节从略。有微分分次 $\mathcal{A}$-模映射
$\mathcal{P}\to\mathcal{M}$：当 $a<0$ 时，它在直和项
$\mathcal{P}_{-a}[-a]$ 上为零；当 $a=0$ 时，它是给定映射
$\mathcal{P}_0\to\mathcal{M}$。注意
$$
\mathcal{P} = \colim_i F_i\mathcal{P}
$$
其中 $F_i\mathcal{P}\subset\mathcal{P}$ 是微分分次
$\mathcal{A}$-子模，其底层分次 $\mathcal{A}$-模为
$$
F_i\mathcal{P} =
\bigoplus\nolimits_{i \geq -a \geq 0} \mathcal{P}_{-a}[-a]
$$
立即可见，映射
$$
0 \to F_1\mathcal{P} \to F_2\mathcal{P} \to F_3\mathcal{P} \to
\ldots \to \mathcal{P}
$$
都是容许单态射，且有容许短正合序列
$$
0 \to F_i\mathcal{P} \to
F_{i + 1}\mathcal{P} \to
\mathcal{P}_{i + 1}[i + 1] \to 0
$$
由归纳法及引理 \ref{sdga-lemma-good-admissible-ses}，
$F_i\mathcal{P}$ 是良好微分分次 $\mathcal{A}$-模。
由于 $\mathcal{P}=\colim F_i\mathcal{P}$，由引理
\ref{sdga-lemma-good-direct-sum} 可知 $\mathcal{P}$ 良好。

\medskip\noindent
最后还需证明 $\mathcal{P}\to\mathcal{M}$ 是拟同构。
若 $\mathcal{C}$ 有足够多的点，则逐茎检验后，该结论由
《同调》引理 \ref{homology-lemma-good-resolution-gives-qis}
这一初等结果得出。一般情形可如下论证
（这个证明实在太长了；也可以像初等证明那样使用局部截面，
但那种论证也相当长）。由于滤过余极限在阿贝尔群层范畴上正合，有
$$
H^d(\mathcal{P}) = \colim H^d(F_i\mathcal{P})
$$
我们断言，对每个 $i\geq0$ 和 $d\in\mathbf{Z}$：
(a) 有短正合序列
$$
0 \to H^d(\mathcal{M}_{i + 1}[i]) \to
H^d(F_i\mathcal{P}) \to H^d(\mathcal{M}) \to 0
$$
其中第二个箭头来自
$F_i\mathcal{P}\to\mathcal{P}\to\mathcal{M}$；且
(b) 复合
$$
H^d(\mathcal{M}_{i + 1}[i]) \to
H^d(F_i\mathcal{P}) \to
H^d(F_{i + 1}\mathcal{P})
$$
为零。显然，该断言足以完成证明。

\medskip\noindent
证明该断言。对任意 $i\geq0$，有映射
$\mathcal{M}_{i+1}[i]\to F_i\mathcal{P}$，它来自
$\mathcal{M}_{i+1}$ 作为 $f_i$ 的核到 $\mathcal{P}_i$ 的包含。
考虑定义 $C_i$ 的 $\mathcal{O}$-模复形短正合序列
$$
0 \to \mathcal{M}_{i + 1}[i] \to F_i\mathcal{P} \to C_i \to 0
$$
注意 $C_0=\mathcal{M}_0=\mathcal{M}$。还要注意，
$C_i$ 是与双复形 $C_i^{\bullet,\bullet}$ 相伴的全复形；
该双复形的各列为
$$
\mathcal{M}_i = \mathcal{P}_i/\mathcal{M}_{i + 1},
\mathcal{P}_{i - 1}, \ldots, \mathcal{P}_0
$$
并分别处于次数 $-i,-i+1,\ldots,0$。有双复形映射
$C_i^{\bullet,\bullet}\to C_{i-1}^{\bullet,\bullet}$：
在次数 $-i$ 的列上为 $0$，在次数 $-i+1$ 上为满射
$\mathcal{P}_{i-1}\to\mathcal{M}_{i-1}$，在其余各列上为恒等映射。
因此有复形映射
$$
C_i \longrightarrow C_{i - 1}
$$
这些映射是满拟同构，因为其核是如下双复形的全复形：
次数 $-i,-i+1$ 的两列都是 $\mathcal{M}_i$，
而这两列之间是恒等映射。利用由此得到的等同
$H^d(C_i) = H^d(C_{i - 1} = \ldots = H^d(\mathcal{M})$
，已经可知由上述复形短正合序列得到长正合序列
$$
H^d(\mathcal{M}_{i + 1}[i]) \to
H^d(F_i\mathcal{P}) \to H^d(\mathcal{M}) \to
H^{d + 1}(\mathcal{M}_{i + 1}[i])
$$
。另一方面，还有交换图
$$
\xymatrix{
\mathcal{M}_{i + 2}[i + 1] \ar[r]_a & T_{i + 1} \ar[r] &
F_{i + 1}\mathcal{P} \ar[r] & C_{i + 1} \ar[d] \\
& \mathcal{M}_{i + 1}[i] \ar[r] \ar[u]^b &
F_i\mathcal{P} \ar[u] \ar[r] &
C_i
}
$$
其中 $T_{i+1}$ 是如下双复形的全复形：各列
$\mathcal{P}_{i+1},\mathcal{M}_{i+1}$ 分别置于次数
$-i-1,-i$。换言之，$T_{i+1}$ 是映射
$\mathcal{P}_{i+1}\to\mathcal{M}_{i+1}$ 的锥的一个平移。
于是 $a$ 是拟同构，而映射 $a^{-1}\circ b$ 是
$D(\mathcal{O})$ 中与短正合序列
$$
0 \to \mathcal{M}_{i + 2} \to \mathcal{P}_{i + 1} \to \mathcal{M}_{i + 1} \to 0
$$
相伴的杰出三角的第三个映射的一个平移。映射
$H^d(\mathcal{P}_{i+1})\to H^d(\mathcal{M}_{i+1})$
是满射，因为选取映射时要求
$\Ker(\text{d}_{\mathcal{P}_{i + 1}}) \to
\Ker(\text{d}_{\mathcal{M}_{i+1}})$ 为满射。因此
$a^{-1}\circ b$ 在上同调层上为零。这证明了断言的 (b)。
由于 $T_{i+1}$ 是复形满射
$F_{i+1}\mathcal{P}\to C_i$ 的核，得到上同调长正合序列的映射
$$
\xymatrix{
H^d(T_{i + 1}) \ar[r] &
H^d(F_{i + 1}\mathcal{P}) \ar[r] &
H^d(\mathcal{M}) \ar[r] &
H^{d + 1}(T_{i + 1}) \\
H^d(\mathcal{M}_{i + 1}[i]) \ar[r] \ar[u] &
H^d(F_i\mathcal{P}) \ar[r] \ar[u] &
H^d(\mathcal{M}) \ar[r] \ar[u] &
H^{d + 1}(\mathcal{M}_{i + 1}[i]) \ar[u]
}
$$
由上述讨论，外侧两个竖直映射为零。因此映射
$H^d(F_{i+1}\mathcal{P})\to H^d(\mathcal{M})$ 是满射，
断言的 (a) 随之成立。更确切地说，由此得到 $i>0$
时的断言；$i=0$ 的情形留给读者
（事实上，为得到本引理，只需对所有 $i\gg0$ 证明该断言）。
\end{proof}

\begin{lemma}
\label{sdga-lemma-acyclic-good}
设 $(\mathcal{C},\mathcal{O})$ 是环化位点，
$\mathcal{A}$ 是 $(\mathcal{C},\mathcal{O})$ 上的微分分次代数层。
设 $\mathcal{P}$ 是良好的无环右微分分次 $\mathcal{A}$-模。
\begin{enumerate}
\item 对任意微分分次左 $\mathcal{A}$-模 $\mathcal{N}$，
张量积 $\mathcal{P}\otimes_\mathcal{A}\mathcal{N}$ 无环；
\item 对任意环化拓扑斯态射
$(f, f^\sharp) : (\Sh(\mathcal{C}'), \mathcal{O}')
\to (\Sh(\mathcal{C}), \mathcal{O})$、任意微分分次
$\mathcal{O}'$-代数 $\mathcal{A}'$，以及任意微分分次
$f^{-1}\mathcal{O}$-代数映射
$\varphi:f^{-1}\mathcal{A}\to\mathcal{A}'$，
拉回 $f^*\mathcal{P}$ 无环且良好。
\end{enumerate}
\end{lemma}

\begin{proof}
证明 (1)。由引理 \ref{sdga-lemma-resolve}，可以选取良好左微分分次
$\mathcal{Q}$ 以及拟同构 $\mathcal{Q}\to\mathcal{N}$。于是
$\mathcal{P} \otimes_\mathcal{A} \mathcal{Q}$
无环，因为 $\mathcal{Q}$ 良好。令 $\mathcal{N}'$
为映射 $\mathcal{Q}\to\mathcal{N}$ 的锥。则
$\mathcal{P}\otimes_\mathcal{A}\mathcal{N}'$ 无环，因为
$\mathcal{P}$ 良好，而 $\mathcal{N}'$ 作为拟同构的锥是无环的。
按我们的三角结构构造，在
$K(\textit{Mod}(\mathcal{A},\text{d}))$ 中有杰出三角
$$
\mathcal{Q} \to \mathcal{N} \to \mathcal{N}' \to \mathcal{Q}[1]
$$
。由于 $\mathcal{P}\otimes_\mathcal{A}-$ 把杰出三角映到杰出三角，
在 $K(\textit{Mod}(\mathcal{O}))$ 中得到杰出三角
$$
\mathcal{P} \otimes_\mathcal{A} \mathcal{Q} \to
\mathcal{P} \otimes_\mathcal{A} \mathcal{N} \to
\mathcal{P} \otimes_\mathcal{A} \mathcal{N}' \to
\mathcal{P} \otimes_\mathcal{A} \mathcal{Q}[1]
$$
从而得出结论。

\medskip\noindent
证明 (2)。由良好模的定义，$f^*\mathcal{P}$ 良好。回忆
$f^*\mathcal{P} = f^{-1}\mathcal{P} \otimes_{f^{-1}\mathcal{A}} \mathcal{A}'$.
此时 $f^{-1}\mathcal{P}$ 是良好的无环微分分次
$f^{-1}\mathcal{A}$-模（因为 $f^{-1}$ 正合）。
因此由 (1) 可知 $f^*\mathcal{P}$ 无环。
\end{proof}






\section{分次模的微分分次包}
\label{sdga-section-dg-hull}

\noindent
分次模 $\mathcal{N}$ 的微分分次包，是应用下列引理中的函子 $G$ 所得的对象。

\begin{lemma}
\label{sdga-lemma-dg-hull}
设 $(\mathcal{C},\mathcal{O})$ 是环化位点，
$\mathcal{A}$ 是 $(\mathcal{C},\mathcal{O})$ 上的微分分次代数层。
遗忘函子
$F : \textit{Mod}(\mathcal{A}, \text{d}) \to \textit{Mod}(\mathcal{A})$
有左伴随
$G : \textit{Mod}(\mathcal{A})\to\textit{Mod}(\mathcal{A},\text{d})$。
\end{lemma}

\begin{proof}
为证明 $G$ 存在，可以使用伴随函子定理；参见《范畴》定理
\ref{categories-theorem-adjoint-functor}
（注意我们交换了 $F$ 与 $G$ 的角色）。由引理
\ref{sdga-lemma-dgm-abelian}，$F$ 满足正合性条件。
集合论条件可如下验证：设给定分次 $\mathcal{A}$-模
$\mathcal{N}$。对任意映射
$$
\varphi : \mathcal{N} \longrightarrow F(\mathcal{M})
$$
可以考虑满足
$\Im(\varphi)\subset F(\mathcal{M}')$ 的最小微分分次
$\mathcal{A}$-子模 $\mathcal{M}'\subset\mathcal{M}$。
显然，$\mathcal{M}'$ 是分次 $\mathcal{A}$-模映射
$$
\mathcal{N} \oplus
\mathcal{N}[-1] \otimes_\mathcal{O} \mathcal{A}
\longrightarrow
\mathcal{M}
$$
的像，该映射定义为
$$
(n, \sum n_i \otimes a_i) \longmapsto
\varphi(n) + \sum \text{d}(\varphi(n_i)) a_i
$$
，因为容易看出此映射的像是 $\mathcal{M}$ 的微分分次子模。
因此，这些 $\mathcal{M}'$ 的可能同构类数目有界，结论得证。
\end{proof}

\noindent
设 $(\mathcal{C},\mathcal{O})$ 是环化位点，
$\mathcal{A}$ 是 $(\mathcal{C},\mathcal{O})$ 上的微分分次代数层，
$\mathcal{M}$ 是微分分次 $\mathcal{A}$-模，并设在
$\textit{Mod}(\mathcal{A})$ 中有短正合序列
$$
0 \to \mathcal{N} \to F(\mathcal{M}) \to \mathcal{N}' \to 0
$$
。于是得到典范的分次 $\mathcal{A}$-模同态
$$
\overline{\text{d}} : \mathcal{N} \to \mathcal{N}'[1]
$$
，定义如下：给定 $\mathcal{N}$ 的局部截面 $x$，
把 $x$ 视为 $\mathcal{M}$ 的局部截面，并以
$\overline{\text{d}}(x)$ 表示
$\text{d}_\mathcal{M}(x)$ 在 $\mathcal{N}'$ 中的像。

\begin{lemma}
\label{sdga-lemma-dg-hull-acyclic}
引理 \ref{sdga-lemma-dg-hull} 中的函子 $F,G$ 具有如下性质。
给定分次 $\mathcal{A}$-模 $\mathcal{N}$，有
\begin{enumerate}
\item 余单位 $\mathcal{N}\to F(G(\mathcal{N}))$ 是单射；
\item 映射
$\overline{\text{d}}:\mathcal{N}\to
\Coker(\mathcal{N}\to F(G(\mathcal{N})))[1]$ 是同构；
\item $G(\mathcal{N})$ 是无环微分分次 $\mathcal{A}$-模。
\end{enumerate}
\end{lemma}

\begin{proof}
注意性质 (3) 是性质 (1) 与 (2) 的推论。事实上，若 $s$ 是
$F(G(\mathcal{N}))$ 的非零局部截面且 $\text{d}(s)=0$，
则 $s$ 不可能属于
$\mathcal{N}\to F(G(\mathcal{N}))$ 的像。因此，$s$
在余核中的像 $\overline{s}$ 可写为
$\overline{\text{d}}(s')$，其中 $s'$ 是 $\mathcal{N}$
的某个局部截面。于是 $s=\text{d}(s')$，因为差
$s-\text{d}(s')$ 仍属于 $\text{d}$ 的核，并且包含在余单位的像中。

\medskip\noindent
暂时以 $\mathcal{A}_{gr}$、$\mathcal{A}_{dg}$ 分别表示把层
$\mathcal{A}$ 视为其自身上的（右）分次模与（右）微分分次模。
本引理最重要的情形是理解 $G(\mathcal{A}_{gr})$。
当然，$G(\mathcal{A}_{gr})$ 是
$\textit{Mod}(\mathcal{A},\text{d})$ 中表示函子
$$
\mathcal{M} \longmapsto
\Hom_{\textit{Mod}(\mathcal{A})}(\mathcal{A}_{gr}, F(\mathcal{M})) =
\Gamma(\mathcal{C}, \mathcal{M})
$$
的对象。由评注 \ref{sdga-remark-cone-identity} 可知，此函子由
$C[-1]$ 表示，其中 $C$ 是 $\mathcal{A}_{dg}$ 恒等映射的锥。
在 $\textit{Mod}(\mathcal{A},\text{d})$ 中有短正合序列
$$
0 \to \mathcal{A}_{dg}[-1] \to C[-1] \to \mathcal{A}_{dg} \to 0
$$
，它在 $\textit{Mod}(\mathcal{A})$ 中由余单位
$\mathcal{A}_{gr}\to F(C[-1])$ 分裂。因此
$G(\mathcal{A}_{gr})$ 满足性质 (1) 与 (2)。

\medskip\noindent
设 $U$ 是 $\mathcal{C}$ 的对象。以
$j_U:\mathcal{C}/U\to\mathcal{C}$ 表示局部化态射，
以 $\mathcal{A}_U$ 表示 $\mathcal{A}$ 到 $U$ 的限制。
用记号 $\mathcal{A}_{U,gr}$ 表示把 $\mathcal{A}_U$
视为分次 $\mathcal{A}_U$-模。以
$F_U:\textit{Mod}(\mathcal{A}_U,\text{d})\to
\textit{Mod}(\mathcal{A}_U)$ 表示遗忘函子，以 $G_U$
表示其伴随。此时有交换图
$$
\vcenter{
\xymatrix{
\textit{Mod}(\mathcal{A}, \text{d}) \ar[d]_{j_U^*} \ar[r]_F &
\textit{Mod}(\mathcal{A}) \ar[d]^{j_U^*} \\
\textit{Mod}(\mathcal{A}_U, \text{d}) \ar[r]^{F_U} &
\textit{Mod}(\mathcal{A}_U)
}
}
\quad\text{且}\quad
\vcenter{
\xymatrix{
\textit{Mod}(\mathcal{A}_U, \text{d}) \ar[r]_{F_U} \ar[d]_{j_{U!}} &
\textit{Mod}(\mathcal{A}_U) \ar[d]^{j_{U!}} \\
\textit{Mod}(\mathcal{A}, \text{d}) \ar[r]^F &
\textit{Mod}(\mathcal{A})
}
}
$$
，这来自第 \ref{sdga-section-functoriality-graded}、
\ref{sdga-section-functoriality-dg}、\ref{sdga-section-localize-graded}
及 \ref{sdga-section-localize-dg} 节中对 $j_U^*$ 与 $j_{U!}$ 的构造。
由伴随的唯一性得到 $j_{U!}\circ G_U=G\circ j_{U!}$。
由于 $j_{U!}$ 是正合函子，证明前一部分已经得到余单位
$\mathcal{A}_{U,gr}\to F_U(G_U(\mathcal{A}_{U,gr}))$
的性质 (1) 与 (2)，故这些性质也适用于余单位
$j_{U!}\mathcal{A}_{U, gr} \to F(G(j_{U!}\mathcal{A}_{U, gr})) =
j_{U!}F_U(G_U(\mathcal{A}_{U, gr}))$.

\medskip\noindent
在引理 \ref{sdga-lemma-gm-grothendieck-abelian} 的证明中，
我们已经看到 $\textit{Mod}(\mathcal{A})$ 的任意对象都是若干
$j_{U!}\mathcal{A}_{U,gr}$ 副本之直和的商。
由于 $G$ 是左伴随，它与直和交换。因此，只要性质 (1) 与 (2)
对一族对象成立，就对它们的直和成立。由此，
$\textit{Mod}(\mathcal{A})$ 的每个对象 $\mathcal{N}$
都可置于正合序列
$$
\mathcal{N}_1 \to \mathcal{N}_0 \to \mathcal{N} \to 0
$$
中，使得 (1) 与 (2) 对 $\mathcal{N}_1$ 和 $\mathcal{N}_0$
成立。利用 $G$ 的右正合性，从中推出 $\mathcal{N}$
的 (1) 与 (2)，留给读者。
\end{proof}









\section{K-内射微分分次模}
\label{sdga-section-K-injective}

\noindent
本节是在微分分次代数层上的微分分次模层这一背景下，
与《内射对象》第 \ref{injectives-section-K-injective} 节对应的内容。

\begin{lemma}
\label{sdga-lemma-characterize-injectives}
设 $(\mathcal{C},\mathcal{O})$ 是环化位点，$\mathcal{A}$
是 $(\mathcal{C},\mathcal{O})$ 上的分次代数层。
存在集合 $T$，并且对每个 $t\in T$ 存在分次
$\mathcal{A}$-模的单射 $\mathcal{N}_t\to\mathcal{N}'_t$，
使得 $\textit{Mod}(\mathcal{A})$ 的对象 $\mathcal{I}$
是内射对象，当且仅当对每个实线图
$$
\xymatrix{
\mathcal{N}_t \ar[r] \ar[d] & \mathcal{I} \\
\mathcal{N}'_t \ar@{..>}[ru]
}
$$
在 $\textit{Mod}(\mathcal{A})$ 中存在使图交换的虚线箭头。
\end{lemma}

\begin{proof}
这在任意格罗滕迪克阿贝尔范畴中成立，参见《内射对象》引理
\ref{injectives-lemma-characterize-injective}。
由引理 \ref{sdga-lemma-gm-grothendieck-abelian}，
范畴 $\textit{Mod}(\mathcal{A})$ 是格罗滕迪克阿贝尔范畴。
\end{proof}

\begin{definition}
\label{sdga-definition-graded-injective}
设 $(\mathcal{C},\mathcal{O})$ 是环化位点，
$(\mathcal{A},\text{d})$ 是 $(\mathcal{C},\mathcal{O})$
上的微分分次代数层。若把 $\mathcal{M}$ 视为分次
$\mathcal{A}$-模时，它是分次 $\mathcal{A}$-模范畴
$\textit{Mod}(\mathcal{A})$ 的内射对象，则称微分分次
$\mathcal{A}$-模 $\mathcal{I}$ {\it 分次内射}
\footnote{这可能不是标准术语。}。
\end{definition}

\begin{remark}
\label{sdga-remark-why-graded-injective}
设 $(\mathcal{C},\mathcal{O})$ 是环化位点，
$(\mathcal{A},\text{d})$ 是 $(\mathcal{C},\mathcal{O})$
上的微分分次代数层。设 $\mathcal{I}$ 是分次内射微分分次
$\mathcal{A}$-模，并设
$$
0 \to \mathcal{M}_1 \to \mathcal{M}_2 \to \mathcal{M}_3 \to 0
$$
是微分分次 $\mathcal{A}$-模的短正合序列。
由于 $\mathcal{I}$ 分次内射，得到
$\Gamma(\mathcal{C},\mathcal{O})$-模复形的短正合序列
$$
0 \to
\Hom_{\textit{Mod}^{dg}(\mathcal{A}, \text{d})}(\mathcal{M}_3, \mathcal{I})
\to
\Hom_{\textit{Mod}^{dg}(\mathcal{A}, \text{d})}(\mathcal{M}_2, \mathcal{I})
\to
\Hom_{\textit{Mod}^{dg}(\mathcal{A}, \text{d})}(\mathcal{M}_1, \mathcal{I})
\to 0
$$
。取上同调得到同伦范畴中同态群的长正合序列
$$
\xymatrix{
\Hom_{K(\textit{Mod}(\mathcal{A}, \text{d}))}(\mathcal{M}_3, \mathcal{I})
\ar[d] &
\Hom_{K(\textit{Mod}(\mathcal{A}, \text{d}))}(\mathcal{M}_3, \mathcal{I})[1]
\ar[d] \\
\Hom_{K(\textit{Mod}(\mathcal{A}, \text{d}))}(\mathcal{M}_2, \mathcal{I})
\ar[d] &
\Hom_{K(\textit{Mod}(\mathcal{A}, \text{d}))}(\mathcal{M}_2, \mathcal{I})[1]
\ar[d] \\
\Hom_{K(\textit{Mod}(\mathcal{A}, \text{d}))}(\mathcal{M}_1, \mathcal{I})
\ar[ruu]
&
\Hom_{K(\textit{Mod}(\mathcal{A}, \text{d}))}(\mathcal{M}_1, \mathcal{I})[1]
}
$$
。关键在于，即使没有假设原短正合序列容许，我们仍得到这一序列
（所以一般而言，该短正合序列并不在同伦范畴中定义杰出三角）。
\end{remark}

\begin{lemma}
\label{sdga-lemma-product-graded-injective}
设 $(\mathcal{C},\mathcal{O})$ 是环化位点，
$(\mathcal{A},\text{d})$ 是 $(\mathcal{C},\mathcal{O})$
上的微分分次代数层。设 $T$ 是集合，并且对每个 $t\in T$，
$\mathcal{I}_t$ 是分次内射微分分次 $\mathcal{A}$-模。
则 $\prod\mathcal{I}_t$ 是分次内射微分分次 $\mathcal{A}$-模。
\end{lemma}

\begin{proof}
这是因为内射对象的直积仍内射，参见《同调》引理
\ref{homology-lemma-product-injectives}；还因为通过遗忘函子，
$\textit{Mod}(\mathcal{A},\text{d})$ 中的直积与
$\textit{Mod}(\mathcal{A})$ 中的直积相容。
\end{proof}

\begin{lemma}
\label{sdga-lemma-characterize-graded-injectives-in-dg}
设 $(\mathcal{C},\mathcal{O})$ 是环化位点，
$(\mathcal{A},\text{d})$ 是 $(\mathcal{C},\mathcal{O})$
上的微分分次代数层。存在集合 $T$，并且对每个 $t\in T$
存在无环微分分次 $\mathcal{A}$-模的单射
$\mathcal{M}_t\to\mathcal{M}'_t$，使得对
$\textit{Mod}(\mathcal{A},\text{d})$ 的对象 $\mathcal{I}$，
下列条件等价：
\begin{enumerate}
\item $\mathcal{I}$ 分次内射；
\item 对每个实线图
$$
\xymatrix{
\mathcal{M}_t \ar[r] \ar[d] & \mathcal{I} \\
\mathcal{M}'_t \ar@{..>}[ru]
}
$$
在 $\textit{Mod}(\mathcal{A},\text{d})$ 中存在使图交换的虚线箭头。
\end{enumerate}
\end{lemma}

\begin{proof}
取引理 \ref{sdga-lemma-characterize-injectives} 中的
$T$ 与 $\mathcal{N}_t\to\mathcal{N}'_t$。
以 $F:\textit{Mod}(\mathcal{A},\text{d})\to
\textit{Mod}(\mathcal{A})$ 表示遗忘函子，
以 $G$ 表示引理 \ref{sdga-lemma-dg-hull} 中 $F$ 的左伴随函子。置
$$
\mathcal{M}_t = G(\mathcal{N}_t) \to
G(\mathcal{N}'_t) = \mathcal{M}'_t
$$
由引理 \ref{sdga-lemma-dg-hull-acyclic}，这是无环微分分次
$\mathcal{A}$-模的单射。由于 $G$ 是 $F$ 的左伴随，
下图中存在虚线箭头
$$
\xymatrix{
\mathcal{M}_t \ar[r] \ar[d] & \mathcal{I} \\
\mathcal{M}'_t \ar@{..>}[ru]
}
$$
当且仅当下图中存在虚线箭头：
$$
\xymatrix{
\mathcal{N}_t \ar[r] \ar[d] & F(\mathcal{I}) \\
\mathcal{N}'_t \ar@{..>}[ru]
}
$$
因此，结论由箭头族
$\mathcal{N}_t\to\mathcal{N}_t'$ 的选取方式得出。
\end{proof}


\begin{lemma}
\label{sdga-lemma-small-acyclics}
设 $(\mathcal{C},\mathcal{O})$ 是环化位点，
$(\mathcal{A},\text{d})$ 是 $(\mathcal{C},\mathcal{O})$
上的微分分次代数层。存在集合 $S$，并且对每个 $s$
存在无环微分分次 $\mathcal{A}$-模 $\mathcal{M}_s$，
使得对每个非零无环微分分次 $\mathcal{A}$-模 $\mathcal{M}$，
存在 $s\in S$ 以及 $\textit{Mod}(\mathcal{A},\text{d})$
中的单射 $\mathcal{M}_s\to\mathcal{M}$。
\end{lemma}

\begin{proof}
开始之前回忆，我们的约定保证位点 $\mathcal{C}$
的对象和态射构成集合，并有一个覆盖集合
$\text{Cov}(\mathcal{C})$。若 $\mathcal{F}$
是微分分次 $\mathcal{A}$-模，定义 $|\mathcal{F}|$ 为
$$
\coprod\nolimits_{(U, n)} \mathcal{F}^n(U)
$$
的基数之和，其中 $U$ 遍历 $\mathcal{C}$ 的对象，
$n\in\mathbf{Z}$。选取无限基数 $\kappa$，使其大于
$|\Ob(\mathcal{C})|$、$|\text{Arrows}(\mathcal{C})|$、
$|\text{Cov}(\mathcal{C})|$、所有
$\{U_i\to U\}_{i\in I}\in\text{Cov}(\mathcal{C})$
所对应的 $\sup|I|$，以及 $|\mathcal{A}|$。

\medskip\noindent
设 $\mathcal{F}\subset\mathcal{M}$ 是微分分次
$\mathcal{A}$-模的包含。设给定集合 $K$，并且对每个 $k\in K$
给定三元组 $(U_k,n_k,x_k)$，其中 $U_k$ 是 $\mathcal{C}$
的对象，$n_k$ 是整数，而
$x_k\in\mathcal{M}^{n_k}(U_k)$ 是截面。可以考虑包含
$\mathcal{F}$ 以及所有 $k\in K$ 对应截面 $x_k$ 的最小微分分次
$\mathcal{A}$-子模 $\mathcal{F}'\subset\mathcal{M}$。可以把
$$
(\mathcal{F}')^n(U) \subset \mathcal{M}^n(U)
$$
描述为由满足如下条件的元素 $x\in\mathcal{M}^n(U)$ 组成的集合：
存在 $\{f_i:U_i\to U\}_{i\in I}\in\text{Cov}(\mathcal{C})$，
使得对每个 $i\in I$，存在有限集 $T_i$ 及态射
$g_t:U_i\to U_{k_t}$，且
$$
f_i^*x = y_i +
\sum\nolimits_{t \in T_i} a_{it}g_t^*x_{k_t} + b_{it}g_t^*\text{d}(x_{k_t})
$$
，其中 $y_i\in\mathcal{F}^n(U)$，并且
$a_{it}\in\mathcal{A}^{n-n_{k_t}}(U_i)$、
$b_{it}\in\mathcal{A}^{n-n_{k_t}-1}(U_i)$。
（细节从略；提示：这些截面显然属于 $\mathcal{F}'$，
反过来则证明这条规则定义微分分次 $\mathcal{A}$-子模。）
由此描述可知
$|\mathcal{F}'| \leq \max(|\mathcal{F}|, |K|, \kappa)$.

\medskip\noindent
设 $\mathcal{M}$ 是非零无环微分分次 $\mathcal{A}$-模。
则可找到整数 $n$ 以及 $\mathcal{M}^n$ 在
$\mathcal{C}$ 某个对象 $U$ 上的非零截面 $x$。令
$$
\mathcal{F}_0 \subset \mathcal{M}
$$
为包含 $x$ 的最小微分分次 $\mathcal{A}$-子模。
由上一段，$|\mathcal{F}_0|\leq\kappa$。归纳地，给定
$\mathcal{F}_0,\ldots,\mathcal{F}_n$ 后，如下定义
$\mathcal{F}_{n+1}$。考虑三元组集合
$$
L = \{(U, n, x)\}
\{U_i \to U\}_{i \in I}, (x_i)_{i \in I})\}
$$
，其中 $U$ 是 $\mathcal{C}$ 的对象，$n\in\mathbf{Z}$，
$x\in\mathcal{F}_n(U)$ 且 $\text{d}(x)=0$。
由于 $\mathcal{M}$ 无环，对每个三元组
$l=(U_l,n_l,x_l)\in L$，可以选取
$\{(U_{l, i} \to U_l\}_{i \in I_l} \in \text{Cov}(\mathcal{C})$，以及
$x_{l, i} \in \mathcal{M}^{n_l - 1}(U_{l, i})$，使得
$\text{d}(x_{l, i}) = x|_{U_{l, i}}$。然后置
$$
K = \{(U_{l, i}, n_l - 1, x_{l, i}) \mid l \in L, i \in I_l\}
$$
并令 $\mathcal{F}_{n+1}$ 为 $\mathcal{M}$ 中包含
$\mathcal{F}_n$ 及所有截面 $x_{l,i}$ 的最小微分分次
$\mathcal{A}$-子模。由于
$|K|\leq\max(\kappa,|\mathcal{F}_n|)$，归纳可得
$|\mathcal{F}_{n+1}|\leq\kappa$。

\medskip\noindent
按构造，包含 $\mathcal{F}_n\to\mathcal{F}_{n+1}$
在上同调层上诱导零映射。因此
$\mathcal{F}=\bigcup\mathcal{F}_n$ 是非零无环子模，且
$|\mathcal{F}|\leq\kappa$。由于满足 $|\mathcal{F}|$
有界的微分分次 $\mathcal{A}$-模 $\mathcal{F}$
的同构类只构成一个集合，结论得证。
\end{proof}

\begin{definition}
\label{sdga-definition-K-injective}
设 $(\mathcal{C},\mathcal{O})$ 是环化位点，
$(\mathcal{A},\text{d})$ 是 $(\mathcal{C},\mathcal{O})$
上的微分分次代数层。若对每个无环微分分次 $\mathcal{M}$ 都有
$$
\Hom_{K(\textit{Mod}(\mathcal{A}, \text{d}))}(\mathcal{M}, \mathcal{I}) = 0
$$
，则称微分分次 $\mathcal{A}$-模 $\mathcal{I}$ {\it K-内射}。
\end{definition}

\noindent
请注意这与《导出范畴》定义
\ref{derived-definition-K-injective} 的相似之处。

\begin{lemma}
\label{sdga-lemma-product-K-injective}
设 $(\mathcal{C},\mathcal{O})$ 是环化位点，
$(\mathcal{A},\text{d})$ 是 $(\mathcal{C},\mathcal{O})$
上的微分分次代数层。设 $T$ 是集合，并且对每个 $t\in T$，
$\mathcal{I}_t$ 是 K-内射微分分次 $\mathcal{A}$-模。
则 $\prod\mathcal{I}_t$ 是 K-内射微分分次 $\mathcal{A}$-模。
\end{lemma}

\begin{proof}
设 $\mathcal{K}$ 是无环微分分次 $\mathcal{A}$-模。则有
$$
\Hom_{\textit{Mod}^{dg}(\mathcal{A}, \text{d})}(\mathcal{K},
\prod\nolimits_{t \in T} \mathcal{I}_t)
=
\prod\nolimits_{t \in T}
\Hom_{\textit{Mod}^{dg}(\mathcal{A}, \text{d})}(\mathcal{K}, \mathcal{I}_t)
$$
，因为在 $\textit{Mod}(\mathcal{A},\text{d})$ 中取直积，
与到分次 $\mathcal{A}$-模范畴的遗忘函子交换。
由于取直积在阿贝尔群范畴上是正合函子，结论得证。
\end{proof}

\begin{lemma}
\label{sdga-lemma-first-property-dg-injective}
设 $(\mathcal{C},\mathcal{O})$ 是环化位点，
$(\mathcal{A},\text{d})$ 是 $(\mathcal{C},\mathcal{O})$
上的微分分次代数层。设 $\mathcal{I}$ 是
$\textit{Mod}(\mathcal{A},\text{d})$ 中既 K-内射又分次内射的对象。
对 $\textit{Mod}(\mathcal{A},\text{d})$ 中任意实线图
$$
\xymatrix{
\mathcal{M} \ar[r]_a \ar[d]_b & \mathcal{I} \\
\mathcal{M}' \ar@{..>}[ru]
}
$$
，若 $b$ 是单射且 $\mathcal{M}$ 无环，则存在使图交换的虚线箭头。
\end{lemma}

\begin{proof}
由于 $\mathcal{M}$ 无环且 $\mathcal{I}$ K-内射，
存在次数为 $-1$ 的分次 $\mathcal{A}$-模映射
$h:\mathcal{M}\to\mathcal{I}$，使得 $a=\text{d}(h)$。
由于 $\mathcal{I}$ 分次内射且 $b$ 是单射，
存在次数为 $-1$ 的分次 $\mathcal{A}$-模映射
$h':\mathcal{M}'\to\mathcal{I}$，使得 $h=h'\circ b$。
于是可取 $a'=\text{d}(h')$ 为虚线箭头。
\end{proof}

\begin{lemma}
\label{sdga-lemma-second-property-dg-injective}
设 $(\mathcal{C},\mathcal{O})$ 是环化位点，
$(\mathcal{A},\text{d})$ 是 $(\mathcal{C},\mathcal{O})$
上的微分分次代数层。设 $\mathcal{I}$ 是
$\textit{Mod}(\mathcal{A},\text{d})$ 中既 K-内射又分次内射的对象。
对 $\textit{Mod}(\mathcal{A},\text{d})$ 中任意实线图
$$
\xymatrix{
\mathcal{M} \ar[r]_a \ar[d]_b & \mathcal{I} \\
\mathcal{M}' \ar@{..>}[ru]
}
$$
，若 $b$ 是拟同构，则存在使图在同伦意义下交换的虚线箭头。
\end{lemma}

\begin{proof}
把 $\mathcal{M}'$ 替换为 $\mathcal{M}'$ 与
$\mathcal{M}$ 恒等映射之锥（后者无环）的直和后，
可设 $b$ 也是单射。此时 $b$ 的余核 $\mathcal{Q}$ 无环。因此
$$
\Hom_{K(\textit{Mod}(\mathcal{A}, \text{d}))}(\mathcal{Q}, \mathcal{I}) =
\Hom_{K(\textit{Mod}(\mathcal{A}, \text{d}))}(\mathcal{Q}, \mathcal{I})[1] = 0
$$
，因为 $\mathcal{I}$ K-内射。又因 $\mathcal{I}$ 分次内射，
由评注 \ref{sdga-remark-why-graded-injective} 可知
$$
\Hom_{K(\textit{Mod}(\mathcal{A}, \text{d}))}(\mathcal{M}', \mathcal{I})
\longrightarrow
\Hom_{K(\textit{Mod}(\mathcal{A}, \text{d}))}(\mathcal{M}, \mathcal{I})
$$
是双射，证明完成。
\end{proof}

\begin{lemma}
\label{sdga-lemma-better-set-of-monos}
设 $(\mathcal{C},\mathcal{O})$ 是环化位点，
$(\mathcal{A},\text{d})$ 是 $(\mathcal{C},\mathcal{O})$
上的微分分次代数层。存在集合 $R$，并且对每个 $r\in R$
存在无环微分分次 $\mathcal{A}$-模的单射
$\mathcal{M}_r\to\mathcal{M}'_r$，使得对
$\textit{Mod}(\mathcal{A},\text{d})$ 的对象 $\mathcal{I}$，
下列条件等价：
\begin{enumerate}
\item $\mathcal{I}$ 既 K-内射又分次内射；
\item 对每个实线图
$$
\xymatrix{
\mathcal{M}_r \ar[r] \ar[d] & \mathcal{I} \\
\mathcal{M}'_r \ar@{..>}[ru]
}
$$
在 $\textit{Mod}(\mathcal{A},\text{d})$ 中存在使图交换的虚线箭头。
\end{enumerate}
\end{lemma}

\begin{proof}
取引理 \ref{sdga-lemma-characterize-graded-injectives-in-dg} 中的
$T$ 与 $\mathcal{M}_t\to\mathcal{M}'_t$，并取引理
\ref{sdga-lemma-small-acyclics} 中的 $S$ 与 $\mathcal{M}_s$。
选取无环微分分次 $\mathcal{A}$-模的单射
$\mathcal{M}_s\to\mathcal{M}'_s$，使其零同伦。
这是可能的，因为可取 $\mathcal{M}'_s$ 为恒等映射的锥；
在此情形下，甚至 $\mathcal{M}'_s$ 上的恒等映射也零同伦，
参见《微分分次代数》引理 \ref{dga-lemma-id-cone-null}；
由第 \ref{sdga-section-conclude-triangulated} 节的讨论，该引理适用。
我们断言，$R=T\coprod S$ 连同给定映射满足要求。

\medskip\noindent
蕴含 (1) $\Rightarrow$ (2) 由引理
\ref{sdga-lemma-first-property-dg-injective} 得出。

\medskip\noindent
假设 (2)。首先，由引理
\ref{sdga-lemma-characterize-graded-injectives-in-dg} 可知
$\mathcal{I}$ 分次内射。其次，设 $\mathcal{M}$
是无环微分分次 $\mathcal{A}$-模。需要证明
$$
\Hom_{K(\textit{Mod}(\mathcal{A}, \text{d}))}(\mathcal{M}, \mathcal{I}) = 0
$$
其证明与《内射对象》引理
\ref{injectives-lemma-characterize-K-injective} 的证明完全相同。

\medskip\noindent
我们将按序数 $\alpha$ 归纳构造无环微分分次子模
$\mathcal{K}_\alpha\subset\mathcal{M}$。当 $\alpha=0$ 时，
置 $\mathcal{K}_0=0$。当 $\alpha>0$ 时，按如下方式进行：
\begin{enumerate}
\item 若 $\alpha=\beta+1$ 且 $\mathcal{K}_\beta=\mathcal{M}$，
则取 $\mathcal{K}_\alpha=\mathcal{K}_\beta$。
\item 若 $\alpha=\beta+1$ 且 $\mathcal{K}_\beta\ne\mathcal{M}$，
则 $\mathcal{M}/\mathcal{K}_\beta$ 是非零无环微分分次
$\mathcal{A}$-模。选取微分分次 $\mathcal{A}$-子模
$\mathcal{N}_\alpha\subset\mathcal{M}/\mathcal{K}_\beta$，
使其同构于某个 $s\in S$ 所对应的 $\mathcal{M}_s$；
参见引理 \ref{sdga-lemma-small-acyclics}。最后令
$\mathcal{K}_\alpha\subset\mathcal{M}$ 为 $\mathcal{N}_\alpha$ 的逆像。
\item 若 $\alpha$ 是极限序数，则置
$\mathcal{K}_\beta = \colim \mathcal{K}_\alpha$.
\end{enumerate}
显然，对某个充分大的序数 $\alpha$，有
$\mathcal{M}=\mathcal{K}_\alpha$。我们将对 $\alpha$
作超限归纳，证明
$$
\Hom_{K(\textit{Mod}(\mathcal{A}, \text{d}))}(\mathcal{K}_\alpha, \mathcal{I})
$$
为零。当 $\alpha=0$ 时成立，因为 $\mathcal{K}_0=0$。
假设对 $\beta$ 成立，且 $\alpha=\beta+1$。
在上述 (1) 情形中结论显然。在 (2) 情形中有短正合序列
$$
0 \to \mathcal{K}_\beta \to \mathcal{K}_\alpha \to \mathcal{N}_\alpha \to 0
$$
。由评注 \ref{sdga-remark-why-graded-injective} 以及已经证明的
$\mathcal{I}$ 分次内射性，得到正合序列
$$
\Hom_{K(\textit{Mod}(\mathcal{A}, \text{d}))}(\mathcal{K}_\beta, \mathcal{I})
\to
\Hom_{K(\textit{Mod}(\mathcal{A}, \text{d}))}(\mathcal{K}_\alpha, \mathcal{I})
\to
\Hom_{K(\textit{Mod}(\mathcal{A}, \text{d}))}(\mathcal{N}_\alpha, \mathcal{I})
$$
。左项由归纳假设为零。右项由假设 (2) 为零：任意映射
$\mathcal{M}_s\to\mathcal{I}$ 都经 $\mathcal{M}'_s$ 分解，
因而零同伦。所以中间群也为零。最后，设 $\alpha$ 是极限序数。
由于作为分次 $\mathcal{A}$-模还有
$\mathcal{K}_\alpha = \colim \mathcal{K}_\alpha$
，可见
$$
\Hom_{\textit{Mod}^{dg}(\mathcal{A}, \text{d})}
(\mathcal{K}_\alpha, \mathcal{I})
= \lim_{\beta < \alpha}
\Hom_{\textit{Mod}^{dg}(\mathcal{A}, \text{d})}
(\mathcal{K}_\beta, \mathcal{I})
$$
是阿贝尔群复形的等式。这些复形的上同调群计算
$K(\textit{Mod}(\mathcal{A},\text{d}))$ 中各平移之间的态射。
由于 $\mathcal{I}$ 分次内射，由评注
\ref{sdga-remark-why-graded-injective}，此复形系统中的转移映射为满射。
此外，对极限序数 $\beta\leq\alpha$，极限与该处的值相等。
故可应用《同调》引理 \ref{homology-lemma-ML-over-ordinals} 得出结论。
\end{proof}

\begin{lemma}
\label{sdga-lemma-functor-set-of-monos}
设 $(\mathcal{C},\mathcal{O})$ 是环化位点，
$(\mathcal{A},\text{d})$ 是 $(\mathcal{C},\mathcal{O})$
上的微分分次代数层。设 $R$ 是集合，并且对每个 $r\in R$
给定无环微分分次 $\mathcal{A}$-模的单射
$\mathcal{M}_r\to\mathcal{M}'_r$。存在函子
$M:\textit{Mod}(\mathcal{A},\text{d})\to
\textit{Mod}(\mathcal{A},\text{d})$ 与自然变换
$j:\text{id}\to M$，使得
\begin{enumerate}
\item $j_\mathcal{M}:\mathcal{M}\to M(\mathcal{M})$
是单射且为拟同构；
\item 对每个实线图
$$
\xymatrix{
\mathcal{M}_r \ar[r] \ar[d] & \mathcal{M} \ar[d]^{j_\mathcal{M}} \\
\mathcal{M}'_r \ar@{..>}[r] & M(\mathcal{M})
}
$$
在 $\textit{Mod}(\mathcal{A},\text{d})$ 中存在使图交换的虚线箭头。
\end{enumerate}
\end{lemma}

\begin{proof}
把 $M(\mathcal{M})$ 定义为下图中的推出：
$$
\xymatrix{
\bigoplus_{(r, \varphi)} \mathcal{M}_r \ar[r] \ar[d] &
\mathcal{M} \ar[d] \\
\bigoplus_{(r, \varphi)} \mathcal{M}'_r \ar[r] &
M(\mathcal{M})
}
$$
其中直和遍历满足 $r\in R$ 以及
$\varphi\in
\Hom_{\textit{Mod}(\mathcal{A}, \text{d})}(\mathcal{M}_r, \mathcal{M})$.
的所有偶对 $(r,\varphi)$。由于单射的推出仍为单射，
$\mathcal{M}\to M(\mathcal{M})$ 是单射。由于左侧竖直箭头
的余核无环，$\mathcal{M}\to M(\mathcal{M})$
与之同构的余核也无环，所以该映射是拟同构。
性质 (2) 由构造成立。此过程确实可构成函子的验证从略。
\end{proof}

\begin{theorem}
\label{sdga-theorem-qis-into-dg-injective}
设 $(\mathcal{C},\mathcal{O})$ 是环化位点，
$(\mathcal{A},\text{d})$ 是 $(\mathcal{C},\mathcal{O})$
上的微分分次代数层。对每个微分分次 $\mathcal{A}$-模
$\mathcal{M}$，存在拟同构 $\mathcal{M}\to\mathcal{I}$，
其中 $\mathcal{I}$ 是既分次内射又 K-内射的微分分次
$\mathcal{A}$-模。此外，该构造关于 $\mathcal{M}$ 具有函子性。
\end{theorem}

\begin{proof}
取引理 \ref{sdga-lemma-better-set-of-monos} 中找到的集合 $R$
以及 $\textit{Mod}(\mathcal{A},\text{d})$ 的态射
$\mathcal{M}_r\to\mathcal{M}'_r$。取引理
\ref{sdga-lemma-functor-set-of-monos} 利用这些数据构造的
$M$ 与自然变换 $\text{id}\to M$。使用超限递归，
如下定义函子序列 $M_\alpha$ 以及当 $\alpha<\beta$
时的自然变换 $M_\beta\to M_\alpha$：
\begin{enumerate}
\item $M_0=\text{id}$；
\item $M_{\alpha+1}=M\circ M_\alpha$；当 $\beta<\alpha+1$ 时，
自然变换 $M_\beta\to M_{\alpha+1}$ 来自已经构造的
$M_\beta\to M_\alpha$，以及由 $\text{id}\to M$ 所得的映射
$M_\alpha\to M\circ M_\alpha$；
\item 若 $\alpha$ 是极限序数，则
$M_\alpha=\colim_{\beta<\alpha}M_\beta$，并以余投射作为
当 $\alpha<\beta$ 时的变换 $M_\beta\to M_\alpha$。
\end{enumerate}
注意，对每个微分分次 $\mathcal{A}$-模，映射
$\mathcal{M} \to M_\beta(\mathcal{M}) \to M_\alpha(\mathcal{M})$
都是单的拟同构（因为滤过余极限正合）。

\medskip\noindent
回忆 $\textit{Mod}(\mathcal{A},\text{d})$ 是格罗滕迪克阿贝尔范畴。
因此，由《内射对象》命题
\ref{injectives-proposition-objects-are-small}
（应用于所有 $r\in R$ 对应的 $\mathcal{M}_r$ 的直和），
存在极限序数 $\alpha$，使每个 $r\in R$ 对应的
$\mathcal{M}_r$ 关于单射都是 $\alpha$-小的。
我们断言 $\mathcal{M}\to M_\alpha(\mathcal{M})$
正是所需的函子性嵌入，它把 $\mathcal{M}$
嵌入一个分次内射且 K-内射的模中。

\medskip\noindent
事实上，任意映射
$\mathcal{M}_r\to M_\alpha(\mathcal{M})$ 都经某个
$\beta<\alpha$ 对应的 $M_\beta(\mathcal{M})$ 分解。
而由 $M$ 的构造，这意味着映射
$\mathcal{M}_r\to M_{\beta+1}(\mathcal{M})
=M(M_\beta(\mathcal{M}))$ 经 $\mathcal{M}'_r$ 分解。
由于
$M_\beta(\mathcal{M})\subset M_{\beta+1}(\mathcal{M})
\subset M_\alpha(\mathcal{M})$，便得到所需的到
$M_\alpha(\mathcal{M})$ 的分解。最后由引理
\ref{sdga-lemma-better-set-of-monos} 中对
$R$ 及 $\mathcal{M}_r\to\mathcal{M}'_r$ 的选取得出结论。
\end{proof}










\section{导出范畴}
\label{sdga-section-derived}

\noindent
本节对应《微分分次代数》第 \ref{dga-section-derived} 节。

\medskip\noindent
设 $(\mathcal{C},\mathcal{O})$ 是环化位点，
$(\mathcal{A},\text{d})$ 是 $(\mathcal{C},\mathcal{O})$
上的微分分次代数层。我们将通过逆转
$K(\textit{Mod}(\mathcal{A},\text{d}))$ 中的拟同构来构造
导出范畴 $D(\mathcal{A},\text{d})$。

\begin{lemma}
\label{sdga-lemma-cohomology-homological}
设 $(\mathcal{C},\mathcal{O})$ 是环化位点，
$(\mathcal{A},\text{d})$ 是 $(\mathcal{C},\mathcal{O})$
上的微分分次代数层。第 \ref{sdga-section-modules} 节的函子
$H^0 : \textit{Mod}(\mathcal{A}, \text{d}) \to \textit{Mod}(\mathcal{O})$
经函子
$$
H^0 : K(\textit{Mod}(\mathcal{A}, \text{d})) \to \textit{Mod}(\mathcal{O})
$$
分解；后一函子在《导出范畴》定义
\ref{derived-definition-homological} 的意义下是同调函子。
\end{lemma}

\begin{proof}
由定义立即得到交换图
$$
\xymatrix{
\textit{Mod}(\mathcal{A}, \text{d}) \ar[r] \ar[d] &
K(\textit{Mod}(\mathcal{A}, \text{d})) \ar[d] \\
\text{Comp}(\mathcal{O}) \ar[r] &
K(\textit{Mod}(\mathcal{O}))
}
$$
由于 $H^0(\mathcal{M})$ 定义为 $\mathcal{M}$ 的底层
$\mathcal{O}$-模复形的第零上同调层，本引理由
$\mathcal{O}$-模复形的情形得出；后者是《导出范畴》引理
\ref{derived-lemma-cohomology-homological} 的特例。
\end{proof}

\begin{lemma}
\label{sdga-lemma-acyclics}
设 $(\mathcal{C},\mathcal{O})$ 是环化位点，
$(\mathcal{A},\text{d})$ 是 $(\mathcal{C},\mathcal{O})$
上的微分分次代数层。同伦范畴
$K(\textit{Mod}(\mathcal{A},\text{d}))$ 中由无环模组成的
全子范畴 $\text{Ac}$，是
$K(\textit{Mod}(\mathcal{A},\text{d}))$
的严格全饱和三角子范畴。
\end{lemma}

\begin{proof}
当然，$K(\textit{Mod}(\mathcal{A},\text{d}))$ 的对象
$\mathcal{M}$ 属于 $\text{Ac}$，当且仅当对所有 $i$，
$H^i(\mathcal{M})=H^0(\mathcal{M}[i])$ 均为零。
本引理由此、引理 \ref{sdga-lemma-cohomology-homological}
及《导出范畴》引理 \ref{derived-lemma-homological-functor-kernel}
得出。另见《导出范畴》定义 \ref{derived-definition-saturated}、
\ref{derived-definition-triangulated-subcategory}
及引理 \ref{derived-lemma-triangulated-subcategory}。
\end{proof}

\begin{lemma}
\label{sdga-lemma-qis}
设 $(\mathcal{C},\mathcal{O})$ 是环化位点，
$(\mathcal{A},\text{d})$ 是 $(\mathcal{C},\mathcal{O})$
上的微分分次代数层。考虑由拟同构组成的子类
$\text{Qis} \subset \text{Arrows}(K(\textit{Mod}(\mathcal{A}, \text{d})))$
。这是一个饱和乘法系，并与
$K(\textit{Mod}(\mathcal{A},\text{d}))$ 上的三角结构相容。
\end{lemma}

\begin{proof}
注意，若 $f,g:\mathcal{M}\to\mathcal{N}$ 是
$\textit{Mod}(\mathcal{A},\text{d})$ 中彼此同伦的态射，
则 $f$ 是拟同构当且仅当 $g$ 是拟同构。事实上，由引理
\ref{sdga-lemma-cohomology-homological}，映射
$H^i(f)=H^0(f[i])$ 与 $H^i(g)=H^0(g[i])$ 相同。
因此，说同伦范畴 $K(\textit{Mod}(\mathcal{A},\text{d}))$
的某个态射是拟同构并无歧义。关于``乘法系''、``饱和''
及``与三角结构相容''的定义，参见《导出范畴》定义
\ref{derived-definition-localization}，以及《范畴》定义
\ref{categories-definition-multiplicative-system}、
\ref{categories-definition-saturated-multiplicative-system}。

\medskip\noindent
为实际证明本引理，考虑三角范畴间正合函子的复合
$$
K(\textit{Mod}(\mathcal{A}, \text{d}))
\longrightarrow
K(\textit{Mod}(\mathcal{O}))
\longrightarrow
D(\mathcal{O})
$$
。注意，$K(\textit{Mod}(\mathcal{A},\text{d}))$ 的态射
$f:\mathcal{M}\to\mathcal{N}$ 属于 $\text{Qis}$，
当且仅当它在 $D(\mathcal{O})$ 中映为同构。
故本引理由《导出范畴》引理
\ref{derived-lemma-triangle-functor-localize} 得出。
\end{proof}

\noindent
在引理 \ref{sdga-lemma-qis} 的情形下，可以应用《导出范畴》命题
\ref{derived-proposition-construct-localization}，
得到三角范畴间的正合函子
$$
Q :
K(\textit{Mod}(\mathcal{A}, \text{d}))
\longrightarrow
\text{Qis}^{-1}K(\textit{Mod}(\mathcal{A}, \text{d}))
$$
。然而，由于 $\textit{Mod}(\mathcal{A},\text{d})$ 是``大''
范畴，即其对象构成真类，给定 $\mathcal{M}$ 与 $\mathcal{N}$ 后，
$\text{Qis}^{-1}K(\textit{Mod}(\mathcal{A},\text{d}))$
的构造是否产生态射的一个{\it 集合}
$$
\Mor_{\text{Qis}^{-1}K(\textit{Mod}(\mathcal{A}, \text{d}))}
(\mathcal{M}, \mathcal{N})
$$
，并非立即清楚。不过，借助我们对 K-内射复形的构造，
答案是肯定的。事实上，由定理
\ref{sdga-theorem-qis-into-dg-injective}，可选取拟同构
$s_0:\mathcal{N}\to\mathcal{I}$，其中 $\mathcal{I}$
是既分次内射又 K-内射的微分分次 $\mathcal{A}$-模。
回忆，上述集合的元素是偶对
$(f : \mathcal{M} \to \mathcal{N}', s : \mathcal{N} \to \mathcal{N}')$
的等价类，其中 $f$ 是
$K(\textit{Mod}(\mathcal{A},\text{d}))$ 的任意态射，
$s$ 是拟同构；参见《范畴》第
\ref{categories-section-localization} 节对左分式演算的描述。
由引理 \ref{sdga-lemma-second-property-dg-injective}，可选取下图中的虚线箭头
$$
\xymatrix{
\mathcal{M} \ar[rd]^f & &
\mathcal{N} \ar[ld]_s \ar[rd]^{s_0} \\
& \mathcal{N}' \ar@{..>}[rr]^{s'} & & \mathcal{I}
}
$$
使图在同伦范畴中交换。因此，偶对 $(f,s)$ 等价于
$(s'\circ f,s_0)$，从而这些等价类构成集合。

\begin{definition}
\label{sdga-definition-derived-category}
设 $(\mathcal{C},\mathcal{O})$ 是环化位点，
$(\mathcal{A},\text{d})$ 是 $(\mathcal{C},\mathcal{O})$
上的微分分次代数层。令 $\text{Qis}$ 如引理
\ref{sdga-lemma-qis} 中所述。{\it $(\mathcal{A},\text{d})$ 的导出范畴}
是上文详细讨论的三角范畴
$$
D(\mathcal{A}, \text{d}) =
\text{Qis}^{-1}K(\textit{Mod}(\mathcal{A}, \text{d}))
$$
\end{definition}

\noindent
下面证明关于这一构造的一些事实。

\begin{lemma}
\label{sdga-lemma-kernel-localization}
在定义 \ref{sdga-definition-derived-category} 中，局部化函子
$Q : K(\textit{Mod}(\mathcal{A}, \text{d})) \to D(\mathcal{A}, \text{d})$
的核是引理 \ref{sdga-lemma-acyclics} 中的范畴 $\text{Ac}$。
\end{lemma}

\begin{proof}
这立即来自《导出范畴》引理
\ref{derived-lemma-kernel-localization}，以及如下事实：
$0\to\mathcal{M}$ 是拟同构，当且仅当 $\mathcal{M}$ 无环。
\end{proof}

\begin{lemma}
\label{sdga-lemma-H0-over-D}
在定义 \ref{sdga-definition-derived-category} 中，函子
$H^0 : K(\textit{Mod}(\mathcal{A}, \text{d})) \to
\textit{Mod}(\mathcal{O})$ 经同调函子
$H^0 : D(\mathcal{A}, \text{d}) \to \textit{Mod}(\mathcal{O})$.
分解。
\end{lemma}

\begin{proof}
立即来自《导出范畴》引理
\ref{derived-lemma-universal-property-localization}。
\end{proof}

\noindent
下面给出此前承诺的计算导出范畴中态射集的引理。

\begin{lemma}
\label{sdga-lemma-hom-derived}
设 $(\mathcal{C},\mathcal{O})$ 是环化位点，
$(\mathcal{A},\text{d})$ 是 $(\mathcal{C},\mathcal{O})$
上的微分分次代数层。设 $\mathcal{M}$ 与 $\mathcal{N}$
是微分分次 $\mathcal{A}$-模，并设
$\mathcal{N}\to\mathcal{I}$ 是拟同构，其中 $\mathcal{I}$
是既分次内射又 K-内射的微分分次 $\mathcal{A}$-模。则
$$
\Hom_{D(\mathcal{A}, \text{d})}(\mathcal{M}, \mathcal{N}) =
\Hom_{K(\textit{Mod}(\mathcal{A}, \text{d}))}(\mathcal{M}, \mathcal{I})
$$
\end{lemma}

\begin{proof}
由于 $\mathcal{N}\to\mathcal{I}$ 是拟同构，有
$$
\Hom_{D(\mathcal{A}, \text{d})}(\mathcal{M}, \mathcal{N}) =
\Hom_{D(\mathcal{A}, \text{d})}(\mathcal{M}, \mathcal{I})
$$
。在定义 \ref{sdga-definition-derived-category} 之前的讨论中，
利用引理 \ref{sdga-lemma-second-property-dg-injective} 已经得到：
$D(\mathcal{A},\text{d})$ 中任意态射
$\mathcal{M}\to\mathcal{I}$ 都可由
$K(\textit{Mod}(\mathcal{A},\text{d}))$ 中的态射
$f:\mathcal{M}\to\mathcal{I}$ 表示。现在，若
$f,f':\mathcal{M}\to\mathcal{I}$ 是
$K(\textit{Mod}(\mathcal{A},\text{d}))$ 中的两个态射，
则它们在 $D(\mathcal{A},\text{d})$ 中定义同一态射，
当且仅当存在 $K(\textit{Mod}(\mathcal{A},\text{d}))$
中的拟同构 $g:\mathcal{I}\to\mathcal{K}$，使
$g\circ f=g\circ f'$；参见《范畴》引理
\ref{categories-lemma-equality-morphisms-left-localization}。
然而，由引理 \ref{sdga-lemma-second-property-dg-injective}，
存在映射 $h:\mathcal{K}\to\mathcal{I}$，使得在
$K(\textit{Mod}(\mathcal{A},\text{d}))$ 中
$h\circ g=\text{id}_\mathcal{I}$。故
$g\circ f=g\circ f'$ 蕴含 $f=f'$，证明完成。
\end{proof}

\begin{lemma}
\label{sdga-lemma-derived-products}
设 $(\mathcal{C},\mathcal{O})$ 是环化位点，
$(\mathcal{A},\text{d})$ 是 $(\mathcal{C},\mathcal{O})$
上的微分分次代数层。则
\begin{enumerate}
\item $D(\mathcal{A},\text{d})$ 既有直和也有直积；
\item 直和通过取微分分次 $\mathcal{A}$-模的直和得到；
\item 直积通过取 K-内射微分分次模的直积得到。
\end{enumerate}
\end{lemma}

\begin{proof}
我们将使用如下事实：
$\textit{Mod}(\mathcal{A},\text{d})$ 是具有任意直和与直积的
阿贝尔范畴，而且这些运算在
$K(\textit{Mod}(\mathcal{A},\text{d}))$ 中给出直和与直积。
参见引理 \ref{sdga-lemma-dgm-abelian} 与
\ref{sdga-lemma-homotopy-direct-sums}。

\medskip\noindent
设 $\mathcal{M}_j$ 是一族微分分次 $\mathcal{A}$-模。
把直和 $\mathcal{M}=\bigoplus\mathcal{M}_j$
视为微分分次 $\mathcal{A}$-模。对微分分次
$\mathcal{A}$-模 $\mathcal{N}$，选取拟同构
$\mathcal{N}\to\mathcal{I}$，其中 $\mathcal{I}$
作为微分分次 $\mathcal{A}$-模既分次内射又 K-内射；
参见定理 \ref{sdga-theorem-qis-into-dg-injective}。
利用引理 \ref{sdga-lemma-hom-derived}，有
\begin{align*}
\Hom_{D(\mathcal{A}, \text{d})}(\mathcal{M}, \mathcal{N})
& =
\Hom_{K(\mathcal{A}, \text{d})}(\mathcal{M}, \mathcal{I}) \\
& =
\prod \Hom_{K(\mathcal{A}, \text{d})}(\mathcal{M}_j, \mathcal{I}) \\
& =
\prod \Hom_{D(\mathcal{A}, \text{d})}(\mathcal{M}_j, \mathcal{I})
\end{align*}
由此得到 $D(A,\text{d})$ 中按本引理 (2) 所述的直和。

\medskip\noindent
设 $\mathcal{M}_j$ 是一族微分分次 $\mathcal{A}$-模。
对每个 $j$，选取拟同构
$\mathcal{M}\to\mathcal{I}_j$，其中 $\mathcal{I}_j$
作为微分分次 $\mathcal{A}$-模既分次内射又 K-内射。
考虑微分分次 $\mathcal{A}$-模的直积
$\mathcal{I}=\prod\mathcal{I}_j$。由引理
\ref{sdga-lemma-product-K-injective} 与
\ref{sdga-lemma-product-graded-injective}，$\mathcal{I}$
作为微分分次 $\mathcal{A}$-模既分次内射又 K-内射。
对微分分次 $\mathcal{A}$-模 $\mathcal{N}$，
利用引理 \ref{sdga-lemma-hom-derived} 有
\begin{align*}
\Hom_{D(\mathcal{A}, \text{d})}(\mathcal{N}, \mathcal{I})
& =
\Hom_{K(\mathcal{A}, \text{d})}(\mathcal{N}, \mathcal{I}) \\
& =
\prod \Hom_{K(\mathcal{A}, \text{d})}(\mathcal{N}, \mathcal{I}_j) \\
& =
\prod \Hom_{D(\mathcal{A}, \text{d})}(\mathcal{N}, \mathcal{M}_j)
\end{align*}
由此得到 $D(\mathcal{A},\text{d})$ 中按本引理 (3) 所述的直积。
\end{proof}







\section{典范 $\delta$-函子}
\label{sdga-section-canonical-delta-functor}

\noindent
设 $(\mathcal{C},\mathcal{O})$ 是环化位点，
$(\mathcal{A},\text{d})$ 是 $(\mathcal{C},\mathcal{O})$
上的微分分次代数层。考虑函子
$\textit{Mod}(\mathcal{A}, \text{d}) \to
K(\textit{Mod}(\mathcal{A}, \text{d}))$.
一般而言，该函子{\bf 不是} $\delta$-函子。
然而，函子
$\textit{Mod}(\mathcal{A},\text{d})\to D(A,\text{d})$
却是 $\delta$-函子。为说明这一点，必须定义与阿贝尔范畴
$\textit{Mod}(\mathcal{A},\text{d})$ 中短正合序列
$$
0 \to \mathcal{K} \xrightarrow{a} \mathcal{L} \xrightarrow{b} \mathcal{M} \to 0
$$
相伴的态射 $\delta$。考虑态射 $a$ 的锥 $C(a)$，
连同其典范态射 $i:\mathcal{L}\to C(a)$ 与
$p:C(a)\to\mathcal{K}[1]$；参见定义 \ref{sdga-definition-cone}。
存在微分分次 $\mathcal{A}$-模同态
$$
q : C(a) \longrightarrow \mathcal{M}
$$
，它由《微分分次代数》引理 \ref{dga-lemma-cone}
应用于图
$$
\xymatrix{
\mathcal{K} \ar[r]_a \ar[d] &
\mathcal{L} \ar[d]^b \\
0 \ar[r] &
\mathcal{M}
}
$$
得到（由第 \ref{sdga-section-conclude-triangulated} 节的讨论，
我们可以使用该引理）。映射 $q$ 是拟同构；例如，这是因为
在 $\mathcal{O}$-模复形的态射范畴中该结论成立，参见
《导出范畴》第 \ref{derived-section-canonical-delta-functor} 节的讨论。
根据《微分分次代数》引理 \ref{dga-lemma-cone-homotopy}
（由第 \ref{sdga-section-conclude-triangulated} 节的讨论，
我们可以使用该引理），三角
$$
(\mathcal{K}, \mathcal{L}, C(a), a, i, -p)
$$
是 $K(\textit{Mod}(\mathcal{A},\text{d}))$ 中的杰出三角。
由于局部化函子
$K(\textit{Mod}(\mathcal{A},\text{d}))\to
D(\mathcal{A},\text{d})$ 正合，
$(\mathcal{K},\mathcal{L},C(a),a,i,-p)$ 也是
$D(\mathcal{A},\text{d})$ 中的杰出三角。
由于 $q$ 是拟同构，它在 $D(\mathcal{A},\text{d})$ 中是同构。
因此推出
$$
(\mathcal{K}, \mathcal{L}, \mathcal{M}, a, b, -p \circ q^{-1})
$$
是 $D(\mathcal{A},\text{d})$ 中的杰出三角。这提示了下列引理。

\begin{lemma}
\label{sdga-lemma-derived-canonical-delta-functor}
设 $(\mathcal{C},\mathcal{O})$ 是环化位点，
$(\mathcal{A},\text{d})$ 是 $(\mathcal{C},\mathcal{O})$
上的微分分次代数层。局部化函子
$\textit{Mod}(\mathcal{A}, \text{d}) \to D(\mathcal{A}, \text{d})$
具有 $\delta$-函子的自然结构，其中
$$
\delta_{\mathcal{K} \to \mathcal{L} \to \mathcal{M}} = - p \circ q^{-1}
$$
，$p$ 与 $q$ 如上所述。
\end{lemma}

\begin{proof}
我们已经看到，对任意复形短正合序列，这一选择给出杰出三角。
还需证明此构造的函子性；参见《导出范畴》定义
\ref{derived-definition-delta-functor}。
稍作推演，这可由《微分分次代数》引理 \ref{dga-lemma-cone}
得出（由第 \ref{sdga-section-conclude-triangulated} 节的讨论，
我们可以使用该引理）。可比较《导出范畴》引理
\ref{derived-lemma-derived-canonical-delta-functor}。
\end{proof}

\begin{lemma}
\label{sdga-lemma-homotopy-colimit}
设 $(\mathcal{C},\mathcal{O})$ 是环化位点，
$(\mathcal{A},\text{d})$ 是 $(\mathcal{C},\mathcal{O})$
上的微分分次代数层。设 $\mathcal{M}_n$
是微分分次 $\mathcal{A}$-模系统。则
$D(\mathcal{A},\text{d})$ 中的导出余极限
$\text{hocolim}\mathcal{M}_n$ 由微分分次模
$\colim\mathcal{M}_n$ 表示。
\end{lemma}

\begin{proof}
置 $\mathcal{M}=\colim\mathcal{M}_n$。由《导出范畴》引理
\ref{derived-lemma-compute-colimit}
（应用于底层 $\mathcal{O}$-模复形），有微分分次
$\mathcal{A}$-模的正合序列
$$
0 \to \bigoplus \mathcal{M}_n \to \bigoplus \mathcal{M}_n \to \mathcal{M} \to 0
$$
。由引理 \ref{sdga-lemma-derived-products}，这些直和也是
$D(\mathcal{A},\text{d})$ 中的直和。因此，结论来自
《导出范畴》定义 \ref{derived-definition-derived-colimit}
中的导出余极限定义，以及复形短正合序列给出杰出三角这一事实
（引理 \ref{sdga-lemma-derived-canonical-delta-functor}）。
\end{proof}






\section{导出拉回}
\label{sdga-section-derived-pullback}

\noindent
设 $(f, f^\sharp) : (\Sh(\mathcal{C}), \mathcal{O}_\mathcal{C})
\to (\Sh(\mathcal{D}), \mathcal{O}_\mathcal{D})$
是环化拓扑斯态射。设 $\mathcal{A}$ 是微分分次
$\mathcal{O}_\mathcal{C}$-代数，$\mathcal{B}$ 是微分分次
$\mathcal{O}_\mathcal{D}$-代数。设给定微分分次
$f^{-1}\mathcal{O}_\mathcal{D}$-代数映射
$$
\varphi : f^{-1}\mathcal{B} \to \mathcal{A}
$$
。由标量限制与标量扩张的伴随性，这等价于微分分次
$\mathcal{O}_\mathcal{C}$-代数映射
$\varphi:f^*\mathcal{B}\to\mathcal{A}$；等价地，
$\varphi$ 可视为微分分次 $\mathcal{O}_\mathcal{D}$-代数映射
$$
\varphi : \mathcal{B} \to f_*\mathcal{A}
$$
。参见评注 \ref{sdga-remark-functoriality-dga}。

\medskip\noindent
除上述数据外，再设 $\mathcal{A}'$ 是另一个微分分次
$\mathcal{O}_\mathcal{C}$-代数，$\mathcal{N}$ 是微分分次
$(\mathcal{A},\mathcal{A}')$-双模。在此情形下可考虑函子
$$
\textit{Mod}(\mathcal{B}, \text{d})
\longrightarrow
\textit{Mod}(\mathcal{A}', \text{d}),\quad
\mathcal{M} \longmapsto f^*\mathcal{M} \otimes_{\mathcal{A}} \mathcal{N}
$$
注意，由第 \ref{sdga-section-functoriality-dg} 与
\ref{sdga-section-dg-bimodules} 节的讨论，它延拓为微分分次范畴间的函子
$$
\textit{Mod}^{dg}(\mathcal{B}, \text{d})
\longrightarrow
\textit{Mod}^{dg}(\mathcal{A}', \text{d}),\quad
\mathcal{M} \longmapsto f^*\mathcal{M} \otimes_{\mathcal{A}} \mathcal{N}
$$
。形式地随之得到三角范畴间的正合函子
\begin{equation}
\label{sdga-equation-pullback}
K(\textit{Mod}(\mathcal{B}, \text{d}))
\longrightarrow
K(\textit{Mod}(\mathcal{A}', \text{d})),\quad
\mathcal{M} \longmapsto f^*\mathcal{M} \otimes_{\mathcal{A}} \mathcal{N}
\end{equation}
。

\begin{lemma}
\label{sdga-lemma-derived-tensor-product}
在上述情形下，函子 (\ref{sdga-equation-pullback}) 与局部化函子
$K(\textit{Mod}(\mathcal{A}', \text{d})) \to D(\mathcal{A}', \text{d})$
的复合有左导出延拓
$D(\mathcal{B},\text{d})\to D(\mathcal{A}',\text{d})$；
它在良好右微分分次 $\mathcal{B}$-模 $\mathcal{P}$ 上的值为
$f^*\mathcal{P}\otimes_\mathcal{A}\mathcal{N}$。
\end{lemma}

\begin{proof}
回忆，对任意（右）微分分次 $\mathcal{B}$-模 $\mathcal{M}$，
存在拟同构 $\mathcal{P}\to\mathcal{M}$，其中 $\mathcal{P}$
是良好微分分次 $\mathcal{B}$-模；参见引理 \ref{sdga-lemma-resolve}。
因此，由《导出范畴》引理
\ref{derived-lemma-find-existence-computes}，只需证明：
给定良好微分分次 $\mathcal{B}$-模间的拟同构
$\mathcal{P}\to\mathcal{P}'$，诱导映射
$$
f^*\mathcal{P} \otimes_\mathcal{A} \mathcal{N}
\longrightarrow
f^*\mathcal{P}' \otimes_\mathcal{A} \mathcal{N}
$$
是拟同构。由引理 \ref{sdga-lemma-good-admissible-ses}，
$\mathcal{P}\to\mathcal{P}'$ 的锥 $\mathcal{P}''$
是良好微分分次 $\mathcal{A}$-模。由于在
$K(\textit{Mod}(\mathcal{B},\text{d}))$ 中有杰出三角
$$
\mathcal{P} \to \mathcal{P}' \to \mathcal{P}'' \to \mathcal{P}[1]
$$
，在 $K(\textit{Mod}(\mathcal{A}',\text{d}))$ 中得到杰出三角
$$
f^*\mathcal{P} \otimes_\mathcal{A} \mathcal{N} \to
f^*\mathcal{P}' \otimes_\mathcal{A} \mathcal{N} \to
f^*\mathcal{P}'' \otimes_\mathcal{A} \mathcal{N} \to
f^*\mathcal{P}[1] \otimes_\mathcal{A} \mathcal{N}
$$
。由引理 \ref{sdga-lemma-acyclic-good}，微分分次模
$f^*\mathcal{P}'' \otimes_\mathcal{A} \mathcal{N}$
无环，证明完成。
\end{proof}

\begin{definition}
\label{sdga-definition-pullback}
导出张量积与导出拉回。
\begin{enumerate}
\item 设 $(\mathcal{C},\mathcal{O})$ 是环化位点，
$\mathcal{A},\mathcal{B}$ 是微分分次 $\mathcal{O}$-代数，
$\mathcal{N}$ 是微分分次 $(\mathcal{A},\mathcal{B})$-双模。
引理 \ref{sdga-lemma-derived-tensor-product} 中构造的函子
$D(\mathcal{A},\text{d})\to D(\mathcal{B},\text{d})$
称为{\it 导出张量积}，记作
$-\otimes_\mathcal{A}^\mathbf{L}\mathcal{N}$。
\item 设 $(f, f^\sharp) : (\Sh(\mathcal{C}), \mathcal{O}_\mathcal{C})
\to (\Sh(\mathcal{D}), \mathcal{O}_\mathcal{D})$
是环化拓扑斯态射，$\mathcal{A}$ 是微分分次
$\mathcal{O}_\mathcal{C}$-代数，$\mathcal{B}$ 是微分分次
$\mathcal{O}_\mathcal{D}$-代数，并设
$\varphi:\mathcal{B}\to f_*\mathcal{A}$ 是微分分次
$\mathcal{O}_\mathcal{D}$-代数同态。引理
\ref{sdga-lemma-derived-tensor-product} 中构造的函子
$D(\mathcal{B},\text{d})\to D(\mathcal{A},\text{d})$
称为{\it 导出拉回}，记作 $Lf^*$。
\end{enumerate}
\end{definition}

\noindent
有了这些术语，可以表述一些显然的相容性。

\begin{lemma}
\label{sdga-lemma-compose-pullback-tensor}
引理 \ref{sdga-lemma-derived-tensor-product} 中的函子
$D(\mathcal{B},\text{d})\to D(\mathcal{A}',\text{d})$ 等于
$\mathcal{M} \mapsto
Lf^*\mathcal{M} \otimes_\mathcal{A}^\mathbf{L} \mathcal{N}$.
\end{lemma}

\begin{proof}
这是立即的，因为可用良好微分分次模表示对象来计算这些函子，
并且当 $\mathcal{P}$ 是良好微分分次 $\mathcal{B}$-模时，
$f^*\mathcal{P}$ 是良好微分分次 $\mathcal{A}$-模。
\end{proof}

\begin{lemma}
\label{sdga-lemma-compose-pullback}
设 $(f, f^\sharp) : (\Sh(\mathcal{C}), \mathcal{O})
\to (\Sh(\mathcal{C}'), \mathcal{O}')$ 与
$(g, g^\sharp) : (\Sh(\mathcal{C}'), \mathcal{O}')
\to (\Sh(\mathcal{C}''), \mathcal{O}'')$
是环化拓扑斯态射。设 $\mathcal{A},\mathcal{A}',\mathcal{A}''$
分别是微分分次 $\mathcal{O}$-代数、
$\mathcal{O}'$-代数、$\mathcal{O}''$-代数。设
$\varphi : \mathcal{A}' \to f_*\mathcal{A}$ 与
$\varphi' : \mathcal{A}'' \to g_*\mathcal{A}'$
分别是微分分次 $\mathcal{O}'$-代数同态与
$\mathcal{O}''$-代数同态。则
$L(g \circ f)^* = Lf^* \circ Lg^* :
D(\mathcal{A}'', \text{d}) \to D(\mathcal{A}, \text{d})$.
\end{lemma}

\begin{proof}
这是立即的，因为可用良好微分分次模表示对象来计算这些函子，
并且当 $\mathcal{P}$ 是良好微分分次 $\mathcal{A}$-模时，
$f^*\mathcal{P}$ 是良好微分分次 $\mathcal{A}'$-模。
\end{proof}

\noindent
设 $(\mathcal{C},\mathcal{O})$ 是环化位点，
$\mathcal{A},\mathcal{B}$ 是微分分次 $\mathcal{O}$-代数。
设 $\mathcal{N}\to\mathcal{N}'$ 是微分分次
$(\mathcal{A},\mathcal{B})$-双模同态。于是得到典范映射
$$
t :
\mathcal{M} \otimes_\mathcal{A}^\mathbf{L} \mathcal{N}
\longrightarrow
\mathcal{M} \otimes_\mathcal{A}^\mathbf{L} \mathcal{N}'
$$
，它关于 $D(\mathcal{A},\text{d})$ 中的 $\mathcal{M}$ 具有函子性，
并定义三角范畴间正合函子
$D(\mathcal{A},\text{d})\to D(\mathcal{B},\text{d})$
之间的自然变换。$t$ 在良好微分分次
$\mathcal{A}$-模 $\mathcal{P}$ 上的值是显然的映射
$$
\mathcal{P} \otimes_\mathcal{A}^\mathbf{L} \mathcal{N} =
\mathcal{P} \otimes_\mathcal{A} \mathcal{N}
\longrightarrow
\mathcal{P} \otimes_\mathcal{A} \mathcal{N}' =
\mathcal{P} \otimes_\mathcal{A}^\mathbf{L} \mathcal{N}'
$$

\begin{lemma}
\label{sdga-lemma-tensor-symmetry}
在上述情形下，若 $\mathcal{N}\to\mathcal{N}'$
在上同调层上是同构，则 $t$ 是函子之间的同构
$(- \otimes_\mathcal{A}^\mathbf{L} \mathcal{N}) \to
(- \otimes_\mathcal{A}^\mathbf{L} \mathcal{N}')$.
\end{lemma}

\begin{proof}
只需证明，对任意良好微分分次 $\mathcal{A}$-模
$\mathcal{P}$，映射
$\mathcal{P} \otimes_\mathcal{A} \mathcal{N} \to
\mathcal{P} \otimes_\mathcal{A} \mathcal{N}'$
在上同调层上是同构。为此，令 $\mathcal{N}''$
为把 $\mathcal{N}\to\mathcal{N}'$ 视为微分分次左
$\mathcal{A}$-模映射时的锥；参见定义 \ref{sdga-definition-cone}。
（当然，$\mathcal{N}''$ 也是双模，但我们不需要这一点。）
由张量积构造的函子性（它是微分分次范畴间的函子），
$\mathcal{P}\otimes_\mathcal{A}\mathcal{N}''$
作为 $\mathcal{O}$-模复形，是映射
$\mathcal{P} \otimes_\mathcal{A} \mathcal{N} \to
\mathcal{P} \otimes_\mathcal{A} \mathcal{N}'$.
的锥。因此只需证明
$\mathcal{P} \otimes_\mathcal{A} \mathcal{N}''$
无环。这是因为 $\mathcal{P}$ 良好，而 $\mathcal{N}''$
作为拟同构的锥是无环的。
\end{proof}

\begin{lemma}
\label{sdga-lemma-good-on-other-side}
设 $(\mathcal{C},\mathcal{O})$ 是环化位点，
$\mathcal{A},\mathcal{B}$ 是微分分次 $\mathcal{O}$-代数，
$\mathcal{N}$ 是微分分次 $(\mathcal{A},\mathcal{B})$-双模。
若 $\mathcal{N}$ 作为微分分次左 $\mathcal{A}$-模良好，
则对所有微分分次 $\mathcal{A}$-模 $\mathcal{M}$，有
$\mathcal{M}\otimes_\mathcal{A}^\mathbf{L}\mathcal{N}
=\mathcal{M}\otimes_\mathcal{A}\mathcal{N}$。
\end{lemma}

\begin{proof}
设 $\mathcal{P}\to\mathcal{M}$ 是拟同构，其中 $\mathcal{P}$
是良好（右）微分分次 $\mathcal{A}$-模。为证明本引理，
需要说明
$\mathcal{P} \otimes_\mathcal{A} \mathcal{N} \to
\mathcal{M} \otimes_\mathcal{A} \mathcal{N}$
是拟同构。映射 $\mathcal{P}\to\mathcal{M}$ 的锥 $C$
是无环右微分分次 $\mathcal{A}$-模。由于假设
$\mathcal{N}$ 作为微分分次左 $\mathcal{A}$-模良好，
$C\otimes_\mathcal{A}\mathcal{N}$ 无环。而作为
$\mathcal{O}$-模复形，$C\otimes_\mathcal{A}\mathcal{N}$
是映射
$\mathcal{P}\otimes_\mathcal{A}\mathcal{N}\to
\mathcal{M}\otimes_\mathcal{A}\mathcal{N}$ 的锥，故结论成立。
\end{proof}

\begin{lemma}
\label{sdga-lemma-compose-tensor}
设 $(\mathcal{C},\mathcal{O})$ 是环化位点，
$\mathcal{A},\mathcal{A}',\mathcal{A}''$ 是微分分次
$\mathcal{O}$-代数。设 $\mathcal{N}$ 与 $\mathcal{N}'$
分别是微分分次 $(\mathcal{A},\mathcal{A}')$-双模和
$(\mathcal{A}',\mathcal{A}'')$-双模。假设典范映射
$$
\mathcal{N} \otimes_{\mathcal{A}'}^\mathbf{L} \mathcal{N}'
\longrightarrow
\mathcal{N} \otimes_{\mathcal{A}'} \mathcal{N}'
$$
在 $D(\mathcal{A}'',\text{d})$ 中是拟同构。则
$$
(\mathcal{M}
\otimes_\mathcal{A}^\mathbf{L} \mathcal{N})
\otimes_{\mathcal{A}'}^\mathbf{L} \mathcal{N}'
=
\mathcal{M}
\otimes_\mathcal{A}^\mathbf{L}
(\mathcal{N} \otimes_{\mathcal{A}'} \mathcal{N}')
$$
是函子
$D(\mathcal{A},\text{d})\to D(\mathcal{A}'',\text{d})$ 的等式。
\end{lemma}

\begin{proof}
选取良好微分分次 $\mathcal{A}$-模 $\mathcal{P}$
以及拟同构 $\mathcal{P}\to\mathcal{M}$；参见引理
\ref{sdga-lemma-resolve}。则
$$
\mathcal{M}
\otimes_\mathcal{A}^\mathbf{L}
(\mathcal{N} \otimes_{\mathcal{A}'} \mathcal{N}') =
\mathcal{P} \otimes_\mathcal{A} \mathcal{N}
\otimes_{\mathcal{A}'} \mathcal{N}'
$$
并且
$$
(\mathcal{M}
\otimes_\mathcal{A}^\mathbf{L} \mathcal{N})
\otimes_{\mathcal{A}'}^\mathbf{L} \mathcal{N}' =
(\mathcal{P} \otimes_\mathcal{A} \mathcal{N})
\otimes_{\mathcal{A}'}^\mathbf{L} \mathcal{N}'
$$
因此需要证明典范映射
$$
(\mathcal{P} \otimes_\mathcal{A} \mathcal{N})
\otimes_{\mathcal{A}'}^\mathbf{L} \mathcal{N}'
\longrightarrow
\mathcal{P} \otimes_\mathcal{A} \mathcal{N}
\otimes_{\mathcal{A}'} \mathcal{N}'
$$
是拟同构。选取拟同构 $\mathcal{Q}\to\mathcal{N}'$，
其中 $\mathcal{Q}$ 是良好左微分分次 $\mathcal{A}'$-模
（引理 \ref{sdga-lemma-resolve}）。由引理
\ref{sdga-lemma-good-on-other-side}，上面的映射作为
$\mathcal{O}$-模导出范畴中的映射是
$$
\mathcal{P} \otimes_\mathcal{A} \mathcal{N}
\otimes_{\mathcal{A}'} \mathcal{Q}
\longrightarrow
\mathcal{P} \otimes_\mathcal{A} \mathcal{N}
\otimes_{\mathcal{A}'} \mathcal{N}'
$$
。依假设，
$\mathcal{N}\otimes_{\mathcal{A}'}\mathcal{Q}\to
\mathcal{N}\otimes_{\mathcal{A}'}\mathcal{N}'$ 是拟同构；
又因 $\mathcal{P}$ 是良好微分分次 $\mathcal{A}$-模，
由引理 \ref{sdga-lemma-tensor-symmetry}，此映射是拟同构
（左端与右端分别计算
$\mathcal{P}\otimes_\mathcal{A}^\mathbf{L}
(\mathcal{N}\otimes_{\mathcal{A}'}\mathcal{Q})$ 和
$\mathcal{P}\otimes_\mathcal{A}^\mathbf{L}
(\mathcal{N}\otimes_{\mathcal{A}'}\mathcal{N}')$；
也可以直接重复该引理证明中的论证）。
\end{proof}






\section{导出直像}
\label{sdga-section-derived-pushforward}

\noindent
足够多 K-内射对象的存在性保证：可以对同伦范畴上的任意正合函子
取右导出函子。

\begin{lemma}
\label{sdga-lemma-right-derived}
设 $(\mathcal{C},\mathcal{O})$ 是环化位点，
$(\mathcal{A},\text{d})$ 是 $(\mathcal{C},\mathcal{O})$
上的微分分次代数层。则三角范畴间的任意正合函子
$$
T : K(\textit{Mod}(\mathcal{A}, \text{d})) \longrightarrow \mathcal{D}
$$
都有右导出延拓
$RT:D(\mathcal{A},\text{d})\to\mathcal{D}$；
它在既分次内射又 K-内射的微分分次
$\mathcal{A}$-模 $\mathcal{I}$ 上的值为 $T(\mathcal{I})$。
\end{lemma}

\begin{proof}
由定理 \ref{sdga-theorem-qis-into-dg-injective}，对任意（右）微分分次
$\mathcal{A}$-模 $\mathcal{M}$，存在拟同构
$\mathcal{M}\to\mathcal{I}$，其中 $\mathcal{I}$
是既分次内射又 K-内射的微分分次 $\mathcal{A}$-模。
因此，由《导出范畴》引理
\ref{derived-lemma-find-existence-computes}，只需证明：
若 $\mathcal{I}\to\mathcal{I}'$ 是两个既分次内射又 K-内射的
微分分次 $\mathcal{A}$-模之间的拟同构，则
$T(\mathcal{I})\to T(\mathcal{I}')$ 是同构。这确实成立，
因为由引理 \ref{sdga-lemma-hom-derived} 可知
$\mathcal{I}\to\mathcal{I}'$ 在
$K(\textit{Mod}(\mathcal{A},\text{d}))$ 中是同构
（也可由引理 \ref{sdga-lemma-second-property-dg-injective} 推出）。
\end{proof}

\noindent
我们已经见过若干可应用此结果的函子。下面给出两个例子。

\begin{definition}
\label{sdga-definition-pushforward}
导出内 Hom 与导出直像。
\begin{enumerate}
\item 设 $(\mathcal{C},\mathcal{O})$ 是环化位点，
$\mathcal{A},\mathcal{B}$ 是微分分次 $\mathcal{O}$-代数，
$\mathcal{N}$ 是微分分次 $(\mathcal{A},\mathcal{B})$-双模。
内 Hom 函子 $\SheafHom_\mathcal{B}^{dg}(\mathcal{N},-)$
的右导出延拓
$$
R\SheafHom_\mathcal{B}(\mathcal{N}, -) :
D(\mathcal{B}, \text{d})
\longrightarrow
D(\mathcal{A}, \text{d})
$$
称为{\it 导出内 Hom}。
\item 设 $(f, f^\sharp) : (\Sh(\mathcal{C}), \mathcal{O}_\mathcal{C})
\to (\Sh(\mathcal{D}), \mathcal{O}_\mathcal{D})$
是环化拓扑斯态射，$\mathcal{A}$ 是微分分次
$\mathcal{O}_\mathcal{C}$-代数，$\mathcal{B}$ 是微分分次
$\mathcal{O}_\mathcal{D}$-代数，并设
$\varphi:\mathcal{B}\to f_*\mathcal{A}$ 是微分分次
$\mathcal{O}_\mathcal{D}$-代数同态。直像 $f_*$ 的右导出延拓
$$
Rf_* :
D(\mathcal{A}, \text{d})
\longrightarrow
D(\mathcal{B}, \text{d})
$$
称为{\it 导出直像}。
\end{enumerate}
\end{definition}

\noindent
事实证明，
$Rf_* : D(\mathcal{A}, \text{d}) \to D(\mathcal{B}, \text{d})$
与底层 $\mathcal{O}$-模复形上的导出直像一致；参见引理
\ref{sdga-lemma-pushforward-agrees}。

\medskip\noindent
这些函子分别是导出拉回与导出张量积的伴随。

\begin{lemma}
\label{sdga-lemma-derived-adjoint-tensor-hom}
设 $(\mathcal{C},\mathcal{O})$ 是环化位点，
$\mathcal{A},\mathcal{B}$ 是微分分次 $\mathcal{O}$-代数，
$\mathcal{N}$ 是微分分次 $(\mathcal{A},\mathcal{B})$-双模。则
$$
R\SheafHom_\mathcal{B}(\mathcal{N}, -) :
D(\mathcal{B}, \text{d})
\longrightarrow
D(\mathcal{A}, \text{d})
$$
是
$$
- \otimes_\mathcal{A}^\mathbf{L} \mathcal{N} :
D(\mathcal{A}, \text{d})
\longrightarrow
D(\mathcal{B}, \text{d})
$$
的右伴随。
\end{lemma}

\begin{proof}
这由《导出范畴》引理
\ref{derived-lemma-pre-derived-adjoint-functors-general}
和引理 \ref{sdga-lemma-tensor-hom-adjunction-dg} 得出。
\end{proof}

\begin{lemma}
\label{sdga-lemma-derived-adjoint-push-pull}
设 $(f, f^\sharp) : (\Sh(\mathcal{C}), \mathcal{O}_\mathcal{C})
\to (\Sh(\mathcal{D}), \mathcal{O}_\mathcal{D})$
是环化拓扑斯态射，$\mathcal{A}$ 是微分分次
$\mathcal{O}_\mathcal{C}$-代数，$\mathcal{B}$ 是微分分次
$\mathcal{O}_\mathcal{D}$-代数，并设
$\varphi:\mathcal{B}\to f_*\mathcal{A}$ 是微分分次
$\mathcal{O}_\mathcal{D}$-代数同态。则
$$
Rf_* : 
D(\mathcal{A}, \text{d})
\longrightarrow
D(\mathcal{B}, \text{d})
$$
是
$$
Lf^* :
D(\mathcal{B}, \text{d})
\longrightarrow
D(\mathcal{A}, \text{d})
$$
的右伴随。
\end{lemma}

\begin{proof}
这由《导出范畴》引理
\ref{derived-lemma-pre-derived-adjoint-functors-general}
和引理 \ref{sdga-lemma-adjunction-push-pull-dg} 得出。
\end{proof}

\noindent
下面讨论第 \ref{sdga-section-derived-pullback} 节所考虑的情形。

\medskip\noindent
设 $(f, f^\sharp) : (\Sh(\mathcal{C}), \mathcal{O}_\mathcal{C})
\to (\Sh(\mathcal{D}), \mathcal{O}_\mathcal{D})$
是环化拓扑斯态射。设 $\mathcal{A}$ 是微分分次
$\mathcal{O}_\mathcal{C}$-代数，$\mathcal{B}$ 是微分分次
$\mathcal{O}_\mathcal{D}$-代数。设给定微分分次
$f^{-1}\mathcal{O}_\mathcal{D}$-代数映射
$$
\varphi : f^{-1}\mathcal{B} \to \mathcal{A}
$$
。由标量限制与标量扩张的伴随性，这等价于微分分次
$\mathcal{O}_\mathcal{C}$-代数映射
$\varphi:f^*\mathcal{B}\to\mathcal{A}$；等价地，
$\varphi$ 可视为微分分次 $\mathcal{O}_\mathcal{D}$-代数映射
$$
\varphi : \mathcal{B} \to f_*\mathcal{A}
$$
。参见评注 \ref{sdga-remark-functoriality-dga}。

\medskip\noindent
除上述数据外，再设 $\mathcal{A}'$ 是另一个微分分次
$\mathcal{O}_\mathcal{C}$-代数，$\mathcal{N}$ 是微分分次
$(\mathcal{A},\mathcal{A}')$-双模。在此情形下可考虑函子
$$
\textit{Mod}(\mathcal{A}', \text{d})
\longrightarrow
\textit{Mod}(\mathcal{B}, \text{d}),\quad
\mathcal{M} \longmapsto
f_*\SheafHom^{dg}_{\mathcal{A}'}(\mathcal{N}, \mathcal{M})
$$
注意，由第 \ref{sdga-section-functoriality-dg} 与
\ref{sdga-section-dg-bimodules} 节的讨论，它延拓为微分分次范畴间的函子
$$
\textit{Mod}^{dg}(\mathcal{A}', \text{d})
\longrightarrow
\textit{Mod}^{dg}(\mathcal{B}, \text{d}),\quad
\mathcal{M} \longmapsto
f_*\SheafHom^{dg}_{\mathcal{A}'}(\mathcal{N}, \mathcal{M})
$$
。形式地随之得到三角范畴间的正合函子
\begin{equation}
\label{sdga-equation-pushforward}
K(\textit{Mod}(\mathcal{A}', \text{d}))
\longrightarrow
K(\textit{Mod}(\mathcal{B}, \text{d})),\quad
\mathcal{M} \longmapsto
f_*\SheafHom^{dg}_{\mathcal{A}'}(\mathcal{N}, \mathcal{M})
\end{equation}
。

\begin{lemma}
\label{sdga-lemma-compose-pushforward-hom}
在上述情形下，以
$RT:D(\mathcal{A}',\text{d})\to D(\mathcal{B},\text{d})$
表示 (\ref{sdga-equation-pushforward}) 的右导出延拓。则
$$
RT(\mathcal{M}) = Rf_* R\SheafHom(\mathcal{N}, \mathcal{M})
$$
关于 $\mathcal{M}$ 函子性地成立。
\end{lemma}

\begin{proof}
由引理 \ref{sdga-lemma-tensor-hom-adjunction-dg} 与
\ref{sdga-lemma-adjunction-push-pull-dg}，函子
(\ref{sdga-equation-pushforward}) 是函子
(\ref{sdga-equation-pullback}) 的右伴随。由《导出范畴》引理
\ref{derived-lemma-pre-derived-adjoint-functors-general}，
函子 $RT$ 是引理 \ref{sdga-lemma-derived-tensor-product}
中函子的右伴随；由引理 \ref{sdga-lemma-compose-pullback-tensor}，
后者等于
$Lf^*(-)\otimes_\mathcal{A}^\mathbf{L}\mathcal{N}$。
由引理 \ref{sdga-lemma-derived-adjoint-tensor-hom} 与
\ref{sdga-lemma-derived-adjoint-push-pull}，函子
$Lf^*(-)\otimes_\mathcal{A}^\mathbf{L}\mathcal{N}$ 的右伴随为
$Rf_* R\SheafHom(\mathcal{N}, -)$
。故由伴随的唯一性得出结论。
\end{proof}

\begin{lemma}
\label{sdga-lemma-compose-pushforward}
设 $(f, f^\sharp) : (\Sh(\mathcal{C}), \mathcal{O})
\to (\Sh(\mathcal{C}'), \mathcal{O}')$ 与
$(g, g^\sharp) : (\Sh(\mathcal{C}'), \mathcal{O}')
\to (\Sh(\mathcal{C}''), \mathcal{O}'')$
是环化拓扑斯态射。设 $\mathcal{A},\mathcal{A}',\mathcal{A}''$
分别是微分分次 $\mathcal{O}$-代数、
$\mathcal{O}'$-代数、$\mathcal{O}''$-代数，并设
$\varphi : \mathcal{A}' \to f_*\mathcal{A}$ 与
$\varphi' : \mathcal{A}'' \to g_*\mathcal{A}'$
分别是微分分次 $\mathcal{O}'$-代数同态与
$\mathcal{O}''$-代数同态。则
$R(g \circ f)_* = Rg_* \circ Rf_* :
D(\mathcal{A}, \text{d}) \to D(\mathcal{A}'', \text{d})$.
\end{lemma}

\begin{proof}
由引理 \ref{sdga-lemma-compose-pullback}、
\ref{sdga-lemma-derived-adjoint-push-pull} 及伴随的唯一性得出。
\end{proof}

\begin{lemma}
\label{sdga-lemma-compose-hom}
设 $(\mathcal{C},\mathcal{O})$ 是环化位点，
$\mathcal{A},\mathcal{A}',\mathcal{A}''$ 是微分分次
$\mathcal{O}$-代数。设 $\mathcal{N}$ 与 $\mathcal{N}'$
分别是微分分次 $(\mathcal{A},\mathcal{A}')$-双模和
$(\mathcal{A}',\mathcal{A}'')$-双模。假设典范映射
$$
\mathcal{N} \otimes_{\mathcal{A}'}^\mathbf{L} \mathcal{N}'
\longrightarrow
\mathcal{N} \otimes_{\mathcal{A}'} \mathcal{N}'
$$
在 $D(\mathcal{A}'',\text{d})$ 中是拟同构。则
$$
R\SheafHom_{\mathcal{A}''}
(\mathcal{N} \otimes_{\mathcal{A}'} \mathcal{N}', -)
=
R\SheafHom_{\mathcal{A}'}(\mathcal{N},
R\SheafHom_{\mathcal{A}''}(\mathcal{N}', -))
$$
是函子
$D(\mathcal{A}'',\text{d})\to D(\mathcal{A},\text{d})$ 的等式。
\end{lemma}

\begin{proof}
由引理 \ref{sdga-lemma-compose-tensor}、
\ref{sdga-lemma-derived-adjoint-tensor-hom} 及伴随的唯一性得出。
\end{proof}

\begin{lemma}
\label{sdga-lemma-pushforward-agrees}
设 $(f, f^\sharp) : (\Sh(\mathcal{C}), \mathcal{O}_\mathcal{C})
\to (\Sh(\mathcal{D}), \mathcal{O}_\mathcal{D})$
是环化拓扑斯态射。设 $\mathcal{A}$ 是微分分次
$\mathcal{O}_\mathcal{C}$-代数，$\mathcal{B}$ 是微分分次
$\mathcal{O}_\mathcal{D}$-代数，并设
$\varphi:\mathcal{B}\to f_*\mathcal{A}$ 是微分分次
$\mathcal{O}_\mathcal{D}$-代数同态。下图交换：
$$
\xymatrix{
D(\mathcal{A}, \text{d}) \ar[d]_{Rf_*} \ar[rr]_{forget} & &
D(\mathcal{O}_\mathcal{C}) \ar[d]^{Rf_*} \\
D(\mathcal{B}, \text{d}) \ar[rr]^{forget} & &
D(\mathcal{O}_\mathcal{D})
}
$$
\end{lemma}

\begin{proof}
除去若干范畴的等同，本引理由引理
\ref{sdga-lemma-compose-pushforward} 立即得出。

\medskip\noindent
把 $\mathcal{O}_\mathcal{C}$ 置于次数 $0$ 并赋予零微分，
便可将 $(\mathcal{O}_\mathcal{C},0)$ 视为微分分次
$\mathcal{O}_\mathcal{C}$-代数。显然有
$$
\textit{Mod}(\mathcal{O}_\mathcal{C}, 0) =
\text{Comp}(\mathcal{O}_\mathcal{C})
\quad\text{且}\quad
D(\mathcal{O}_\mathcal{C}, 0) = D(\mathcal{O}_\mathcal{C})
$$
在此等同下，遗忘函子
$\textit{Mod}(\mathcal{A}, \text{d}) \to
\text{Comp}(\mathcal{O}_\mathcal{C})$
是第 \ref{sdga-section-functoriality-dg} 节定义的``直像''
$\text{id}_{\mathcal{C},*}$；它对应于环化拓扑斯的恒等态射
$\text{id}_\mathcal{C} : (\mathcal{C}, \mathcal{O}_\mathcal{C}) \to
(\mathcal{C}, \mathcal{O}_\mathcal{C})$，以及微分分次
$\mathcal{O}_\mathcal{C}$-代数映射
$(\mathcal{O}_\mathcal{C},0)\to(\mathcal{A},\text{d})$。
由于 $\text{id}_{\mathcal{C},*}$ 正合，立即可知
$$
R\text{id}_{\mathcal{C}, *} = forget :
D(\mathcal{A}, \text{d}) \longrightarrow
D(\mathcal{O}_\mathcal{C}, 0) = D(\mathcal{O}_\mathcal{C})
$$
完全相同的论证表明
$$
R\text{id}_{\mathcal{D}, *} = forget :
D(\mathcal{B}, \text{d}) \longrightarrow
D(\mathcal{O}_\mathcal{D}, 0) = D(\mathcal{O}_\mathcal{D})
$$
此外，《位点上的上同调》第
\ref{sites-cohomology-section-unbounded} 节中
$Rf_* : D(\mathcal{O}_\mathcal{C}) \to D(\mathcal{O}_\mathcal{D})$
的构造，与定义 \ref{sdga-definition-pushforward} 中
$Rf_* : D(\mathcal{O}_\mathcal{C}, 0) \to D(\mathcal{O}_\mathcal{D}, 0)$
的构造一致，因为两个函子都定义为底层模复形之直像的右导出延拓。
由引理 \ref{sdga-lemma-compose-pushforward}，
$Rf_*\circ R\text{id}_{\mathcal{C},*}$ 与
$R\text{id}_{\mathcal{D},*}\circ Rf_*$ 都是
$f_*\circ forget=forget\circ f_*$ 的导出函子，
故由伴随的唯一性二者相等。
\end{proof}

\begin{lemma}
\label{sdga-lemma-cohomology-ext}
设 $(\mathcal{C},\mathcal{O})$ 是环化位点，
$\mathcal{A}$ 是微分分次 $\mathcal{O}$-代数，
$\mathcal{M}$ 是微分分次 $\mathcal{A}$-模，且
$n\in\mathbf{Z}$。则
$$
H^n(\mathcal{C}, \mathcal{M}) =
\Hom_{D(\mathcal{A}, \text{d})}(\mathcal{A}, \mathcal{M}[n])
$$
，其中左端是把 $\mathcal{M}$ 视为 $\mathcal{O}$-模复形所得的上同调。
\end{lemma}

\begin{proof}
为证明该公式，注意
$$
R\Gamma(\mathcal{C}, \mathcal{M}) = \Gamma(\mathcal{C}, \mathcal{I})
$$
，其中 $\mathcal{M}\to\mathcal{I}$ 是到既分次内射又 K-内射的
微分分次 $\mathcal{A}$-模 $\mathcal{I}$ 的拟同构
（合用引理 \ref{sdga-lemma-right-derived} 与
\ref{sdga-lemma-pushforward-agrees}）。由引理
\ref{sdga-lemma-hom-derived}，有
$$
\Hom_{D(\mathcal{A}, \text{d})}(\mathcal{A}, \mathcal{M}[n]) =
\Hom_{K(\textit{Mod}(\mathcal{A}, \text{d}))}(\mathcal{M}, \mathcal{I}[n]) =
H^0(\Gamma(\mathcal{C}, \mathcal{I}[n])) =
H^n(\Gamma(\mathcal{C}, \mathcal{I}))
$$
合并这两个结果即得所需等式。
\end{proof}








\section{导出范畴的等价}
\label{sdga-section-equivalence}

\noindent
本节是《微分分次代数》第 \ref{dga-section-equivalence} 节的类似物。

\begin{lemma}
\label{sdga-lemma-qis-equivalence}
设 $(\mathcal{C},\mathcal{O})$ 是环化位点。若
$\varphi:\mathcal{A}\to\mathcal{B}$ 是微分分次
$\mathcal{O}$-代数的同态，并在上同调层上诱导同构，则
$$
D(\mathcal{A}, \text{d}) \longrightarrow D(\mathcal{B}, \text{d}), \quad
\mathcal{M}
\longmapsto
\mathcal{M} \otimes_\mathcal{A}^\mathbf{L} \mathcal{B}
$$
是范畴等价。
\end{lemma}

\begin{proof}
回忆限制函子
$$
\textit{Mod}^{dg}(\mathcal{B}, \text{d}) \to
\textit{Mod}^{dg}(\mathcal{A}, \text{d}),\quad
\mathcal{N} \mapsto res_\varphi \mathcal{N}
$$
是函子
$$
\textit{Mod}^{dg}(\mathcal{A}, \text{d}) \to
\textit{Mod}^{dg}(\mathcal{B}, \text{d}),\quad
\mathcal{M} \mapsto \mathcal{M} \otimes_\mathcal{A} \mathcal{B}
$$
的右伴随；见第 \ref{sdga-section-dg-bimodules} 节。由于限制把拟同构
映为拟同构，故它显然有左导出延拓（仍由限制给出）。由
《导出范畴》引理
\ref{derived-lemma-pre-derived-adjoint-functors-general}，此函子是
$- \otimes_\mathcal{A}^\mathbf{L} \mathcal{B}$ 的右伴随。伴随态射
$$
\mathcal{M} \to
res_\varphi(\mathcal{M} \otimes_\mathcal{A}^\mathbf{L} \mathcal{B})
$$
在 $D(\mathcal{A}, \text{d})$ 中是同构，因为按照假设，
$\mathcal{A}\to\mathcal{B}$ 是（左）微分分次
$\mathcal{A}$-模的拟同构。特别地，本引理中的函子全忠实；见
《范畴》引理 \ref{categories-lemma-adjoint-fully-faithful}。
显然，限制函子
$D(\mathcal{B}, \text{d}) \to D(\mathcal{A}, \text{d})$
的核为零。于是结论由《导出范畴》引理
\ref{derived-lemma-fully-faithful-adjoint-kernel-zero} 得出。
\end{proof}













\section{微分分次代数的分辨}
\label{sdga-section-resolution-dgas}

\noindent
本节是《微分分次代数》第 \ref{dga-section-resolution-dgas} 节的类似物。

\medskip\noindent
设 $(\mathcal{C},\mathcal{O})$ 是环化位点。与评注
\ref{sdga-remark-sheaf-graded-sets} 一样，考虑 $\mathcal{C}$ 上的分次集合层
$\mathcal{S}$。把 $r$ 重自积
$\mathcal{S} \times \ldots \times \mathcal{S}$
看作按如下规则分次的集合层：
$\deg(s_1 \cdot \ldots \cdot s_r) = \sum \deg(s_i)$。
这里，给定局部截面 $s_i\in\mathcal{S}(U)$，$i=1,\ldots,r$，
用 $s_1\cdot\ldots\cdot s_r$ 表示
$\mathcal{S} \times \ldots \times \mathcal{S}$
在 $U$ 上的相应截面。记 $\mathcal{O}\langle\mathcal{S}\rangle$
为 $\mathcal{S}$ 上的自由分次 $\mathcal{O}$-代数。更准确地，置
$$
\mathcal{O}\langle \mathcal{S} \rangle =
\mathcal{O} \oplus
\bigoplus\nolimits_{r \geq 1}
\mathcal{O}[\mathcal{S} \times \ldots \times \mathcal{S}]
$$
，记号同评注 \ref{sdga-remark-sheaf-graded-sets}。按串接定义乘法：
$$
(s_1 \cdot \ldots \cdot s_r)  (s'_1 \cdot \ldots \cdot s'_{r'}) =
s_1 \cdot \ldots s_r \cdot s'_1 \cdot \ldots \cdot s'_{r'}
$$
，便使其成为分次 $\mathcal{O}$-代数层。可以令
$\mathcal{S}$ 的所有局部截面 $s$ 满足 $\text{d}(s)=0$，
再依 Leibniz 法则唯一延拓，从而在
$\mathcal{O}\langle\mathcal{S}\rangle$ 上赋予微分；不过考虑其他微分
同样重要。

\medskip\noindent
事实上，设给定如下系统：
\begin{enumerate}
\item 对 $i=0,1,2,\ldots$，给定分次集合层 $\mathcal{S}_i$；
\item 对 $i=0,1,2,\ldots$，给定态射
$$
\delta_{i + 1} : \mathcal{S}_{i + 1}
\longrightarrow
\mathcal{A}_i =
\mathcal{O}\langle \mathcal{S}_0 \amalg \ldots \amalg \mathcal{S}_i\rangle
$$
，它是次数为 $1$ 的分次集合层态射，且其像包含于靶上归纳定义的微分之核中。
\end{enumerate}
更准确地，先置
$\mathcal{A}_0 = \mathcal{O}\langle \mathcal{S}_0 \rangle$
，并在其上赋予满足 Leibniz 法则且使 $\mathcal{S}$ 的任意局部截面
$s$ 满足 $\text{d}(s)=0$ 的唯一微分。归纳地，设
$\mathcal{A}_i$ 上已给定微分 $\text{d}$。将它延拓为
$\mathcal{A}_{i+1}$ 上满足 Leibniz 法则并使
$$
\text{d}(s) = \delta(s)
$$
的唯一微分，其中若 $s$ 属于
$\mathcal{S}_0\amalg\ldots\amalg\mathcal{S}_{i+1}$ 的直和项
$\mathcal{S}_j$，则 $\delta(s)=\delta_j(s)$。此定义恰因
$\delta(s)$ 属于归纳定义的微分之核而有意义。

\begin{lemma}
\label{sdga-lemma-special-good}
在上述情形中，微分分次 $\mathcal{O}$-代数
$$
\mathcal{A} = \colim \mathcal{A}_i
$$
具有如下性质：对环化拓扑斯的任意态射
$(f, f^\sharp) : (\Sh(\mathcal{C}'), \mathcal{O}')
\to (\Sh(\mathcal{C}), \mathcal{O})$
，拉回 $f^*\mathcal{A}$ 作为分次 $\mathcal{O}'$-模是平坦的，
并且作为 $\mathcal{O}'$-模复形是 K-平坦的。
\end{lemma}

\begin{proof}
注意 $f^*\mathcal{A}=\colim f^*\mathcal{A}_i$，且
$$
f^*\mathcal{A}_i = \mathcal{O}'\langle
f^{-1}\mathcal{S}_0 \amalg \ldots \amalg f^{-1}\mathcal{S}_i\rangle
$$
，其微分由上述归纳过程以 $f^{-1}\delta_{i+1}$ 给出。因此，只需证明
$\mathcal{A}$ 作为分次 $\mathcal{O}$-模平坦，并且作为
$\mathcal{O}$-模复形 K-平坦。为此，只需证明每个
$\mathcal{A}_i$ 作为分次 $\mathcal{O}$-模平坦，并且作为
$\mathcal{O}$-模复形 K-平坦；可比较引理
\ref{sdga-lemma-good-direct-sum}。


\medskip\noindent
对 $i\geq1$，记
$\mathcal{S}=\mathcal{S}_0\amalg\ldots\amalg\mathcal{S}_i$，
于是作为分次 $\mathcal{O}$-代数，有
$\mathcal{A}_i=\mathcal{O}\langle\mathcal{S}\rangle$。我们要用
微分分次 $\mathcal{O}$-子模构造该代数的一个滤过。

\medskip\noindent
置 $W=\mathbf{Z}_{\geq0}^{i+1}$，并赋予字典序。也就是说，
对 $W$ 中的 $w=(w_0,\ldots w_i)$ 与
$w'=(w'_0,\ldots,w'_i)$，规定
$$
w > w' \Leftrightarrow
\exists j,\ 0 \leq j \leq i :
w_i = w'_i,\ w_{i - 1} = w'_{i - 1},\ \ldots ,
\ w_{j + 1} = w'_{j + 1},\ w_j > w'_j
$$
，依此类推。设给定
$s = s_1 \cdot \ldots \cdot s_r$，它是
$\mathcal{S}\times\ldots\times\mathcal{S}$ 在 $U$ 上的截面。若对每个
$a$，存在唯一的 $0\leq j_a\leq i$ 使
$s_a\in\mathcal{S}_{j_a}(U)$，便称{\it $s$ 的权已定义}。此时定义其权为
$$
w(s) = (w_0(s), \ldots, w_i(s)) \in W,\quad
w_j(s) = |\{a \mid j_a = j\}|
$$
$\mathcal{S}\times\ldots\times\mathcal{S}$ 的任意截面之权在局部均有定义。
读者容易验证，我们得到不交并分解
$$
\mathcal{S} \times \ldots \times \mathcal{S} =
\coprod\nolimits_{w \in W} \left(
\mathcal{S} \times \ldots \times \mathcal{S}\right)_w
$$
，其中各子层由给定权的截面组成。当然，对给定的 $r$，只会出现满足
$\sum_{0\leq j\leq i}w_j=r$ 的 $w\in W$。相应地得到分解
$$
\mathcal{A}_i = \mathcal{O} \oplus
\bigoplus\nolimits_{r \geq 1}
\bigoplus\nolimits_{w \in W}
\mathcal{O}[\left(\mathcal{S} \times \ldots \times \mathcal{S}\right)_w]
$$
证明的其余部分依赖如下显然观察：给定 $r$、$w$ 以及
$\left(\mathcal{S}\times\ldots\times\mathcal{S}\right)_w$ 的局部截面
$s=s_1\cdot\ldots\cdot s_r$，有
$$
\text{d}(s) \text{ 是如下对象的局部截面：} \mathcal{O} \oplus
\bigoplus\nolimits_{r' \geq 1}
\bigoplus\nolimits_{w' \in W,\ w' < w}
\mathcal{O}[\left(\mathcal{S} \times \ldots \times \mathcal{S}\right)_{w'}]
$$
这是因为在下列各表达式中
$$
(-1)^{\deg(s_1) + \ldots + \deg(s_{a - 1})}
s_1 \cdot \ldots s_{a - 1} \cdot \delta(s_a) \cdot
s_{a + 1} \cdot \ldots \cdot s_r
$$
（它们之和给出 $\text{d}(s)$），元素 $\delta(s_a)$ 在局部是诸元素
$s'_1\cdot\ldots\cdot s'_{r'}$ 的 $\mathcal{O}$-线性组合，其中
$s'_{a'}$ 属于 $\mathcal{S}_{j'_a}$，对某个
$0\leq j'_{a'}<j_a$；这里 $j_a$ 由 $s_a$ 是
$\mathcal{S}_{j_a}$ 的截面所确定。

\medskip\noindent
这意味着如下事实。对 $w\in W$，置
$$
F_w \mathcal{A}_i = \mathcal{O} \oplus \bigoplus\nolimits_{r \geq 1}
\bigoplus\nolimits_{w' \in W,\ w' \leq w}
\mathcal{O}[\left(\mathcal{S} \times \ldots \times \mathcal{S}\right)_{w'}]
$$
由上述观察，这是一个微分分次 $\mathcal{O}$-子模。得到微分分次
$\mathcal{A}$-模的容许短正合列
$$
0 \to \colim_{w' < w} F_{w'}\mathcal{A}_i \to
F_w\mathcal{A}_i \to
\bigoplus\nolimits_{r \geq 1}
\mathcal{O}[\left(\mathcal{S} \times \ldots \times \mathcal{S}\right)_w]
\to 0
$$
，其中右端的微分为零。

\medskip\noindent
现在对有序集 $W$ 作超限归纳以完成证明。微分分次复形
$F_0\mathcal{A}_0$ 是直和项 $\mathcal{O}$，因而 K-平坦且分次平坦。
对 $w\in W$，若结论对所有 $w'<w$ 的
$F_{w'}\mathcal{A}_i$ 成立，则由引理
\ref{sdga-lemma-good-direct-sum}、\ref{sdga-lemma-good-admissible-ses} 与
\ref{sdga-lemma-free-graded-module-good}，结论对 $w$ 也成立。最后，
$\mathcal{A}_i=\colim_{w\in W}F_w\mathcal{A}_i$，故得结论。
\end{proof}

\begin{lemma}
\label{sdga-lemma-good-dga}
设 $(\mathcal{C},\mathcal{O})$ 是环化位点，
$(\mathcal{B},\text{d})$ 是微分分次 $\mathcal{O}$-代数。存在微分分次
$\mathcal{O}$-代数的拟同构
$(\mathcal{A}, \text{d}) \to (\mathcal{B}, \text{d})$，使得
$\mathcal{A}$ 分次平坦，并且作为 $\mathcal{O}$-模复形
K-平坦；而且经环化拓扑斯的任意态射拉回后，同样的结论仍成立。
\end{lemma}

\begin{proof}
证明与引理 \ref{sdga-lemma-resolve} 的第一个证明完全相同，只是现在使用
自由分次代数而非自由分次模。

\medskip\noindent
我们像引理 \ref{sdga-lemma-special-good} 那样，通过构造
$$
\mathcal{A}_0 \to \mathcal{A}_1 \to \mathcal{A}_2 \to \ldots \to \mathcal{B}
$$
来构造 $\mathcal{A}=\colim\mathcal{A}_i$。令 $\mathcal{S}_0$ 为
次数 $n$ 部分是 $\Ker(\text{d}_\mathcal{B}^n)$ 的分次集合层
（评注 \ref{sdga-remark-sheaf-graded-sets}）。考虑微分分次模的同态
$$
\mathcal{A}_0 = \mathcal{O}\langle \mathcal{S}_0 \rangle
\longrightarrow
\mathcal{B}
$$
，该映射把 $\mathcal{S}_0$ 的局部截面 $s$ 映到
$\mathcal{A}^{\deg(s)}$ 中相应的局部截面（它属于微分之核，故我们的映射
确实是微分分次代数的映射）。按构造，所诱导的上同调层映射
$H^n(\mathcal{A}_0)\to H^n(\mathcal{B})$ 是满射，因而对所有 $i$
仍有同样结论。

\medskip\noindent
构造的归纳步骤如下。给定 $\mathcal{A}_i\to\mathcal{B}$，记
$\mathcal{S}_{i+1}$ 为次数 $n$ 部分是
$$
\Ker(\text{d}_{\mathcal{A}_i}^{n + 1})
\times_{\mathcal{B}^{n + 1}, \text{d}}
\mathcal{B}^n
$$
的分次集合层。它带有典范映射
$$
\delta_{i + 1} : \mathcal{S}_{i + 1} \longrightarrow
\mathcal{A}_i
$$
，按构造其像包含于 $\text{d}_{\mathcal{A}_i}$ 之核。因此
$\mathcal{A}_{i + 1}=\mathcal{O}\langle
\mathcal{S}_0\amalg\ldots\mathcal{S}_{i + 1}\rangle$ 带有延拓
$\mathcal{A}_i$ 上微分的微分；见本节开头的讨论。从
$\mathcal{A}_{i+1}$ 到 $\mathcal{B}$ 的映射是唯一的分次代数映射，
它在 $\mathcal{A}_i$ 上限制为给定映射，并把
$\mathcal{S}_{i+1}$ 的局部截面 $s=(a,b)$ 映到
$\mathcal{B}$ 中的 $b$。这与微分相容，恰因为 $\text{d}(b)$ 是
$a$ 在 $\mathcal{B}$ 中的像。

\medskip\noindent
映射 $\mathcal{A}\to\mathcal{B}$ 是拟同构：有
$H^n(\mathcal{A})=\colim H^n(\mathcal{A}_i)$，并且对每个 $i$，映射
$H^n(\mathcal{A}_i)\to H^n(\mathcal{B})$ 是满射，而按构造其核被映射
$H^n(\mathcal{A}_i)\to H^n(\mathcal{A}_{i+1})$ 零化。最后，
$\mathcal{A}$ 的平坦性条件已在引理 \ref{sdga-lemma-special-good} 中证明。
\end{proof}






\section{杂项}
\label{sdga-section-misc}

\noindent
设 $(f,f^\sharp):(\Sh(\mathcal{C}),\mathcal{O})
\to(\Sh(\mathcal{C}'),\mathcal{O}')$ 是环化拓扑斯的态射，且
$\mathcal{A}$ 是微分分次 $\mathcal{O}$-代数层。利用复合\footnote{更准确地，
应写成
$F(\mathcal{A}) \otimes_\mathcal{O}^\mathbf{L} F(\mathcal{A})
\to F(\mathcal{A} \otimes_\mathcal{O} \mathcal{A}) \to F(\mathcal{A})$
，其中 $F$ 表示到 $\mathcal{O}$-模复形的遗忘函子。还要注意，
$\mathcal{A}\otimes_\mathcal{O}\mathcal{A}$ 表示第
\ref{sdga-section-tensor-product-dg} 节的张量积，故
$F(\mathcal{A} \otimes_\mathcal{O} \mathcal{A}) =
\text{Tot}(F(\mathcal{A}) \otimes_\mathcal{O} F(\mathcal{A}))$.
该序列的第一个箭头是从两个 $\mathcal{O}$-模复形的导出张量积到
$\mathcal{O}$-模复形通常张量积的典范映射。}
$$
\mathcal{A} \otimes_\mathcal{O}^\mathbf{L} \mathcal{A}
\longrightarrow
\mathcal{A} \otimes_\mathcal{O} \mathcal{A}
\longrightarrow
\mathcal{A}
$$
以及相对杯积（见《站点上的上同调》评注
\ref{sites-cohomology-remark-cup-product} 与第
\ref{sites-cohomology-section-cup-product} 节），得到乘法\footnote{这里及下文，
$Rf_*:D(\mathcal{O})\to D(\mathcal{O}')$ 是《站点上的上同调》第
\ref{sites-cohomology-section-unbounded} 节及以后研究的导出函子。}
$$
\mu :
Rf_*\mathcal{A} \otimes_{\mathcal{O}'}^\mathbf{L} Rf_*\mathcal{A}
\longrightarrow
Rf_*\mathcal{A}
$$
，它位于 $D(\mathcal{O}')$ 中。此乘法具有结合性，即图
$$
\xymatrix{
Rf_*\mathcal{A} \otimes_{\mathcal{O}'}^\mathbf{L} Rf_*\mathcal{A}
\otimes_{\mathcal{O}'}^\mathbf{L} Rf_*\mathcal{A}
\ar[rr]_-{\mu \otimes 1} \ar[d]_{1 \otimes \mu} & &
Rf_*\mathcal{A} \otimes_{\mathcal{O}'}^\mathbf{L} Rf_*\mathcal{A}
\ar[d]^\mu \\
Rf_*\mathcal{A} \otimes_{\mathcal{O}'}^\mathbf{L} Rf_*\mathcal{A}
\ar[rr]^-\mu & &
Rf_*\mathcal{A}
}
$$
在 $D(\mathcal{O}')$ 中交换；这由《站点上的上同调》引理
\ref{sites-cohomology-lemma-cup-product-associative} 得出。
完全类似地，给定右微分分次 $\mathcal{A}$-模 $\mathcal{M}$，得到乘法
$$
\mu_\mathcal{M} :
Rf_*\mathcal{M} \otimes_{\mathcal{O}'}^\mathbf{L} Rf_*\mathcal{A}
\longrightarrow
Rf_*\mathcal{M}
$$
，它位于 $D(\mathcal{O}')$ 中。此乘法与上述 $\mu$ 相容，即图
$$
\xymatrix{
Rf_*\mathcal{M} \otimes_{\mathcal{O}'}^\mathbf{L} Rf_*\mathcal{A}
\otimes_{\mathcal{O}'}^\mathbf{L} Rf_*\mathcal{A}
\ar[rr]_-{\mu_\mathcal{M} \otimes 1} \ar[d]_{1 \otimes \mu} & &
Rf_*\mathcal{M} \otimes_{\mathcal{O}'}^\mathbf{L} Rf_*\mathcal{A}
\ar[d]^{\mu_\mathcal{M}} \\
Rf_*\mathcal{M} \otimes_{\mathcal{O}'}^\mathbf{L} Rf_*\mathcal{A}
\ar[rr]^-{\mu_\mathcal{M}} & &
Rf_*\mathcal{M}
}
$$
在 $D(\mathcal{O}')$ 中交换；这仍由《站点上的上同调》引理
\ref{sites-cohomology-lemma-cup-product-associative} 得出。

\medskip\noindent
上述构造的一个特殊例子是取 $f$ 为到单点拓扑斯 $\Sh(pt)$ 的态射。
此时 $\mu$ 正是杯积映射
$$
R\Gamma(\mathcal{C}, \mathcal{A})
\otimes_{\Gamma(\mathcal{C}, \mathcal{O})}^\mathbf{L}
R\Gamma(\mathcal{C}, \mathcal{A})
\longrightarrow
R\Gamma(\mathcal{C}, \mathcal{A}),
\quad \eta \otimes \theta \mapsto \eta \cup \theta
$$
，而类似地，$\mu_\mathcal{M}$ 是杯积映射
$$
R\Gamma(\mathcal{C}, \mathcal{M})
\otimes_{\Gamma(\mathcal{C}, \mathcal{O})}^\mathbf{L}
R\Gamma(\mathcal{C}, \mathcal{A})
\longrightarrow
R\Gamma(\mathcal{C}, \mathcal{M}),
\quad \eta \otimes \theta \mapsto \eta \cup \theta
$$
一般地，借助《站点上的上同调》评注
\ref{sites-cohomology-remark-before-Leray} 中的同一
$$
R\Gamma(\mathcal{C}, \mathcal{A}) =
R\Gamma(\mathcal{C}', Rf_*\mathcal{A})
\quad\text{且}\quad
R\Gamma(\mathcal{C}, \mathcal{M}) =
R\Gamma(\mathcal{C}', Rf_*\mathcal{M})
$$
，映射 $\mu_\mathcal{M}$ 在上同调上诱导杯积。为此，应用
《站点上的上同调》引理
\ref{sites-cohomology-lemma-compose-cup-product}，其中拓扑斯的第二个态射
取为上述从 $\Sh(\mathcal{C}')$ 到单点拓扑斯的态射。

\medskip\noindent
若 $\mathcal{M}_1\to\mathcal{M}_2$ 是右微分分次
$\mathcal{A}$-模的同态，则图
$$
\xymatrix{
Rf_*\mathcal{M}_1 \otimes_{\mathcal{O}'}^\mathbf{L} Rf_*\mathcal{A}
\ar[rr]_-{\mu_{\mathcal{M}_1}} \ar[d] & &
Rf_*\mathcal{M}_1 \ar[d] \\
Rf_*\mathcal{M}_2 \otimes_{\mathcal{O}'}^\mathbf{L} Rf_*\mathcal{A}
\ar[rr]^-{\mu_{\mathcal{M}_2}} & &
Rf_*\mathcal{M}_2
}
$$
在 $D(\mathcal{O}')$ 中交换；这源于相对杯积的函子性。设有右微分分次
$\mathcal{A}$-模的短正合列
$$
0 \to \mathcal{M}_1 \xrightarrow{a} \mathcal{M}_2 \to \mathcal{M}_3 \to 0
$$
。我们断言图
$$
\xymatrix{
Rf_*\mathcal{M}_3 \otimes_{\mathcal{O}'}^\mathbf{L} Rf_*\mathcal{A}
\ar[rr]_-{\mu_{\mathcal{M}_3}} \ar[d]_{Rf_*\delta \otimes \text{id}} & & 
Rf_*\mathcal{M}_3 \ar[d]^{Rf_*\delta} \\
Rf_*\mathcal{M}_1[1] \otimes_{\mathcal{O}'}^\mathbf{L} Rf_*\mathcal{A}
\ar[rr]^-{\mu_{\mathcal{M}_1[1]}} & &
Rf_*\mathcal{M}_1[1]
}
$$
在 $D(\mathcal{O}')$ 中交换，其中
$\delta:\mathcal{M}_3\to\mathcal{M}_1[1]$ 是给定短正合列所给出的
$D(\mathcal{O})$ 中的态射（见《导出范畴》第
\ref{derived-section-canonical-delta-functor} 节）。若我们的序列作为分次右
$\mathcal{A}$-模列分裂，则这一点很清楚，因为此时 $\delta$ 可由右
$\mathcal{A}$-模映射表示，从而可应用上述讨论。一般情形下，使用
$a$ 的锥和图
$$
\xymatrix{
\mathcal{M}_1 \ar[r]_a \ar[d] &
\mathcal{M}_2 \ar[r]_i \ar[d] &
C(a) \ar[r]_{-p} \ar[d]^q &
\mathcal{M}_1[1] \ar[d] \\
\mathcal{M}_1 \ar[r] &
\mathcal{M}_2 \ar[r] &
\mathcal{M}_3 \ar[r]^\delta &
\mathcal{M}_1[1]
}
$$
。由《导出范畴》引理
\ref{derived-lemma-derived-canonical-delta-functor} 中 $\delta$ 的定义，
右边的方块在 $D(\mathcal{O})$ 中交换。现在，锥 $C(a)$ 具有右微分分次
$\mathcal{A}$-模结构，使得 $i$、$p$、$q$ 是右微分分次
$\mathcal{A}$-模的同态；见定义 \ref{sdga-definition-cone}。因此由上可知，
对应图对态射 $q$ 与 $-p$ 交换。由于 $q$ 在 $D(\mathcal{O})$ 中是同构，
遂如所需得出对 $\delta$ 也有同样结论。

\medskip\noindent
在上述情形中，给定右微分分次 $\mathcal{A}$-模 $\mathcal{M}$，取
$$
\xi \in H^n(\mathcal{C}, \mathcal{M})
$$
。换言之，把 $\mathcal{M}$ 视为 $\mathcal{O}$-模复形时，$\xi$ 是其
次数 $n$ 上同调类。由引理 \ref{sdga-lemma-cohomology-ext}，可以构造右微分分次
$\mathcal{A}$-模的映射
$$
x : \mathcal{A} \rightarrow \mathcal{M}'[n]
\quad\text{且}\quad
s : \mathcal{M} \to \mathcal{M}'
$$
，其中 $s$ 是拟同构，并且 $\xi$ 是
$1\in H^0(\mathcal{C},\mathcal{A})$ 在导出范畴
$D(\mathcal{A},\text{d})$ 中经态射 $s[n]^{-1}\circ x$ 的像，
因而在导出范畴 $D(\mathcal{O})$ 中亦然。因此，$D(\mathcal{O})$
中的相应映射
$$
\xi' = (s[n])^{-1} \circ x : \mathcal{A} \longrightarrow \mathcal{M}[n]
$$
由以下两个性质唯一刻画：
\begin{enumerate}
\item $\xi'$ 可提升为 $D(\mathcal{A},\text{d})$ 中的态射；
\item 在
$H^0(\mathcal{C}, \mathcal{M}[n]) = H^n(\mathcal{C}, \mathcal{M})$
中有 $\xi=\xi'(1)$。
\end{enumerate}
由上面讨论的 $x$ 与 $s$ 同相对杯积的相容性可知，对环化拓扑斯的每个
态射\footnote{例如恒等态射。}
$(f, f^\sharp) : (\Sh(\mathcal{C}), \mathcal{O})
\to (\Sh(\mathcal{C}'), \mathcal{O}')$，$\xi'$ 的导出直像
$$
Rf_*\xi' : Rf_*\mathcal{A}  \longrightarrow Rf_*\mathcal{M}[n]
$$
与上述构造的映射 $\mu$ 和 $\mu_{\mathcal{M}[n]}$ 相容，即图
$$
\xymatrix{
Rf_*\mathcal{A} \otimes_{\mathcal{O}'}^\mathbf{L} Rf_*\mathcal{A}
\ar[rr]_-\mu \ar[d]_{Rf_*\xi' \otimes \text{id}}  & &
Rf_*\mathcal{A} \ar[d]^{Rf_*\xi'} \\
Rf_*\mathcal{M}[n] \otimes_{\mathcal{O}'}^\mathbf{L} Rf_*\mathcal{A}
\ar[rr]^-{\mu_{\mathcal{M}[n]}} & &
Rf_*\mathcal{M}[n]
}
$$
在 $D(\mathcal{O}')$ 中交换。把此相容性用于到单点拓扑斯的映射，
特别得到图
$$
\xymatrix{
R\Gamma(\mathcal{C}, \mathcal{A})
\otimes_{\Gamma(\mathcal{C}, \mathcal{O})}^\mathbf{L}
R\Gamma(\mathcal{C}, \mathcal{A})
\ar[d]_{\xi' \otimes \text{id}} \ar[r] &
R\Gamma(\mathcal{C}, \mathcal{A}) \ar[d]^{\xi'} \\
R\Gamma(\mathcal{C}, \mathcal{M}[n])
\otimes_{\Gamma(\mathcal{C}, \mathcal{O})}^\mathbf{L}
R\Gamma(\mathcal{C}, \mathcal{A})
\ar[r] &
R\Gamma(\mathcal{C}, \mathcal{M}[n])
}
$$
交换。结合 $\xi'(1)=\xi$，这说明所诱导的上同调映射
$$
\xi' : R\Gamma(\mathcal{C}, \mathcal{A}) \to
R\Gamma(\mathcal{C}, \mathcal{M}[n]), \quad \eta \mapsto \xi \cup \eta
$$
如所示由以 $\xi$ 作左杯乘给出。














\section{范畴上的微分分次模}
\label{sdga-section-modules-cohomology}

\noindent
本节续接《站点上的上同调》第
\ref{sites-cohomology-section-modules-cohomology} 节。

\medskip\noindent
设 $\mathcal{C}$ 是范畴。把 $\mathcal{C}$ 视为带混沌拓扑的位点。设
$\mathcal{O}$ 是 $\mathcal{C}$ 上的环层，
$(\mathcal{A},\text{d})$ 是微分分次 $\mathcal{O}$-代数层。换言之，
$\mathcal{O}$ 是范畴 $\mathcal{C}$ 上的环预层，而
$(\mathcal{A},\text{d})$ 是 $\mathcal{C}$ 上的微分分次
$\mathcal{O}$-代数预层；见《范畴》定义
\ref{categories-definition-presheaf}。

\begin{definition}
\label{sdga-definition-cartesian}
在上述情形中，以
{\it $\mathit{QC}(\mathcal{A}, \text{d})$}
表示 $D(\mathcal{A},\text{d})$ 的全子范畴，其对象为满足如下条件的
$M$：对 $\mathcal{C}$ 中所有 $U\to V$，典范映射
$$
R\Gamma(V, M) \otimes_{\mathcal{A}(V)}^\mathbf{L} \mathcal{A}(U)
\longrightarrow
R\Gamma(U, M)
$$
在 $D(\mathcal{A}(U),\text{d})$ 中是同构。
\end{definition}

\begin{lemma}
\label{sdga-lemma-cartesian-good-sub}
在上述情形中，子范畴 $\mathit{QC}(\mathcal{A},\text{d})$ 是
$D(\mathcal{A},\text{d})$ 的严格满的饱和三角子范畴，并在任意直和下保持。
\end{lemma}

\begin{proof}
设 $U$ 是 $\mathcal{C}$ 的对象。由于 $\mathcal{C}$ 上的拓扑是混沌的，
函子 $\mathcal{F}\mapsto\mathcal{F}(U)$ 正合且与直和交换。因此，
正合函子 $M\mapsto R\Gamma(U,M)$ 可如下计算：以任意微分分次
$\mathcal{A}$-模 $\mathcal{M}$ 表示 $K$，再取 $\mathcal{M}(U)$。
故 $R\Gamma(U,-)$ 与直和交换；见引理
\ref{sdga-lemma-derived-products}。类似地，给定 $\mathcal{C}$ 的态射
$U\to V$，导出张量积函子
$- \otimes_{\mathcal{O}(A)}^\mathbf{L} \mathcal{A}(U) :
D(\mathcal{A}(V)) \to D(\mathcal{A}(U))$
正合且与直和交换。本引理由这些观察直接得出；略去细节。
\end{proof}

\begin{remark}
\label{sdga-remark-cartesian-brown}
仍按上述记号，设 $\mathcal{C}$ 是视为带混沌拓扑位点的范畴，
$\mathcal{O}$ 是 $\mathcal{C}$ 上的环层，而
$(\mathcal{A},\text{d})$ 是微分分次 $\mathcal{O}$-代数层。则
《站点上的上同调》命题
\ref{sites-cohomology-proposition-cartesian-brown} 的类似结论对
$\mathit{QC}(\mathcal{A},\text{d})$ 成立，证明几乎完全相同：
\begin{enumerate}
\item 任何把直和变为积的反变上同调函子
$H : \mathit{QC}(\mathcal{A}, \text{d}) \to \textit{Ab}$
都是可表的；
\item 三角范畴之间任何把直和变为直和的正合函子
$F : \mathit{QC}(\mathcal{A}, \text{d}) \to \mathcal{D}$
都有正合右伴随；
\item 含入函子
$\mathit{QC}(\mathcal{A}, \text{d}) \to D(\mathcal{A}, \text{d})$
有正合右伴随。
\end{enumerate}
若今后需要此结论，我们将在这里精确表述并证明它。
\end{remark}

\noindent
设 $u:\mathcal{C}'\to\mathcal{C}$ 是范畴之间的函子。若把
$\mathcal{C}$ 与 $\mathcal{C}'$ 视为带混沌拓扑的位点，则 $u$ 是连续且
余连续的函子。因此得到拓扑斯的态射
$g:\Sh(\mathcal{C}')\to\Sh(\mathcal{C})$；见《位点》引理
\ref{sites-lemma-cocontinuous-morphism-topoi}。再设给定
$\mathcal{C}$ 上的环层 $\mathcal{O}$、$\mathcal{C}'$ 上的环层
$\mathcal{O}'$，以及映射
$g^\sharp : g^{-1}\mathcal{O} \to \mathcal{O}'$.
把相应的环化拓扑斯态射简记为
$g : (\Sh(\mathcal{C}'), \mathcal{O}') \to (\Sh(\mathcal{C}), \mathcal{O})$,
；见《位点上的模》第 \ref{sites-modules-section-ringed-topoi} 节。
最后，设 $(\mathcal{A},\text{d})$ 是微分分次 $\mathcal{O}$-代数层，
$(\mathcal{A}',\text{d})$ 是微分分次 $\mathcal{O}'$-代数层，并且给定
微分分次 $\mathcal{O}'$-代数的映射
$\varphi:g^*\mathcal{A}\to\mathcal{A}'$（见第
\ref{sdga-section-functoriality-dg} 节）。

\begin{lemma}
\label{sdga-lemma-cartesian-functorial}
设 $g:(\Sh(\mathcal{C}'),\mathcal{O}')
\to(\Sh(\mathcal{C}),\mathcal{O})$ 与
$\varphi:g^*\mathcal{A}\to\mathcal{A}'$ 如上。则函子
$Lg^* : D(\mathcal{A}, \text{d}) \to D(\mathcal{A}', \text{d})$
把 $\mathit{QC}(\mathcal{A},\text{d})$ 映入
$\mathit{QC}(\mathcal{A}',\text{d})$。
\end{lemma}

\begin{proof}
取 $U'\in\Ob(\mathcal{C}')$，其在 $\mathcal{C}$ 中的像记为
$U=u(U')$。以 $pt$ 表示仅有一个对象和一个态射的范畴。把所示环化拓扑斯
$(\Sh(pt),\mathcal{O}'(U'))$ 与 $(\Sh(pt),\mathcal{O}(U))$
分别赋予微分分次代数 $\mathcal{A}'(U)$ 与 $\mathcal{A}(U)$。
当然，把其上的微分分次模导出范畴分别等同于
$D(\mathcal{A}'(U'),\text{d})$ 与
$D(\mathcal{A}(U),\text{d})$。于是有环化拓扑斯的交换图
$$
\xymatrix{
(\Sh(pt), \mathcal{O}'(U')) \ar[rr]_{U'} \ar[d] & &
(\Sh(\mathcal{C}'), \mathcal{O}') \ar[d]^g \\
(\Sh(pt), \mathcal{O}(U)) \ar[rr]^U & &
(\Sh(\mathcal{C}), \mathcal{O})
}
$$
，各顶点均赋有相应的微分分次代数。沿下方水平态射拉回，把
$D(\mathcal{A},\text{d})$ 中的 $M$ 映为视作
$D(\mathcal{A}(U),\text{d})$ 中对象的 $R\Gamma(U,K)$。
沿左边竖直箭头拉回，把 $M$ 映为
$M \otimes_{\mathcal{A}(U)}^\mathbf{L} \mathcal{A}'(U')$.
沿此图任一方向复合都给出相同结果（引理
\ref{sdga-lemma-compose-pullback}），故得到
$$
R\Gamma(U', Lg^*K) =
R\Gamma(U, K) \otimes_{\mathcal{A}(U)}^\mathbf{L} \mathcal{A}'(U')
$$
最后，设 $f':U'\to V'$ 是 $\mathcal{C}'$ 中的态射，并把它在
$\mathcal{C}$ 中的像记为
$f=u(f'):U=u(U')\to V=u(V')$。若
$K\in\mathit{QC}(\mathcal{A},\text{d})$，则
\begin{align*}
R\Gamma(V', Lg^*K) \otimes_{\mathcal{A}'(V')}^\mathbf{L} \mathcal{A}'(U')
& =
R\Gamma(V, K) \otimes_{\mathcal{A}(V)}^\mathbf{L} \mathcal{A}'(V')
\otimes_{\mathcal{A}'(V')}^\mathbf{L} \mathcal{A}'(U') \\
& =
R\Gamma(V, K) \otimes_{\mathcal{A}(V)}^\mathbf{L} \mathcal{A}'(U') \\
& =
R\Gamma(V, K) \otimes_{\mathcal{A}(V)}^\mathbf{L} \mathcal{A}(U)
\otimes_{\mathcal{A}(U)}^\mathbf{L} \mathcal{A}'(U') \\
& =
R\Gamma(U, K) \otimes_{\mathcal{A}(U)}^\mathbf{L} \mathcal{A}'(U') \\
& =
R\Gamma(U', Lg^*K)
\end{align*}
，正如所需。这里对 $U'$ 与 $V'$ 都使用了上述观察。
\end{proof}






\section{范畴上的微分分次模（续）}
\label{sdga-section-modules-cohomology-bis}

\noindent
我们进一步给出第 \ref{sdga-section-modules-cohomology} 节所引入的拟凝聚模概念的
若干结果。

\begin{lemma}
\label{sdga-lemma-cartesian-cofinal-subcategory}
设 $\mathcal{C}$、$\mathcal{O}$、$\mathcal{A}$ 如第
\ref{sdga-section-modules-cohomology} 节。设
$\mathcal{C}'\subset\mathcal{C}$ 是具有如下性质的全子范畴：对每个
$U\in\Ob(\mathcal{C})$，由箭头 $U\to U'$ 构成的范畴
$U/\mathcal{C}'$ 是余滤过的。以 $\mathcal{O}'$、$\mathcal{A}'$
表示 $\mathcal{O}$、$\mathcal{A}$ 到 $\mathcal{C}'$ 的限制。则限制诱导等价
$\mathit{QC}(\mathcal{A}, \text{d}) \to \mathit{QC}(\mathcal{A}', \text{d})$.
\end{lemma}

\begin{proof}
我们构造该函子的拟逆。设
$M'$ 是 $\mathit{QC}(\mathcal{A}',\text{d})$ 的对象。可以用良好微分分次模
$\mathcal{M}'$ 表示 $M'$；见引理 \ref{sdga-lemma-resolve}。于是对每个
$U'\in\Ob(\mathcal{C}')$，微分分次
$\mathcal{A}'(U')$-模 $\mathcal{M}'(U)$ 是 K-平坦且分次平坦的；
对 $\mathcal{C}'$ 的每个态射 $U'_1\to U'_2$，映射
$$
\mathcal{M}'(U'_2) \otimes_{\mathcal{A}'(U'_2)} \mathcal{A}'(U'_1)
\longrightarrow
\mathcal{M}'(U'_1)
$$
是拟同构（因为源表示导出张量积）。考虑按如下规则定义的微分分次
$\mathcal{A}$-模 $\mathcal{M}$：
$$
\mathcal{M}(U) =
\colim_{U \to U' \in U/\mathcal{C}'}
\mathcal{M}'(U') \otimes_{\mathcal{A}'(U')} \mathcal{A}(U)
$$
。由本引理中的假设，这是复形的滤过余极限。由于
$M'\in\mathit{QC}(\mathcal{A}',\text{d})$，该系统中的所有转移映射都是
拟同构。又因滤过余极限正合，可知对任意态射 $U\to U'$，其中
$U'\in\Ob(\mathcal{C}')$，$D(\mathcal{A}(U),\text{d})$ 中的
$\mathcal{M}(U)$ 同构于
$\mathcal{M}'(U')\otimes_{\mathcal{A}'(U')}\mathcal{A}(U)$。

\medskip\noindent
我们断言 $\mathcal{M}\in\mathit{QC}(\mathcal{A},\text{d})$：事实上，
给定 $\mathcal{C}$ 中的 $U\to V$，选择映射 $V\to V'$，其中
$V'\in\Ob(\mathcal{C}')$。由上可知，映射
$\mathcal{M}(V)\to\mathcal{M}(U)$ 等同于映射
$$
\mathcal{M}'(V') \otimes_{\mathcal{A}'(V')} \mathcal{A}(V)
\longrightarrow
\mathcal{M}'(V') \otimes_{\mathcal{A}'(V')} \mathcal{A}(U)
$$
。由于 $\mathcal{M'}(V')$ 作为微分分次
$\mathcal{A}'(V')$-模是 K-平坦的，断言得证。

\medskip\noindent
自然映射 $\mathcal{M}|_{\mathcal{C}'}\to\mathcal{M}'$ 在
$D(\mathcal{A}',d)$ 中是同构，这立即由上得出。

\medskip\noindent
反之，若有 $\mathit{QC}(\mathcal{A},\text{d})$ 的对象 $E$，则以良好
微分分次模 $\mathcal{E}$ 表示它。置
$\mathcal{M}'=\mathcal{E}|_{\mathcal{C}'}$（这仍是良好微分分次模），
可知存在映射
$$
\mathcal{E} \to \mathcal{M}
$$
，它在 $\mathcal{C}$ 的 $U$ 上由映射
$$
\mathcal{E}(U)
\longrightarrow
\colim_{U \to U' \in U/\mathcal{C}'}
\mathcal{E}(U') \otimes_{\mathcal{A}'(U')} \mathcal{A}(U)
$$
给出；同理这是拟同构。因此限制与上述构造是互为拟逆的函子，正如所需。
\end{proof}

\begin{lemma}
\label{sdga-lemma-cartesian-qis-equivalence}
设 $\mathcal{C}$、$\mathcal{O}$ 如第
\ref{sdga-section-modules-cohomology} 节。设
$\varphi:\mathcal{A}\to\mathcal{B}$ 是微分分次
$\mathcal{O}$-代数的同态，并在上同调层上诱导同构，则引理
\ref{sdga-lemma-qis-equivalence} 的等价
$D(\mathcal{A}, \text{d}) \to D(\mathcal{B}, \text{d})$
诱导等价
$\mathit{QC}(\mathcal{A}, \text{d}) \to
\mathit{QC}(\mathcal{B}, \text{d})$.
\end{lemma}

\begin{proof}
只需证明如下事实：给定 $\mathcal{C}$ 的态射 $U\to V$ 与
$D(\mathcal{A},\text{d})$ 中的 $M$，下列条件等价：
\begin{enumerate}
\item $R\Gamma(V, M) \otimes_{\mathcal{A}(V)}^\mathbf{L} \mathcal{A}(U)
\to \Gamma(U, M)$ 在 $D(\mathcal{A}(U),\text{d})$ 中是同构；
\item $R\Gamma(V, M \otimes_\mathcal{A}^\mathbf{L} \mathcal{B})
\otimes_{\mathcal{B}(V)}^\mathbf{L} \mathcal{B}(U)
\to \Gamma(U, M \otimes_\mathcal{A}^\mathbf{L} \mathcal{B})$
在 $D(\mathcal{B}(U),\text{d})$ 中是同构。
\end{enumerate}
由于 $\mathcal{C}$ 上的拓扑是混沌的，这只归结为
$\mathcal{A}(U)\to\mathcal{B}(U)$ 与
$\mathcal{A}(V)\to\mathcal{B}(V)$ 是拟同构这一事实。略去细节。
\end{proof}






\section{微分分次代数的逆系统}
\label{sdga-section-qc-inverse-system}

\noindent
本节考虑第 \ref{sdga-section-modules-cohomology} 节所讨论情形的如下特殊情况：
\begin{enumerate}
\item $\mathcal{C}$ 是范畴 $\mathbf{N}$，当且仅当 $i\leq j$ 时，
存在唯一态射 $i\to j$；
\item $\mathcal{O}$ 是取值为给定环 $R$ 的常值环（预）层。
\end{enumerate}
在此设定下，微分分次 $\mathcal{O}$-代数层 $\mathcal{A}$ 等同于
微分分次 $R$-代数的逆系统 $(A_n)$。微分分次
$\mathcal{A}$-模层 $\mathcal{M}$ 等同于逆系统 $(M_n)$，其中
$M_n$ 是微分分次 $A_n$-模，而转移映射
$M_{n+1}\to M_n$ 是 $A_{n+1}$-模映射。本节在记号中省略微分，
并以 $\mathit{QC}(\mathcal{A})$ 表示定义
\ref{sdga-definition-cartesian} 中记为
$\mathit{QC}(\mathcal{A},\text{d})$ 的范畴。

\medskip\noindent
设 $\mathcal{B}=(B_n)$ 是微分分次 $R$-代数的另一个逆系统。给定
pro-对象的态射 $\varphi:(A_n)\to(B_n)$，我们将构造从
$\mathit{QC}(\mathcal{A})$ 到 $\mathit{QC}(\mathcal{B})$ 的正合函子。
具体地，按照《范畴》例
\ref{categories-example-pro-morphism-inverse-systems}，态射 $\varphi$
由整数序列 $\ldots\geq m(3)\geq m(2)\geq m(1)$ 及交换图
$$
\xymatrix{
\ldots \ar[r] &
A_{m(3)} \ar[d]^{\varphi_3} \ar[r] &
A_{m(2)} \ar[d]^{\varphi_2} \ar[r] &
A_{m(1)} \ar[d]^{\varphi_1} \\
\ldots \ar[r] &
B_3 \ar[r] &
B_2 \ar[r] &
B_1
}
$$
给出，其中各项是微分分次 $R$-代数。于是，给定表示
$\mathit{QC}(\mathcal{A})$ 中对象的良好微分分次
$\mathcal{A}$-模层 $\mathcal{M}=(M_n)$，可以置
$$
N_n = M_{m(n)} \otimes_{A_{m(n)}} B_n
$$
。此逆系统确定 $\mathit{QC}(\mathcal{B})$ 的对象，因为
$A_{m(n)}$-模 $M_{m(n)}$ 是 K-平坦的；略去细节。所得函子不依赖其构造中
所作选择这一事实也留给读者证明。

\begin{lemma}
\label{sdga-lemma-pro-isomorphic}
在上述情形中，设 $\mathcal{A}=(A_n)$ 和 $\mathcal{B}=(B_n)$ 是微分分次
$R$-代数的逆系统。若 $\varphi:(A_n)\to(B_n)$ 是 pro-对象的同构，则上面构造的
函子 $\mathit{QC}(\mathcal{A})\to\mathit{QC}(\mathcal{B})$ 是等价。
\end{lemma}

\begin{proof}
令 $\psi:(B_n)\to(A_n)$ 为与 $\varphi$ 互逆的 pro-对象态射。根据
《范畴》，例 \ref{categories-example-pro-morphism-inverse-systems} 中的讨论，
可设 $\varphi$ 由如上的一族映射给出，而 $\psi$ 由
$n(1)<n(2)<\ldots$ 以及如下交换图给出：
$$
\xymatrix{
\ldots \ar[r] &
B_{n(3)} \ar[d]^{\psi_3} \ar[r] &
B_{n(2)} \ar[d]^{\psi_2} \ar[r] &
B_{n(1)} \ar[d]^{\psi_1} \\
\ldots \ar[r] &
A_3 \ar[r] &
A_2 \ar[r] &
A_1
}
$$
其中各项是微分分次 $R$-代数。由于 $\varphi\circ\psi=\text{id}$，在可能增大
函数 $n(\cdot)$ 和 $m(\cdot)$ 的取值后，可以假设
$B_{n(m(i))}\to A_{m(i)}\to B_i$ 是转移映射。因此，函子的复合
$$
\mathit{QC}(\mathcal{B}) \to
\mathit{QC}(\mathcal{A}) \to
\mathit{QC}(\mathcal{B})
$$
把良好微分分次 $\mathcal{B}$-模层 $\mathcal{N}=(N_n)$ 送到取值为
$$
N'_i = N_{n(m(i))} \otimes_{B_{n(m(i))}} B_i
$$
的逆系统 $\mathcal{N}'=(N'_i)$；它与 $\mathcal{N}$ 典范拟同构，恰是因为
$\mathcal{N}$ 是 $\mathit{QC}(\mathcal{B})$ 的对象，且对所有 $j$，$N_j$
都是 K-平坦微分分次模。反向复合也有同样结论，故证毕。
\end{proof}

\noindent
令 $\mathcal{C}=\mathbf{N}$，令 $\mathcal{O}$ 为取值于环 $R$ 的常值层，并令
$\mathcal{A}$ 由微分分次 $R$-代数的逆系统 $(A_n)$ 给出。设有两个左微分分次
$\mathcal{A}$-模 $\mathcal{N}$ 和 $\mathcal{N}'$，它们分别由逆系统
$(N_n)$ 和 $(N'_n)$ 给出。因此每个 $N_n$ 和 $N'_n$ 都是左微分分次
$A_n$-模。暂称 $(N_n)$ 和 $(N'_n)$ {\it 在导出范畴中 pro-同构}，如果存在
整数序列
$$
1 = n_0 < n_1 < n_2 < n_3 < \ldots
$$
以及映射
$$
N_{n_{2i}} \to N'_{n_{2i - 1}}
\quad\text{于}\quad
D(A_{n_{2i}}^{opp}, \text{d})
$$
和
$$
N'_{n_{2i + 1}} \to N'_{n_{2i}}
\quad\text{于}\quad
D(A_{n_{2i + 1}}^{opp}, \text{d})
$$
使得复合 $N_{n_{2i}}\to N_{n_{2i-2}}$ 和
$N'_{n_{2i+1}}\to N'_{2i-1}$ 分别由相应系统的转移映射给出。

\begin{lemma}
\label{sdga-lemma-pro-isomorphic-systems-tensor}
若 $(N_n)$ 和 $(N'_n)$ 按上述定义在导出范畴中 pro-同构，则对
$D(\mathbf{N},\mathcal{A})$ 的每个对象 $(M_n)$，有
$$
R\lim (M_n \otimes_{A_n}^\mathbf{L} N_n) =
R\lim (M_n \otimes_{A_n}^\mathbf{L} N'_n)
$$
于 $D(R)$ 中。
\end{lemma}

\begin{proof}
该假设蕴涵 $D(R)$ 中的逆系统
$(M_n\otimes_{A_n}^\mathbf{L}N_n)$ 与逆系统
$(M_n\otimes_{A_n}^\mathbf{L}N'_n)$ pro-同构（通常意义下）。因此结论由
《代数详论》，引理 \ref{more-algebra-lemma-pro-isomorphism-Rlim} 得出。
\end{proof}

\begin{lemma}
\label{sdga-lemma-different-tensors}
\begin{reference}
这是 \cite[引理 3.5.4]{BS} 的一个变体。
\end{reference}
设 $R$ 为环，$f_1,\ldots,f_r\in R$。令 $K_n$ 为
$f_1^n,\ldots,f_r^n$ 上的 Koszul 复形，并将其视为微分分次 $R$-代数。
设 $(M_n)$ 是 $D(\mathbf{N},(K_n))$ 的对象。则对任意 $t\geq1$，有
$$
R\lim (M_n \otimes_R^\mathbf{L} K_t) =
R\lim (M_n \otimes_{K_n}^\mathbf{L} K_t)
$$
于 $D(R)$ 中。
\end{lemma}

\begin{proof}
固定 $t\geq1$。对 $n\geq t$，用 ${}_nK_t$ 表示微分分次 $R$-代数 $K_t$，
将其视为左微分分次 $K_n$-模。注意
$$
M_n \otimes_R^\mathbf{L} K_t =
M_n \otimes_{K_n}^\mathbf{L} (K_n \otimes_R^\mathbf{L} K_t) =
M_n \otimes_{K_n}^\mathbf{L} (K_n \otimes_R K_t)
$$
。因此由引理 \ref{sdga-lemma-pro-isomorphic-systems-tensor}，只需证明
$({}_nK_t)$ 和 $(K_n\otimes_RK_t)$ 在导出范畴中 pro-同构。乘法映射
$$
K_n \otimes_R K_t \longrightarrow {}_nK_t
$$
是左微分分次 $K_n$-模的映射。因此，为完成证明，只需证明对所有
$n\geq1$，存在 $N>n$ 以及映射
$$
{}_NK_t \longrightarrow {}_NK_n \otimes_R K_t
$$
于 $D(K_N^{opp},\text{d})$ 中，使它与乘法映射的复合（沿任一方向）是转移
映射。这已在《除幂代数》，引理 \ref{dpa-lemma-construct-some-maps} 中通过
一个显式构造完成。
\end{proof}

\begin{proposition}
\label{sdga-proposition-derived-complete}
设 $R$ 为诺特环，$I\subset R$ 为理想。以下三个范畴典范等价：
\begin{enumerate}
\item 令 $\mathcal{A}$ 为 $\mathbf{N}$ 上与 $R$-代数逆系统
$A_n=R/I^n$ 对应的 $R$-代数层。范畴 $\mathit{QC}(\mathcal{A})$。
\item 选取 $I$ 的生成元 $f_1,\ldots,f_r$。令 $\mathcal{B}$ 为
$\mathbf{N}$ 上与 $f_1^n,\ldots,f_r^n$ 上的 Koszul 代数逆系统对应的
微分分次 $R$-代数层。范畴 $\mathit{QC}(\mathcal{B})$。
\item 导出完备对象所成的全子范畴
$D_{comp}(R,I)\subset D(R)$；见《代数详论》，
定义 \ref{more-algebra-definition-derived-complete} 及其后的文字。
\end{enumerate}
\end{proposition}

\begin{proof}
考虑环化拓扑斯的显然态射
$f : (\Sh(\mathbf{N}), \mathcal{A}) \to (\Sh(pt), R)$
，并考虑伴随函子 $Lf^*$ 和 $Rf_*$。前者限制为函子
$$
F : D_{comp}(R, I) \longrightarrow \mathit{QC}(\mathcal{A})
$$
，它把由 K-平坦复形 $K^\bullet$ 表示的 $D_{comp}(R,I)$ 中对象 $K$ 送到
$\mathit{QC}(\mathcal{A})$ 的对象 $(K^\bullet\otimes_RR/I^n)$。后者限制为函子
$$
G : \mathit{QC}(\mathcal{A}) \longrightarrow D_{comp}(R, I)
$$
，它把 $\mathit{QC}(\mathcal{A})$ 的对象 $(M_n^\bullet)$ 送到
$R\lim M_n^\bullet$。例如，由《代数详论》，
引理 \ref{more-algebra-lemma-naive-derived-completion}，所得对象是导出完备的。
此外，由《代数详论》，
命题 \ref{more-algebra-proposition-noetherian-naive-completion-is-completion}，
有 $G\circ F=\text{id}$。因此，为看出 $F$ 和 $G$ 是互为拟逆的等价，只需看出
$G$ 的核为零（见《导出范畴》，
引理 \ref{derived-lemma-fully-faithful-adjoint-kernel-zero}）。
然而，直接证明这一点似乎并不容易！

\medskip\noindent
本段将证明 $\mathit{QC}(\mathcal{A})$ 和 $\mathit{QC}(\mathcal{B})$ 等价。
写 $\mathcal{B}=(B_n)$，其中 $B_n$ 是 Koszul 复形，并将它视为集中在次数
$-r,-r+1,\ldots,0$ 的上链复形。由《除幂代数》，
注 \ref{dpa-remark-pro-system-koszul}（但把链复形改成上链复形），可以找到
$1<n_1<n_2<\ldots$ 以及微分分次 $R$-代数的映射
$B_{n_i} \to E_i \to R/(f_1^{n_i}, \ldots, f_r^{n_i})$
和 $E_i\to B_{n_{i-1}}$，使得
$$
\xymatrix{
B_{n_1} \ar[d] &
B_{n_2} \ar[d] \ar[l] &
B_{n_3} \ar[d] \ar[l] &
\ldots \ar[l] \\
E_1 \ar[d] &
E_2 \ar[l] \ar[d] &
E_3 \ar[l] \ar[d] &
\ldots \ar[l] \\
B_1 &
B_{n_1} \ar[l] &
B_{n_2} \ar[l] &
\ldots \ar[l]
}
$$
是微分分次 $R$-代数的交换图，且
$E_i\to R/(f_1^{n_i},\ldots,f_r^{n_i})$ 是拟同构。由此得到：
\begin{enumerate}
\item $\mathit{QC}(\mathcal{B})$ 与 $\mathit{QC}((E_i))$ 之间有一个等价；
\item $\mathit{QC}((E_i))$ 与
$\mathit{QC}((R/(f_1^{n_i},\ldots,f_r^{n_i})))$ 之间有一个等价；
\item $\mathit{QC}((R/(f_1^{n_i},\ldots,f_r^{n_i})))$ 与
$\mathit{QC}(\mathcal{A})$ 之间有一个等价。
\end{enumerate}
事实上，对 (1)，可把引理 \ref{sdga-lemma-pro-isomorphic} 应用于上图；该图表明
$(E_i)$ 与 $(B_n)$ pro-同构。对 (2)，可把引理
\ref{sdga-lemma-cartesian-qis-equivalence} 应用于拟同构的逆系统
$E_i\to R/(f_1^{n_i},\ldots,f_r^{n_i})$。对 (3)，可应用引理
\ref{sdga-lemma-pro-isomorphic} 以及逆系统 $(R/I^n)$ 与
$(R/(f_1^{n_i},\ldots,f_r^{n_i})$ pro-同构这一初等事实。

\medskip\noindent
与证明的第一段完全一样，可以定义伴随函子\footnote{可以证明，通过上述等价，
这些函子与先前定义的 $F$ 和 $G$ 相容。}
$$
F' : D_{comp}(R, I) \longrightarrow \mathit{QC}(\mathcal{B})
\quad\text{及}\quad
G' : \mathit{QC}(\mathcal{B}) \longrightarrow D_{comp}(R, I).
$$
前者把由 K-平坦复形 $K^\bullet$ 表示的 $D_{comp}(R,I)$ 中对象 $K$ 送到
$\mathit{QC}(\mathcal{B})$ 的对象 $(K^\bullet\otimes_RB_n)$。后者把
$\mathit{QC}(\mathcal{B})$ 的对象 $(M_n)$ 送到 $R\lim M_n$。与上面一样论证，
只需证明 $G'$ 的核为零。因此，令 $\mathcal{M}=(M_n)$ 为 $\mathcal{B}$ 上的
良好微分分次模层，它表示 $\mathit{QC}(\mathcal{B})$ 中属于 $G'$ 之核的一个
对象。则
$$
0 = R\lim M_n \Rightarrow
0 = (R\lim M_n) \otimes_R^\mathbf{L} B_t =
R\lim (M_n \otimes_R^\mathbf{L} B_t)
$$
。由引理 \ref{sdga-lemma-different-tensors}，有
$R\lim (M_n \otimes_R^\mathbf{L} B_t) =
R\lim (M_n \otimes_{B_n}^\mathbf{L} B_t)$.
由于 $(M_n)$ 是 $\mathit{QC}(\mathcal{B})$ 的对象，可见逆系统
$M_n\otimes_{B_n}^\mathbf{L}B_t$ 最终为常值 $M_t$。所以如所需，$M_t=0$。
\end{proof}

\begin{remark}
\label{sdga-remark-proregular}
设 $R$ 为环，$f_1,\ldots,f_r\in R$ 为生成理想 $I$ 的元素序列。令 $K_n$ 为
$f_1^n,\ldots,f_r^n$ 上的 Koszul 复形，并将它视为微分分次 $R$-代数。
若对所有 $n$，都存在 $m>n$，使得 $K_m\to K_n$ 在除次数 $0$ 以外的
上同调上诱导零映射，则称 $f_1,\ldots,f_r$ 为{\it 弱 pro-正则序列}。
若是如此，则即使 $R$ 不是诺特环，命题 \ref{sdga-proposition-derived-complete}
证明中的论证仍然成立。特别地可见，只要 $I$ 可由弱 pro-正则序列生成，
$\mathit{QC}(\{R/I^n\})$ 作为 $R$-线性三角范畴就与导出完备对象所成的范畴
$D_{comp}(R,I)$ 等价。如有需要，我们将在此精确陈述并证明之。
\end{remark}
